\documentclass[11pt,reqno]{amsart}
\usepackage[margin=1.05in]{geometry}
\usepackage{amsmath,amssymb,amsthm,mathtools,booktabs,longtable,array,enumitem}
\usepackage[hyphens]{url}
\usepackage[colorlinks=true,linkcolor=blue,citecolor=blue,urlcolor=blue]{hyperref}
\usepackage[expansion=false]{microtype}

\setlist{leftmargin=2em}
\emergencystretch=2em

\theoremstyle{plain}
\newtheorem{theorem}{Theorem}[section]
\newtheorem{proposition}[theorem]{Proposition}
\newtheorem{lemma}[theorem]{Lemma}
\newtheorem{corollary}[theorem]{Corollary}
\newtheorem{fact}[theorem]{Fact}
\theoremstyle{definition}
\newtheorem{definition}[theorem]{Definition}
\newtheorem{remark}[theorem]{Remark}
\newtheorem{assumption}[theorem]{Assumption}
\newtheorem{example}[theorem]{Example}

% Pointer to the script that produced a number
\newcommand{\cert}[2]{\par\smallskip{\raggedright\noindent\emph{Computation:} \texttt{#1}; output \texttt{#2}.\par}}
% Status label for statements that are computer-assisted or numerical
\newcommand{\status}[1]{\textup{\textsc{[#1]}}}

\newcommand{\R}{\mathbb R}
\newcommand{\hG}{h_\Gamma}
\newcommand{\Gh}{\Gamma_h}
\newcommand{\Oh}{\Omega_h}
\newcommand{\Th}{\mathcal T_h}
\newcommand{\FG}{\mathcal F_h^\Gamma}
\newcommand{\Vh}{V_h}
\newcommand{\VR}{V_h^R}
\newcommand{\VO}{V_h^0}
\newcommand{\Pih}{\Pi_h}
\newcommand{\Zh}{Z_h}
\newcommand{\Wh}{W_h}
\newcommand{\gI}{g_I}
\newcommand{\dn}{\partial_n}
\newcommand{\ut}{\tilde u}
\newcommand{\pt}{\tilde p}
\newcommand{\ft}{\tilde f}
\newcommand{\pbar}{\bar p}
\newcommand{\pbG}{\bar p_{\Gamma}}
\newcommand{\ip}[2]{\langle #1,\, #2\rangle}
\newcommand{\tnorm}[1]{|\!|\!|#1|\!|\!|}
\newcommand{\norm}[1]{\lVert #1\rVert}
\newcommand{\GS}{\mathrm{GS}}
\newcommand{\CNS}{\mathrm{CNS}^\ast}
\newcommand{\epsd}{\varepsilon_h^{(3)}}
\newcommand{\curl}{\operatorname{curl}}
\renewcommand{\div}{\operatorname{div}}
\newcommand{\Rc}{\mathcal R}
\newcommand{\Gc}{\mathcal G}
\newcommand{\M}{\mathcal M}
\newcommand{\hh}{\mathfrak h}
\newcommand{\gh}{\hat g}
\newcommand{\Aop}{\mathfrak A}
\newcommand{\Ut}{\tilde U}
\newcommand{\Pt}{\tilde P}
\newcommand{\Yh}{Y_h}
\newcommand{\Tr}{\mathcal T^{\rm tr}_h}
\newcommand{\RM}{\mathcal{RM}}
\DeclareMathOperator{\tr}{tr}
\DeclareMathOperator{\diam}{diam}
\DeclareMathOperator{\dist}{dist}
\DeclareMathOperator{\conv}{conv}

\begin{document}

%% TBD: authorship may become joint; author list, affiliations and CRediT statement to be settled before submission.
\title[Scott--Vogelius--Nitsche on inscribed polyhedra]{Scott--Vogelius--Nitsche on inscribed polyhedra:\\ convergence, sharpness and the penalty threshold in three dimensions}
\author{Jaideep Sai Padhi}
\address{Department of Computer Science, Purdue University, West Lafayette, IN 47907, USA}
\email{jpadhi@purdue.edu}
\date{\today}
\subjclass[2020]{Primary 65N30; Secondary 65N12, 65N15, 76D07}
\keywords{Stokes equations, Navier--Stokes equations, Scott--Vogelius elements, Nitsche's method, curved boundaries, inscribed polyhedra, Alfeld splits, divergence-free finite elements}

\begin{abstract}
Scott--Vogelius elements give exactly divergence-free velocities, but at a curved no-slip wall that is approximated by an inscribed polygon or polyhedron they are hampered by gradient constraints at the wall. Gjerde and Scott proposed to impose the wall condition weakly, by Nitsche's method, and Scott's zero-gradient prize asks whether this Scott--Vogelius--Nitsche method converges with the error $\hG^{3/2}+h^k$ expected from the scalar theory. We answer the question in three dimensions, for a smooth wall surface $\Gamma$ replaced by a triangulated surface $\Gh$ with vertices on $\Gamma$, straight-sided tetrahedra, continuous $P_k$ velocities and discontinuous $P_{k-1}$ pressures.

The method as printed omits the pressure part of the wall traction. When the wall pressure is not constant, the discrete fluid then leaks through the wall with normal velocity $(h/\mu)(p-\pbar)$, where $\mu/h$ is the Nitsche penalty; the error is exactly of order $(h/\mu)^{1/2}$ in energy and $h/\mu$ in $H^1$, and the $H^1$ constant is the Dirichlet energy of the Stokes flow driven by the leak, which we identify for every admissible penalty. Adding the traction with the wall mean of the discrete pressure removed, a device of Frachon, Nilsson and Zahedi, gives an exactly divergence-free method with error $O(\hG^{3/2}+h^k)$. On Alfeld refinements this holds for every $k\ge3$, and the term $\hG^{3/2}$ is sharp: the error is bounded below by $c\,\hG^{3/2}\norm{|\hh|\,\partial_Nu}_{L^2(\Gamma)}-C\hG^2$, where $\hh$ is the second fundamental form. The proofs use a lift theorem for inscribed surfaces, a uniform discrete inf-sup condition, and local witnesses verified by exact computations. We also show that, under local conditions at a smallest wall face, the penalty must grow with the ratio of bulk to wall mesh size, treat drag, lift and torque, and extend the results to steady Navier--Stokes flow under linearised-stability assumptions. A three-dimensional implementation confirms exact incompressibility and the leak law; the sphere tests we could run are still pre-asymptotic for the rate $\hG^{3/2}$.
\end{abstract}

\maketitle

\section{Introduction}\label{sec:intro}

\subsection{Exactly divergence-free elements}
Most mixed finite element methods for incompressible flow impose the constraint $\div u=0$ only weakly, against the discrete pressures. The velocity error then depends on the pressure, and this can be large when pressure gradients are large or the viscosity is small; the review of John, Linke, Merdon, Neilan and Rebholz \cite{JohnEtAl2017} explains why this matters and what can be done about it. One remedy is to choose the velocity and pressure spaces so that the divergence of every discrete velocity is itself a discrete pressure. Then the weak constraint forces $\div u_h=0$ pointwise, mass is conserved exactly, and on fitted meshes with strongly imposed no-slip data the velocity error does not see the pressure.

The oldest such pair is due to Scott and Vogelius \cite{ScottVogelius1985}: continuous piecewise-$P_k$ velocities with discontinuous piecewise-$P_{k-1}$ pressures. In two dimensions they proved $\div\Vh=\Pih$ for $k\ge4$ on meshes without singular vertices, and Guzm\'an and Scott \cite{GuzmanScott2019} later made the stability constant explicit in the local mesh geometry. Lower degrees become possible on split meshes; for example, quadratic velocities and linear pressures are stable on Clough--Tocher splits \cite{ArnoldQin1992}.

Three dimensions are much harder, and the theory is still incomplete. Zhang proved stability of the Scott--Vogelius pair on Alfeld (barycentric) refinements for $k\ge3$ \cite{Zhang2005}, and on Freudenthal meshes for $k\ge6$ \cite{Zhang2011}; Guzm\'an and Neilan \cite{GuzmanNeilan2018} gave a short proof on barycentric refinements in every dimension $d$, for $k\ge d$, built from a \emph{local} stability result on a single split simplex. Fu, Guzm\'an and Neilan \cite{FuGuzmanNeilan2020} constructed exact smooth piecewise-polynomial sequences on Alfeld splits, with commuting projections. On Worsey--Farin splits, Guzm\'an, Lischke and Neilan \cite{GuzmanLischkeNeilan2022} and Fabien, Guzm\'an, Neilan and Zytoon \cite{FabienGuzmanNeilanZytoon2022} obtained divergence-free pairs of low degree, whose pressure spaces carry extra conditions at the singular edges of the split. Neilan \cite{Neilan2015} constructs smooth discrete de Rham complexes in three dimensions, and Farrell, Mitchell and Scott \cite{FarrellMitchellScott2024} formulate conjectures on the Stokes complex on Freudenthal meshes; we know of no unconditional stability result for the Scott--Vogelius pair on general unsplit tetrahedral meshes. Neilan's review \cite{Neilan2020} surveys the field. In this paper, all unconditional results are on Alfeld refinements.

\subsection{Curved walls}
Exact incompressibility interacts badly with curved boundaries. Replacing a smooth domain by an inscribed polygon or polyhedron is the cheapest way to mesh it. For scalar second-order problems this costs order $h^{3/2}$ in $H^1$ \cite{BergerScottStrang1972,Scott1975}, and classical remedies either correct the boundary values \cite{BrambleDupontThomee1972,DupontGuzmanScott2025} or fit the boundary with curved elements. For Scott--Vogelius elements the cost is higher. In two dimensions, a smooth divergence-free field that vanishes on two edges meeting at a polygon vertex has zero gradient there, and discrete divergence-free fields inherit such \emph{zero-gradient} constraints at the vertices of an inscribed wall. If no-slip is imposed strongly on the inscribed polygon, the $H^1$ error can then lock at order $h^{1/2}$ \cite{ScottPrize,GjerdeScott2024}; in the companion paper \cite{PadhiLong2D} we trace the lock to wall vertices that lie in only two triangles.

Three routes avoid the problem. (i) \emph{Fit the geometry.} Neilan and Otus \cite{NeilanOtus2021} proved optimal convergence of Scott--Vogelius elements on isoparametric curved triangulations in two dimensions, Durst and Neilan \cite{DurstNeilan2024} treated general degree on smooth planar domains, and Chen and Liu \cite{ChenLiu2026} have very recently extended the isoparametric approach to three dimensions. (ii) \emph{Correct the boundary condition} on a mesh that does not reach the true boundary: Liu, Neilan and Otus \cite{LiuNeilanOtus2023} combine a divergence-free method with a Taylor boundary correction in two dimensions. Boundary-value corrections go back to Bramble, Dupont and Thom\'ee \cite{BrambleDupontThomee1972}; related ideas underlie the shifted boundary method \cite{MainScovazzi2018} and cut methods with boundary-value correction \cite{BurmanHansboLarson2018,BurmanHansboLarson2019Stokes}. (iii) \emph{Impose the wall condition weakly.} Nitsche's method \cite{Nitsche1971,Stenberg1995,FreundStenberg1995,JuntunenStenberg2009} and its penalty and penalty-free relatives \cite{Babuska1973,Burman2012,BoiveauBurman2016} impose Dirichlet data through a consistent boundary term, and are the standard tool of unfitted and cut finite element methods for Stokes flow \cite{BurmanHansbo2014}, including recent exactly divergence-free cut methods \cite{LiuNeilanOlshanskii2023,FrachonNilssonZahedi2024,BurmanHansboLarson2024}. Gjerde and Scott have used Nitsche's method for slip conditions in Navier--Stokes flow \cite{GjerdeScott2022}.

\subsection{Scott's zero-gradient prize and the two-dimensional answer}
Gjerde and Scott \cite{GjerdeScott2024} computed drag on a cylinder with Scott--Vogelius elements. They kept the inscribed polygon but imposed no-slip by Nitsche's method, with a penalty $\mu/h$ and no pressure term in the boundary form, and observed errors far below those of strong imposition. Scott's zero-gradient prize \cite{ScottPrize} asks for a proof that this Scott--Vogelius--Nitsche method converges, with the error $\hG^{3/2}+h^k$ one expects from the scalar theory ($\hG$ is the size of the wall faces and $h$ the bulk mesh size), and adds: \emph{if this does not hold, why not?}

The two-dimensional answer is given in the companion paper \cite{PadhiLong2D}, and its central mechanism in a short letter \cite{CavalcantePadhi2026}. The mechanism is simple. Nitsche's boundary form comes from integrating the full stress $\sigma(u,p)n=D(u)n-pn$ by parts. The method as printed in \cite[(27)--(28)]{GjerdeScott2024} keeps a viscous flux term, $\partial_nu$, with its symmetric counterpart and the penalty, but has no pressure term $-pn$. At the wall, nothing in the discrete momentum balance then opposes the pressure except the penalty $(\mu/h)\,u_h\cdot n$, so the discrete fluid leaks through the wall with normal velocity approximately $(h/\mu)(p-\pbar)$. The wall mean $\pbar$ is subtracted because an exactly divergence-free velocity has zero net flux through the wall. The printed method therefore converges only like $(h/\mu)^{1/2}$ in energy and $h/\mu$ in $H^1$ unless the wall pressure is constant. Adding back the plain traction $\ip{p_h}{v\cdot n}$ makes the discrete system singular, because constant pressures no longer see the divergence; adding the traction with the wall mean of the discrete pressure removed gives a well-posed, exactly divergence-free method, which we call $\CNS$. Its velocity coincides with that of the zero-net-flux formulation of Frachon, Nilsson and Zahedi \cite{FrachonNilssonZahedi2024}; the device is theirs. In two dimensions $\CNS$ attains $\hG^{3/2}+h^k$, the term $\hG^{3/2}$ is sharp, and the Nitsche penalty must grow with the ratio of bulk to boundary mesh size \cite{PadhiLong2D}. The constant-pressure benchmark of \cite{GjerdeScott2024,ScottPrize}, for which the leak vanishes, was analysed independently by Cavalcante \cite{Cavalcante}.

\subsection{This paper}
We answer the same question in three dimensions. The wall $\Gamma$ is a smooth closed surface, possibly with several components and not necessarily convex (the exterior of the unit ball in a box is the model case). It is replaced by a closed triangulated surface $\Gh$ whose vertices lie on $\Gamma$, and the region between $\Gh$ and a polyhedral outer boundary is meshed by straight-sided tetrahedra. The velocity is continuous piecewise $P_k$, the pressure discontinuous piecewise $P_{k-1}$, and the wall condition is imposed by Nitsche's method. We study the printed method, which we call $\GS$, and the corrected method $\CNS$.

\paragraph{Where $\hG^{3/2}$ comes from.} A wall face $F$ with vertices on $\Gamma$ is a chord of the surface: it lies at a small distance $d$ from $\Gamma$, and $d$ vanishes at the vertices. To leading order $d$ is the quadratic bump
\[
d\approx\tfrac12\sum_{i<j}\hh(e_{ij},e_{ij})\,\mu_i\mu_j
\]
in the barycentric coordinates $\mu_i$ of $F$, where $e_{ij}$ are the edge vectors and $\hh$ is the second fundamental form of $\Gamma$ (Lemma~\ref{geom:face}). The exact velocity vanishes on $\Gamma$, so on $\Gh$ it is approximately $d\,\partial_Nu$, a field of size $\hG^2$ that oscillates on the scale $\hG$. Such a trace has $H^{1/2}$-norm of order $\hG^{3/2}$, and this is the geometric term in the error. The lower bound asks the converse question: can the discrete divergence-free fields near a face see the bump? The bump is described by three numbers per face, the values of $\hh$ on the three edges, so the lower bound reduces to three \emph{edge moments} of the normal derivative on each face. The weight that results is $|\hh|\,\partial_Nu$: on flat parts of the wall the faces lie in $\Gamma$, there is no bump, and no $\hG^{3/2}$ term.

\paragraph{What is new in three dimensions.} Four things change, and they account for most of the length of the paper.
\begin{itemize}
\item \emph{Curvature enters through the second fundamental form.} In two dimensions the sagitta of a chord is a single number times the curvature. On a triangle it is a quadratic form, and the lower bound needs discrete fields that respond to all three edge weights. On Alfeld refinements with $k\ge4$ one macro-element suffices; for $k=3$ it provably does not, and the witnesses must be built on the star of an interior edge.
\item \emph{A face-mean defect.} In two dimensions the tilt between the chord normal and the normal of $\Gamma$ is odd about the chord midpoint, which is also the centre of the chord's circumcircle, and this gives a useful cancellation. On a triangle inscribed in a sphere the face mean of the tilt is proportional to the offset between centroid and circumcentre (Lemma~\ref{geom:sag}), which is nonzero on every non-equilateral face. We show that this defect changes constants but no rate.
\item \emph{No stream functions.} The two-dimensional constructions use stream functions. We replace them by a Piola--Hestenes extension of the flux $2$-form of the velocity (Lemma~\ref{geom:ext}) and by interpolate-and-correct constructions based on the discrete inf-sup condition.
\item \emph{Geometry and stability.} In surface finite elements the closest-point projection $\Gh\to\Gamma$ is assumed to be bijective; we derive it from the volume mesh. Stability of the Scott--Vogelius pair in three dimensions is known only on split meshes, and it must hold uniformly over the family of perturbed domains. On Worsey--Farin splits the discrete inf-sup condition fails for the full pressure space, as is known \cite{GuzmanLischkeNeilan2022,FabienGuzmanNeilanZytoon2022}; this is the three-dimensional counterpart of the singular vertices of two dimensions.
\end{itemize}

\subsection{Main results}
We state the results informally; precise statements, with all hypotheses, are in the sections cited. Throughout, $h$ is the bulk mesh size that appears in Scott's penalty $\mu/h$, $\hG$ is the largest diameter of a wall face, $\gamma:=\mu\hG/h$ is the effective wall penalty, $\rho$ is the ratio of $h$ to the smallest wall-face diameter, $\pbar$ is the mean of $p$ over $\Gamma$ and $G:=\norm{p-\pbar}_{L^2(\Gamma)}$ measures how far the wall pressure is from constant. The exact solution is assumed smooth up to the boundary; at the edges and corners of the outer box this is a hypothesis on the solution, not a consequence of smooth data (Remark~\ref{set:rem:corner}).

\begin{enumerate}[label=\textbf{\arabic*.},leftmargin=1.6em]
\item \emph{Geometry.} If the tetrahedral mesh is shape regular and the discrete domain does not reach the far side of the obstacle (condition (OB) of Section~\ref{sec:setting}), then the closest-point projection maps $\Gh$ homeomorphically onto $\Gamma$ and $\Gh$ is a normal graph over $\Gamma$ of height $O(\hG^2)$ (Theorems~\ref{geom:orient} and~\ref{geom:mesh}). Neither hypothesis can be dropped (Examples~\ref{ex:sliver} and~\ref{ex:filled}).
\item \emph{Stability.} On Alfeld refinements with $k\ge3$ the discrete inf-sup constant is bounded below uniformly in $h$ (Lemma~\ref{st:IS3}). On meshes whose wall faces are subdivided, such as Worsey--Farin splits, the inf-sup condition fails for the full pressure space (Proposition~\ref{st:wf}), as is known.
\item \emph{The printed method leaks.} If the wall pressure is not constant, $\GS$ converges with error exactly of order $(h/\mu)^{1/2}G$ in energy and $h/\mu$ in $H^1$. To leading order the discrete velocity satisfies the leak condition $(\mu/h)\,u_h\cdot n\approx p-\pbar$ on $\Gh$, with constant exactly $1$ in the slip and energy norms (Theorems~\ref{er:A}, \ref{er:C}, \ref{er:B}, Corollary~\ref{er:Bp}).
\item \emph{The leak limit.} For every penalty above the coercivity threshold, the scaled error $(\mu/h)(\ut-u_h)$, minus its geometric part, converges in $H^1$ to the Stokes flow $U$ driven by the Dirichlet data $-(p-\pbar)n$ on $\Gamma$, at rate $\hG^{1/2}$ (Theorem~\ref{lk:up}). On Alfeld refinements with $k\ge4$ the rate is exactly $\frac12$, with a lower bound weighted by the second fundamental form (Theorem~\ref{lk:low}). So the $H^1$ constant of $\GS$ is $\norm{\nabla U}/G$ (Corollary~\ref{lk:gen}).
\item \emph{The corrected method attains Scott's rate.} $\CNS$ is well posed and exactly divergence-free, and for bounded $\gamma\ge\gamma_2$ its $H^1$ error is at most $C(\hG^{3/2}+h^k)$ provided $\hG\ge h^2$ (Theorem~\ref{er:D}); the proviso disappears for outer data with zero discrete net flux. On Alfeld refinements this holds for every $k\ge3$ with no further hypothesis.
\item \emph{Sharpness.} For $\CNS$, and for $\GS$ with any wall pressure,
\[
\norm{\nabla(\ut-u_h)}\ \ge\ c\,\hG^{3/2}\,\bigl\|\,|\hh|\,\partial_Nu\bigr\|_{L^2(\Gamma)}-C\hG^2,
\]
on Alfeld refinements for every $k\ge3$ and on every shape-regular mesh for $k\ge9$ (Theorems~\ref{sh:lb}, \ref{sh:alfeld}, \ref{sh:k3main}, Proposition~\ref{sh:bubble}). So the $\hG^{3/2}$ term of Scott's rate cannot be improved whenever the wall shear is nonzero somewhere on a curved part of the wall.
\item \emph{Drag.} The drag, lift and torque of $\GS$ have the error $-(h/\mu)K_\psi$ to leading order, with $K_\psi$ given by a reciprocity formula (Theorem~\ref{dr:gs}); those of $\CNS$ converge at rate $\frac32$, and at rate $2$ modulo a regularity assumption for a dual problem at the outer boundary (Theorem~\ref{dr:cns}).
\item \emph{Penalty.} A penalty $\mu\ge C\rho$ suffices for coercivity. Conversely, coercivity on the divergence-free space fails for $\mu$ below a constant times $\rho$, under explicit conditions at a smallest wall face: either its Alfeld macro-element has a certified positive single-cell constant (exact certificates for a range of shapes), or the mesh is locally quasi-uniform there (a continuum witness) (Theorems~\ref{pen:single}, \ref{pen:cont}).
\item \emph{Navier--Stokes.} Under discrete linearised-stability conditions, which hold for small data and, for small $h$, on nonsingular solution branches, the upper bounds, the leak law, the leak limit and the drag statements carry over to steady Navier--Stokes flow (Corollary~\ref{ns:cor}, Theorem~\ref{ns:oseen}).
\item \emph{Computations.} A first three-dimensional implementation confirms exact incompressibility, the optimal rates on a manufactured problem and the leak law with constant tending to $1$; the sphere tests we could run are still pre-asymptotic for the rate $\hG^{3/2}$ (Section~\ref{sec:numerics}).
\end{enumerate}
The constants of the lower bounds are explicit on Alfeld refinements with $k\ge4$ (Appendix~\ref{app:constants}). Table~\ref{tab:summary} collects the results with their scope and their status.

\begin{table}[t]
\centering\footnotesize
\caption{Summary of results. ``Exact'' marks proofs that use a finite computation in exact rational or interval arithmetic, reproducible from Appendix~\ref{app:repro}. (IS3) is the uniform discrete inf-sup condition; it is a theorem on Alfeld refinements with $k\ge3$. The other named hypotheses are explained in Table~\ref{tab:hyp}.}\label{tab:summary}
\begin{tabular}{>{\raggedright\arraybackslash}p{0.27\textwidth}>{\raggedright\arraybackslash}p{0.27\textwidth}>{\raggedright\arraybackslash}p{0.23\textwidth}>{\raggedright\arraybackslash}p{0.13\textwidth}}
\toprule
Statement & Size & Scope & Status\\
\midrule
$\pi:\Gh\to\Gamma$ homeomorphism & normal graph, height $O(\hG^2)$ & shape regular $+$ (OB) & proved\\
uniform discrete inf-sup & $\beta$ independent of $h$ & Alfeld, $k\ge3$ & proved\\
$\GS$, energy error & $\asymp(h/\mu)^{1/2}G$ & (IS3) & proved\\
$\GS$, $H^1$ error & $\asymp h/\mu$, constant $\norm{\nabla U}/G$ & (IS3), all $\gamma\ge\gamma_0$ & proved\\
leak limit & rate exactly $\hG^{1/2}$ & upper: (IS3); lower: Alfeld, $k\ge4$ & proved; lower bound exact\\
$\CNS$, $H^1$ error & $\le C(\hG^{3/2}+h^k)$ & (IS3), bounded $\gamma\ge\gamma_2$ & proved\\
lower bound $c\,\hG^{3/2}$ & weight $|\hh|\,\partial_Nu$ & Alfeld $k\ge3$; any mesh $k\ge9$ & proved; exact\\
$\GS$ drag, lift, torque & error $-(h/\mu)K_\psi+\dots$ & (IS3) & proved\\
$\CNS$ drag, lift, torque & rate $\frac32$; rate $2$ & (IS3) & proved; rate $2$ modulo (E$_Z$)\\
penalty threshold $\mu^\ast$ & $c\rho\le\mu^\ast\le C\rho$ & lower: certified cell or (LQ) & proved; single cell exact\\
Navier--Stokes & as for Stokes & (L), (L$^\ast$), (L$_D$), (L$^\dagger$) & conditional\\
Worsey--Farin & inf-sup fails for full $\Pih$ & subdivided wall faces & known; exact count\\
computations & rates on a box, leak law & $k=3,4$ & numerical\\
\bottomrule
\end{tabular}
\end{table}

\subsection{Related work and novelty}\label{sec:novelty}
We list what is known, what we take from others, and what we believe is new.

\paragraph{Known, and used.} Stability of the Scott--Vogelius pair on Alfeld refinements is due to Zhang \cite{Zhang2005} and Guzm\'an and Neilan \cite{GuzmanNeilan2018}; our uniform-in-$h$ version (Lemma~\ref{st:IS3}) combines the local result of \cite[Theorem~3.1]{GuzmanNeilan2018}, Bernardi--Raugel face fluxes \cite{BernardiRaugel1985} and the continuity of the continuous inf-sup constant under Lipschitz perturbations of the domain due to Bernardi, Costabel, Dauge and Girault \cite{BCDG2016}. Chen and Liu use the same combination for polyhedral approximations of curved domains \cite[Lemma~3.9]{ChenLiu2026}, and our argument follows theirs in structure; we claim no novelty for it. The mean-free traction of $\CNS$ is the zero-net-flux device of Frachon, Nilsson and Zahedi \cite{FrachonNilssonZahedi2024}. The failure of the inf-sup condition for the full pressure space on Worsey--Farin splits is known \cite{GuzmanLischkeNeilan2022,FabienGuzmanNeilanZytoon2022}; we only record a short proof and an exact count for our setting. The order $h^{3/2}$ for inscribed polygonal boundaries is classical for scalar problems \cite{BergerScottStrang1972,Scott1975}. The structure of our lift theorem (a proper local homeomorphism with one sheet) is that of Amenta, Choi, Dey and Leekha \cite{AmentaChoiDeyLeekha2000}.

\paragraph{Closest recent work.} Chen and Liu \cite{ChenLiu2026} prove the optimal, pressure-robust rate $h^k$ for Piola-mapped Scott--Vogelius elements with $k\ge3$ on Alfeld splits of \emph{isoparametric} curved tetrahedral meshes, with strongly imposed Dirichlet data and a commuting-interpolant load. That is the right method if one is willing to build curved elements. Our setting is the cheap one that Scott's question is about: straight-sided, affinely mapped tetrahedra, an inscribed triangulated wall, standard Nitsche terms, and no curved element maps, multipliers or boundary corrections. Its best possible rate is $\hG^{3/2}+h^k$, and our contribution is to show that $\CNS$ attains this rate and that it cannot be improved; we do not compete with \cite{ChenLiu2026} on rate. Liu, Neilan and Otus \cite{LiuNeilanOtus2023} reach order $h^k$ in two dimensions on a mesh inside the curved domain with a mean-free boundary multiplier, a zero-flux constraint and a Taylor boundary correction; the divergence-free cut methods \cite{LiuNeilanOlshanskii2023,FrachonNilssonZahedi2024,BurmanHansboLarson2024} are unfitted, and their analyses concern that setting. None of these works treats the Nitsche method on inscribed polyhedra, the leak of the printed method, or lower bounds. Cavalcante \cite{Cavalcante} analyses the two-dimensional constant-pressure benchmark. The two-dimensional versions of our results are in \cite{PadhiLong2D}.

\paragraph{New, to our knowledge.}
\begin{enumerate}[label=(\alph*)]
\item The lift theorem: the closest-point projection $\Gh\to\Gamma$ is a homeomorphism as a \emph{consequence} of shape regularity of the volume mesh and the global condition (OB), with examples showing that neither can be dropped (Section~\ref{sec:geometry}).
\item The three-dimensional leak law of the printed method, with exact constants, and the leak limit uniformly in the penalty, with the exact rate $\hG^{1/2}$ and a lower bound weighted by the second fundamental form (Sections~\ref{sec:error} and~\ref{sec:leak}).
\item The analysis of $\CNS$ on inscribed polyhedra: well-posedness and the rate $\hG^{3/2}+h^k$, unconditional on Alfeld refinements for $k\ge3$.
\item The sharp lower bound $c\,\hG^{3/2}\norm{|\hh|\,\partial_Nu}$, with the reduction to three edge moments per face, the exact witnesses for $k\ge4$, the proof that one macro-element cannot work for $k=3$, and the edge-star construction for $k=3$, including saddle and parabolic points (Section~\ref{sec:sharp} and Appendix~\ref{app:k3}).
\item Drag, lift and torque errors of both methods in three dimensions, and the Navier--Stokes and Oseen versions.
\item Necessity of a penalty proportional to $\rho$ in three dimensions, with exact single-cell certificates and a continuum witness; explicit constants for the lower bounds.
\end{enumerate}
Items (b)--(f) have two-dimensional counterparts in \cite{PadhiLong2D}. Where a two-dimensional proof carries over with only notational changes, we give the three-dimensional replacement of each dimension-dependent ingredient and refer to \cite{PadhiLong2D} for the unchanged algebra.

\subsection{Organisation and conventions}
Section~\ref{sec:setting} fixes the setting and collects the named hypotheses in Table~\ref{tab:hyp}. Section~\ref{sec:geometry} studies inscribed surfaces and proves the lift theorem, Section~\ref{sec:stability} the stability results, Section~\ref{sec:error} the error analysis of both methods, Section~\ref{sec:sharp} the lower bounds, Section~\ref{sec:leak} the leak limit and drag, Section~\ref{sec:penalty} the penalty threshold, Section~\ref{sec:ns} the Navier--Stokes extension and Section~\ref{sec:numerics} the computations. Section~\ref{sec:open} discusses what remains open. Appendix~\ref{app:k3} contains the degree-three witnesses, Appendix~\ref{app:constants} the explicit constants and the details of the computer-assisted certificates, and Appendix~\ref{app:repro} the command behind every computed number.

A statement marked \status{computer-assisted, exact} relies on a finite computation in exact rational arithmetic (or, where stated, interval arithmetic) whose output is a certificate that can be rechecked independently; \status{numerical} marks floating-point evidence that is not part of any proof. Constants $C,c$ never depend on $h,\hG,\mu,\gamma,\rho$ unless this is displayed.

\section{Setting, methods and assumptions}\label{sec:setting}

\subsection{Domain}
\begin{assumption}[Domain]\label{set:dom}
$\Omega\subset\R^3$ is a bounded connected Lipschitz domain with $\partial\Omega=\Gamma\sqcup\Sigma$. The wall $\Gamma=\Gamma_1\sqcup\dots\sqcup\Gamma_m$ is a union of closed, connected, embedded surfaces of class $C^{k+3}$, each a whole component of $\partial\Omega$. The outer boundary $\Sigma\ne\emptyset$ is a finite union of closed polyhedral surfaces (whole components of $\partial\Omega$) and carries strongly imposed Dirichlet data.
\end{assumption}

The model case is $\Omega=R\setminus\overline D$ with $R=(-L,L)^3$, $L>1$, and $D$ the open unit ball.

$N$ denotes the unit normal of $\Gamma$ pointing into $\Omega$ and $n:=-N$ the outward normal of $\Omega$ on $\Gamma$. Let $d$ be the signed distance to $\Gamma$, positive in $\Omega$, $\pi$ the closest-point projection, and $\hh_y(X,Y):=-\ip{D_XN}{Y}$ the second fundamental form with respect to $N$, so that $\hh>0$ where $\Omega$ is locally convex; for the exterior of the unit ball, $\hh=-I$. Put $\kappa_\infty:=\sup_\Gamma|\hh|$. We use two standard facts about the fixed surface $\Gamma$.
\begin{enumerate}
\item[(T)] There is $r>0$ with $r\kappa_\infty\le\frac12$ such that $\Psi(y,s):=y+sN(y)$ is a $C^{k+2}$ diffeomorphism of $\Gamma\times(-r,r)$ onto a tube $\mathcal T_r$, with $d(\Psi(y,s))=s$, $\pi(\Psi(y,s))=y$ and $\nabla d=N\circ\pi$. The volume element is $dx=J\,dA\,ds$ with $J=\det(I-sS_y)\in[\frac14,\frac94]$, $S$ the shape operator. $\mathcal T_r$ has positive distance from $\Sigma$, and the tubes around different $\Gamma_i$ are disjoint.
\item[(Gr)] There is $a_0\le r/4$ with $\kappa_\infty a_0\le\frac1{40}$ such that, for each $y\in\Gamma$, in the frame $(y;\tau_1,\tau_2,N(y))$ with coordinates $(\xi,\eta)\in\R^2\times\R$, the set $\Gamma\cap(-2a_0,2a_0)^3$ is the graph of $f_y\in C^{k+3}$ with $f_y(0)=0$, $\nabla f_y(0)=0$, $D^2f_y(0)=\hh_y$, $|\nabla f_y(\xi)|\le2\kappa_\infty|\xi|$ and $|f_y(\xi)|\le\kappa_\infty|\xi|^2$, and $\Omega\cap(-2a_0,2a_0)^3=\{\eta>f_y(\xi)\}$.
\end{enumerate}

\subsection{Inscribed surface and mesh}
\begin{assumption}[Inscribed surface and mesh]\label{set:mesh}\leavevmode
\begin{enumerate}[label=(G\arabic*)]
\item $\Gh$ is a closed triangulated surface (each edge lies in exactly two faces) whose vertices lie on $\Gamma$. Its faces have diameter in $[c_0\hG,\hG]$ and minimal angle at least $\theta_0>0$.
\item[(G3)] $\hG\le h_\ast$, a constant depending on $\Gamma$, $c_0$, $\theta_0$ and the shape regularity below.
\end{enumerate}
\begin{enumerate}[label=(M0$^3$)]
\item $\Oh$ is a bounded connected open polyhedral set with $\partial\Oh=\Sigma\sqcup\Gh$, and $\Th$ is a conforming, shape-regular tetrahedral mesh of $\Oh$ (inradius $\ge c_s\diam T$) whose boundary faces on $\Gh$ are exactly the faces of $\Gh$. We write $\FG$ for this set of faces.
\end{enumerate}
\begin{enumerate}[label=(OB)]
\item For each $i$, the parallel surface $\Gamma_i^-:=\{y-\frac r2N(y):y\in\Gamma_i\}$ on the obstacle side is not contained in $\Oh$.
\end{enumerate}
\end{assumption}

(OB) holds whenever $\Oh$ agrees with $\Omega$ away from a thin neighbourhood of $\Gamma$, for example for every mesh that only modifies $\Omega$ near $\Gamma$. In the model case, $P_h$ is a convex polyhedron with triangular faces and vertices on the unit sphere, and $\Oh=R\setminus P_h$; then $P_h\supset B(0,1-C\hG^2)\supset\Gamma^-$, so (OB) holds. The two-dimensional analogue of (G1) and (OB) is discussed in \cite{PadhiLong2D}. Section~\ref{sec:geometry} shows that (G1), (M0$^3$) and (OB) imply that $\pi|_{\Gh}$ is a homeomorphism onto $\Gamma$; in surface finite elements this lift property is usually assumed.

Put
\[
h:=\max_{T\in\Th}\diam T,\qquad \hG:=\max_{F\in\FG}\diam F,\qquad \rho:=\frac h{\min_{F\in\FG}\diam F},\qquad \gamma:=\frac{\mu\hG}h .
\]
Here $h$ is the $h$ in Scott's penalty $\mu/h$, and $\gamma$ is the effective wall penalty, $\mu/h=\gamma/\hG$. By (G1), $\rho/h\le1/(c_0\hG)$. No quasi-uniformity is assumed.

\begin{definition}[Alfeld refinement]\label{set:alfeld}
$\Th$ is an \emph{Alfeld refinement} if it is obtained from a conforming shape-regular macro mesh $\Th^M$ of $\Oh$ by splitting every macro-tetrahedron into four at its barycentre. The faces of the macro-tetrahedra on $\Gh$ are not subdivided, so each $F\in\FG$ is a face of exactly one macro-tetrahedron $T_F$ and of one of its four sub-tetrahedra.
\end{definition}

\textbf{(D)} $\ft$ is a smooth bounded extension of $f$ to a fixed neighbourhood of $\overline\Omega$, and $h\le h_\ast$.

\subsection{Spaces and forms}
Fix $k\ge3$. $\Vh$ is the space of continuous piecewise-$P_k$ vector fields on $\Th$, $\VR:=\{v\in\Vh:v|_\Sigma=0\}$ and $\VO:=\{v\in\Vh:v|_{\partial\Oh}=0\}$. $\Pih$ is the space of discontinuous piecewise-$P_{k-1}$ functions and $\Zh:=\{v\in\VR:\div v=0\}$. With $g:=u|_\Sigma$ and $\gI$ its $P_k$ Lagrange interpolant, $\Wh:=\{z\in\Vh:z|_\Sigma=\gI,\ \div z=0\}$.

On $\Gh$, $n$ is the outward unit normal of $\Oh$ (on the face $F$ we also write $n_F$), $\dn v:=(\nabla v)n$ and $\ip\cdot\cdot$ is the $L^2(\Gh)$ pairing. Let
\begin{align*}
a_h(u,v)&:=\tfrac12(Du,Dv),\qquad Du:=\nabla u+(\nabla u)^{\mathsf T},\\
N_h(u,v)&:=a_h(u,v)-\ip{\dn u}{v}-\ip{u}{\dn v}+\frac\mu h\ip uv .
\end{align*}

\begin{definition}[Method $\GS$ {\cite[(28)]{GjerdeScott2024}}, as printed]\label{set:GS}
Find $u_h\in\Vh$ with $u_h|_\Sigma=\gI$ and $p_h\in\Pih$ such that
\[
N_h(u_h,v)-(p_h,\div v)=(\ft,v)\quad\forall v\in\VR,\qquad (q,\div u_h)=0\quad\forall q\in\Pih .
\]
\end{definition}
The transcription conventions (the sign of the pressure, and the test space $\VR$) are those of \cite{PadhiLong2D}.

\begin{definition}[Method $\CNS$]\label{set:CNS}
As $\GS$, with momentum row
\[
N_h(u_h,v)-(p_h,\div v)+\ip{p_h-\pbG(p_h)}{v\cdot n}=(\ft,v),\qquad \pbG(q):=|\Gh|^{-1}\int_{\Gh}q .
\]
\end{definition}
The mean is over all of $\Gh$, also when $\Gamma$ has several components; per-component means are inconsistent (Remark~\ref{er:rem:global}). The added term appears in the momentum row only, so the velocity stays exactly divergence-free. Its velocity is that of the zero-net-flux formulation of \cite[Section~5.1]{FrachonNilssonZahedi2024} (see \cite{PadhiLong2D} for the transcription; the argument is dimension-free).

\begin{enumerate}[label=(IS3)]
\item There is $\beta>0$, independent of $h$, such that for every $q\in\Pih\cap L^2_0(\Oh)$ there is $v\in\VO$ with $\div v=q$ and $\norm{\nabla v}\le\beta^{-1}\norm q$.
\end{enumerate}
On Alfeld refinements with $k\ge3$, (IS3) is a theorem (Lemma~\ref{st:IS3}). On meshes whose wall faces are subdivided, such as Worsey--Farin splits, it is false as formulated (Proposition~\ref{st:wf}). On other meshes it is an assumption; see Remark~\ref{st:rem:status}.

\subsection{Exact solution, extensions and norms}
$(u,p)$ is a Stokes solution, $-\Delta u+\nabla p=f$, $\div u=0$ in $\Omega$, $u|_\Gamma=0$, $u|_\Sigma=g$, with $u\in C^{k+1}(\overline\Omega)$ and $p\in C^k(\overline\Omega)$. Lemma~\ref{geom:ext} provides extensions $\ut\in C^{k+1}$ and $\pt\in C^k$ to a fixed neighbourhood $\mathcal O\supset\overline\Oh$ with $\div\ut=0$. The defect
\[
F:=\ft+\Delta\ut-\nabla\pt
\]
vanishes on $\Omega$ and is supported in $\overline{\Oh\setminus\Omega}$.

\begin{remark}[Edges and corners of the outer box]\label{set:rem:corner}
In the model case $\Sigma=\partial R$ has edges and corners, and near them the smoothness of $(u,p)$ is a hypothesis on the \emph{solution}, not a consequence of smooth data. Along an edge of $R$ the singular exponents include those of the plane Stokes problem in a right angle, the roots of $\sin^2(\lambda\pi/2)=\lambda^2$ with $\operatorname{Re}\lambda>1$, the first being $\lambda_1=2.739593\pm1.119025\,i$ (computed for the companion paper \cite{PadhiLong2D}). So for generic smooth, compatible data one can only expect $u\in H^s$ near an edge for $s<1+2.7396$, unless the singular coefficients vanish; the corners of $R$ add vertex singularities, whose exponents we have not computed \cite{Dauge1989}. The hypothesis $u\in C^{k+1}(\overline\Omega)$ holds for the manufactured solutions of Section~\ref{sec:numerics}. For data-driven problems on meshes that are not graded at the edges of $R$, the terms $h^k$ in the $H^1$ bounds, and every other term that uses the smoothness of $u$ near $\Sigma$, should be read with $h^k$ replaced by $h^{\min(k,2.7396)-\epsilon}$. Near $\Gamma$ nothing changes: the wall is smooth and $\mathcal T_r$ has positive distance from $\Sigma$. The same caveat applies to the global regularity used in Remark~\ref{ns:rem} (Stokes $H^2\times H^1$ on $\Omega$, which for the box minus a ball we can only reduce, by a partition of unity and a Bogovski\u{\i} correction, to the $H^2$ theorem for convex polyhedra attributed to \cite{Dauge1989}, whose statement we have not checked), to the approximation term $E^0_U$ of Section~\ref{sec:leak}, and to assumption (E$_Z$) of Section~\ref{sec:leak}.
\end{remark}

For fields with traces on $\Gh$,
\[
\tnorm v^2:=\norm{\nabla v}^2_{\Oh}+\frac\mu h\norm v^2_{\Gh}+\sum_{F\in\FG}\diam F\,\norm{\dn v}^2_{L^2(F)} .
\]
We write $\pbar:=|\Gamma|^{-1}\int_\Gamma p$ and $G:=\norm{p-\pbar}_{L^2(\Gamma)}$. $C$ and $c$ depend on $k$, the constants in (G1) and (M0$^3$), $\Omega$, and norms of $(u,p,\ut,\pt)$; they never depend on $h,\hG,\mu,\gamma,\rho$ unless displayed. $C_p$ may also depend on $\norm{p-\pbar}_{W^{1,\infty}(\Gamma)}$.

\section{Geometry of inscribed surfaces}\label{sec:geometry}

Every estimate in this paper starts from two geometric facts: on each wall face, the distance to $\Gamma$ is a small quadratic bump, and the closest-point projection identifies $\Gh$ with $\Gamma$. The first is a local computation (Lemma~\ref{geom:face}); on the sphere it also exposes the face-mean defect that distinguishes three dimensions from two (Lemma~\ref{geom:sag}). The second is global, and it is where the volume mesh enters: a surface can be locally close to $\Gamma$ and still fold over itself, and we show that a shape-regular tetrahedral mesh of the fluid region rules this out (Theorems~\ref{geom:orient} and~\ref{geom:mesh}). The section ends with the consequences used later: uniform Korn and trace inequalities on the family $\Oh$, and a divergence-free extension of the exact solution across $\Gamma$.

\subsection{Faces}
For a face $F\in\FG$ with vertices $z_1,z_2,z_3\in\Gamma$ we write $\mu_1,\mu_2,\mu_3$ for its barycentric coordinates, $e_{ij}:=z_i-z_j$, $\ell:=\diam F$, $x_F$ for its centroid, $y_F:=\pi(x_F)$, $P_y$ for the orthogonal projection onto $T_y\Gamma$ and $\hh_F:=\hh_{y_F}\circ P_{y_F}$. For $x\in F$ put $x^\ast:=\pi(x)$.

\begin{lemma}[Faces on a general surface]\label{geom:face}
Assume (G1) and (G3). Then for every $F\in\FG$ and $x\in F$:
\begin{enumerate}[label=(\alph*)]
\item $d(x)=\frac12\sum_{i<j}\hh_F(e_{ij},e_{ij})\,\mu_i\mu_j+O(\ell^3)$; in particular $|d(x)|\le C\ell^2$.
\item There is $\sigma_F\in\{\pm1\}$, independent of $x\in F$, with $|\sigma_Fn_F-n(x^\ast)|\le C\ell$ (where $n=-N$); hence $|n_F\cdot N(x^\ast)|\ge1-C\ell^2$.
\item $\pi|_F$ is a $C^2$ diffeomorphism onto $\pi(F)$, with area Jacobian $1+O(\ell^2)$.
\item With $w_F:=(\hh_F(e_{ij},e_{ij}))_{i<j}\in\R^3$ and $|\hh|$ the operator norm, $|w_F|\ge c(\theta_0)\,\ell^2\,|\hh_{y_F}|-C\ell^3$.
\end{enumerate}
\end{lemma}

\begin{proof}
Use the frame (Gr) at $y:=z_1$. Then $z_j=(\xi_j,f(\xi_j))$ with $|\xi_j|\le\ell$; the vertical parts are at most $\kappa_\infty\ell^2$, so the projected triangle $\tilde F=\conv(\xi_j)$ has minimal angle at least $\theta_0/2$ for $h_\ast$ small, and $F$ is the graph over $\tilde F$ of the linear interpolant $L:=I_1f$, with the same barycentrics.

(a) Write $f=\frac12\xi^{\mathsf T}H\xi+R_3$ with $H=\hh_{z_1}$ and $|R_3|\le C\ell^3$ on $\tilde F$. For a quadratic form, $I_1Q-Q=\sum_{i<j}\mu_i\mu_j\,Q(\xi_i-\xi_j)$; this is the polarised variance identity $\sum_i\mu_i|\xi_i-o|^2-|\xi-o|^2=\sum_{i<j}\mu_i\mu_j|\xi_i-\xi_j|^2$. For $x=(\xi,L(\xi))$ we have $d(x)=(L-f)(\xi)N_3+O((L-f)^2)=(L-f)(\xi)+O(\ell^4)$, since $N_3=(1+|\nabla f|^2)^{-1/2}=1+O(\ell^2)$. Finally $\xi_i-\xi_j=P_{z_1}e_{ij}$, and replacing $(z_1,P_{z_1})$ by $(y_F,P_{y_F})$ changes $\hh(e,e)$ by $O(\ell)|e|^2$ because $\Gamma\in C^3$.

(b) $n_F\parallel(\nabla L,-1)$ and $N(x^\ast)\parallel(-\nabla f(\xi^\ast),1)$. On $\tilde F$, $|\nabla L-\nabla f|\le C\ell/\sin(\theta_0/2)$ (linear interpolation), and $|\xi^\ast-\xi|\le C\ell^2$.

(c) $D\pi(x)=(I-d\,S)^{-1}P_{x^\ast}$. Restricted to the plane of $F$ its Jacobian is $|n_F\cdot N(x^\ast)|(1+O(d))=1+O(\ell^2)$ by (a) and (b). Injectivity on $F$: $\pi|_F$ is a $C^1$-small perturbation of the vertical projection, which is injective.

(d) The map $H\mapsto(H(u_{ij},u_{ij}))_{i<j}$ with $u_{ij}:=P_{y_F}e_{ij}/|P_{y_F}e_{ij}|$, from symmetric forms on $T_{y_F}\Gamma$ to $\R^3$, is invertible with inverse bounded by $C(\theta_0)$: the three directions are pairwise non-parallel with angles bounded below, and three distinct doubled angles give three non-collinear points on a circle. Also $|P_{y_F}e_{ij}|\ge c\ell$.
\end{proof}

For the model case the sagitta is exactly a circumsphere function, and the face mean of the normal tilt measures the offset of the centroid from the circumcentre.

\begin{lemma}[Faces inscribed in the unit sphere]\label{geom:sag}
Let $F$ be a face of a convex polyhedron inscribed in the unit sphere, with circumcentre $o_F$, circumradius $r_F$, outward (from the centre) unit normal $\nu_F$, and $d_F:=(1-r_F^2)^{1/2}$. Write $\ell_{ij}:=|e_{ij}|$, and let $x_c$ be the centroid.
\begin{enumerate}[label=(\alph*)]
\item For $x\in F$, $1-|x|=\frac12(r_F^2-|x-o_F|^2)+\delta(x)=\frac12\sum_{i<j}\ell_{ij}^2\mu_i\mu_j+\delta(x)$ with $0\le\delta\le r_F^4/4$.
\item $n(x^\ast)-n_F=-(x-o_F)+O(r_F^2)$, where $n(x^\ast)=-x^\ast$ and $n_F=-\nu_F$. The leading term is tangential to $F$ and vanishes at the circumcentre.
\item $|F|^{-1}\int_F(n(x^\ast)-n_F)\,dA=-(x_c-o_F)+O(r_F^2)$. The leading term vanishes if and only if $F$ is equilateral.
\item With $e_{12},e_{23},e_{31}$ the cyclically oriented edge vectors, $\sum_{\rm cyc}\ell_{ij}^2e_{ij}=\pm12|F|\,R_{\pi/2}(x_c-o_F)$, where $R_{\pi/2}$ is the rotation by $\pi/2$ in the plane of $F$.
\end{enumerate}
\end{lemma}

\begin{proof}
(a) With $s:=r_F^2-|x-o_F|^2\in[0,r_F^2]$ we have $|x|^2=d_F^2+|x-o_F|^2=1-s$, and $1-\sqrt{1-s}\in[\frac s2,\frac s2+\frac{s^2}4]$ for $s\le\frac12$. For $x=\sum\mu_iz_i$, $\sum_i\mu_i|z_i-o_F|^2-|x-o_F|^2=\sum_{i<j}\mu_i\mu_j\ell_{ij}^2$, and $|z_i-o_F|=r_F$.
(b) $n(x^\ast)-n_F=-(o_F+y)/|x|+\nu_F=-y/|x|+\nu_F(1-d_F/|x|)$ with $y:=x-o_F\perp\nu_F$, $|x|^{-1}=1+O(r_F^2)$ and $0\le1-d_F/|x|\le1-d_F$. (c) integrates (b); centroid and circumcentre coincide exactly for equilateral triangles. (d) With $a+b+c=0$ the cyclic edge vectors, $|a|^2a+|b|^2b+|c|^2c=(|a|^2-|c|^2)a+(|b|^2-|c|^2)b$; the identity with $x_c-o_F$ is an elementary computation in the plane (we checked it symbolically and on random triangles).
\end{proof}
In two dimensions the circumcentre of a chord is its midpoint, so the tilt of a chord is odd about its midpoint; this is the cancellation used in \cite{PadhiLong2D}. By (c) it is lost on every non-equilateral face. We never use it (Lemma~\ref{er:G}).

\subsection{The closest-point projection is a homeomorphism}\label{geom:sec:lift}
Lemma~\ref{geom:face} is a statement about single faces. It does not exclude folds: two faces sharing an edge can have $\sigma_F=-\sigma_{F'}$ and project to the same side of the image of the edge. In surface finite elements the lift $\pi|_{\Gh}$ is therefore assumed to be bijective \cite{Dziuk1988,DziukElliott2013}; in discrete differential geometry, polyhedral surfaces that are normal graphs are assumed as well \cite{HildebrandtPolthierWardetzky2006,MorvanThibert2004}. We show that in our setting it is a consequence of the volume mesh.

For $y\in\Gamma$ let $\ell_y:=\{\Psi(y,s):|s|\le r/2\}$ be the normal segment. Let $S_1$ be the union of the edges of $\Gh$ and $V:=\Gamma\setminus\pi(S_1)$, which is open and dense by Lemma~\ref{geom:face}(c). For $y\in V$, the set $\ell_y\cap\Gh$ is finite and consists of interior points of faces; let $k(y)$ be its cardinality. By Lemma~\ref{geom:face}(a), $\Gh\subset\{|d|\le C_1\hG^2\}\subset\mathcal T_{r/4}$, and a normal segment meets a face at most once and transversally (Lemma~\ref{geom:face}(b)).

\begin{lemma}[Crossing rule]\label{geom:cross}
Let $y\in V$ and $\ell_y\cap\Gh=\{\Psi(y,s_j)\}_{j=1}^k$ with $s_1<\dots<s_k$, the $j$-th point interior to the face $F_j$. On each open interval of $(-\frac r2,\frac r2)\setminus\{s_j\}$ the property $\Psi(y,s)\in\Oh$ is constant. Crossing $s_j$ upwards goes from outside to inside $\Oh$ if $\sigma_{F_j}=+1$, and from inside to outside if $\sigma_{F_j}=-1$.
\end{lemma}

\begin{proof}
$\ell_y$ does not meet $\Sigma$, so the only points of $\partial\Oh$ on it are the $\Psi(y,s_j)$. Near an interior point of $F_j$, $\Oh$ is the open half-space side into which $-n_{F_j}$ points, and $\sigma_{F_j}=+1$ means that this is the side of increasing $s$.
\end{proof}

The proof has two parts. First, if every face is oriented like $\Gamma$ (condition (OR) below), a normal segment can cross $\Gh$ only from outside to inside $\Oh$, hence at most once; a covering-space argument then upgrades this to a global homeomorphism. Second, the tetrahedra touching $\Gh$ force (OR), provided $\Oh$ does not reach the far side of the obstacle.

\begin{theorem}[Orientation implies the lift]\label{geom:orient}
Assume (G1), (G3), that $\Oh$ is a bounded open polyhedral set with $\partial\Oh=\Sigma\sqcup\Gh$ lying on one side of each face, and
\begin{enumerate}[label=(OR)]
\item $\sigma_F=+1$ for every $F\in\FG$, i.e.\ $n_F\cdot n(\pi(x))>0$ on $F$.
\end{enumerate}
Then for each $i$, $\pi$ maps $\Gamma_h^{(i)}:=\Gh\cap\mathcal T_r(\Gamma_i)$ homeomorphically onto $\Gamma_i$, piecewise as a $C^2$ diffeomorphism, and $n_F\cdot n(\pi(x))\ge1-C\hG^2$ on every face.
\end{theorem}

\begin{proof}
\emph{Step 1: $k\le1$ on $V$.} By Lemma~\ref{geom:cross} and (OR) every crossing goes from outside to inside. If $k(y)\ge2$, the status just above $s_1$ is inside and just below $s_2$ is outside, which contradicts constancy on $(s_1,s_2)$.

\emph{Step 2: faces.} $\pi$ is injective on each face (Lemma~\ref{geom:face}(c)).

\emph{Step 3: edges.} Let $x_0$ be interior to the edge $E=F\cap F'$ and $\gamma:=\pi(E)$, a $C^2$ arc. $\pi$ maps half-disc neighbourhoods of $x_0$ in $F$ and in $F'$ onto one-sided neighbourhoods of $\pi(x_0)$ bounded by $\gamma$. If they lay on the same side of $\gamma$, they would overlap in an open set, which contains a point $y\in V$ with a preimage in the interior of $F$ and one in the interior of $F'$, so $k(y)\ge2$. Hence they lie on opposite sides, and $\pi$ is injective near $x_0$ in $F\cup F'$, a neighbourhood of $x_0$ in $\Gh$.

\emph{Step 4: vertices.} Let $v$ be a vertex and $St(v)$ its closed star. On $St(v)$, $|\pi(x)-v|\ge\frac12|x-v|$ (segments stay inside faces), and $\dist(v,\partial St(v))\ge c(\theta_0)c_0\hG$. Fix $\varepsilon<\frac14c(\theta_0)c_0\hG$ such that $B:=\{y\in\Gamma:0<|y-v|<\varepsilon\}$ is a punctured topological disc, and let $A:=\{x\in St(v):0<|\pi(x)-v|<\varepsilon\}$. By Steps~2--3 and invariance of domain, $\pi|_A:A\to B$ is a local homeomorphism; it is proper, since preimages of compact sets are closed in the compact $St(v)$ and contained in $A$. A proper local homeomorphism onto a connected, locally compact Hausdorff space is a finite covering \cite[Theorem~4.22]{Forster1981}. Its number of sheets is at most $k(y)\le1$ for $y\in B\cap V$, so it is one, and $\pi$ is injective near $v$.

\emph{Step 5: global.} By Steps~2--4 and invariance of domain, $\pi:\Gamma_h^{(i)}\to\Gamma_i$ is a local homeomorphism. It is proper ($\Gamma_h^{(i)}$ is compact), so its image is open and closed, hence all of the connected $\Gamma_i$, and it is a covering with a constant number of sheets. That number is $k(y)\le1$ at any $y\in V\cap\Gamma_i$. A continuous bijection of a compact space onto a Hausdorff space is a homeomorphism. The normal bound is Lemma~\ref{geom:face}(b) with $\sigma_F=+1$.
\end{proof}

\begin{corollary}[Normal graph]\label{geom:graph}
Under the assumptions of Theorem~\ref{geom:orient}, $\eta_h:=d\circ(\pi|_{\Gh})^{-1}$ is Lipschitz and piecewise $C^2$ on $\Gamma$, $\Gh=\{\Psi(y,\eta_h(y)):y\in\Gamma\}$, $\norm{\eta_h}_\infty\le C\hG^2$, $\norm{\nabla_\Gamma\eta_h}_\infty\le C\hG$, and
\[
\Oh\cap\mathcal T_{r/2}=\{x\in\mathcal T_{r/2}:d(x)>\eta_h(\pi(x))\}.
\]
In particular $\Gamma_i^-\cap\Oh=\emptyset$ and $\Gamma_i^+:=\{y+\frac r2N(y)\}\subset\Oh$, so (OB) holds. Moreover $|\Omega\triangle\Oh|\le C\hG^2$, and both $\Omega\setminus\Oh$ and $\Oh\setminus\Omega$ lie in $\{|d|\le C\hG^2\}$.
\end{corollary}

\begin{proof}
The bounds on $\eta_h$ follow from Lemma~\ref{geom:face}: on $\pi(F)$, $\nabla_\Gamma\eta_h=D(\pi|_F)^{-\mathsf T}\nabla_Fd$ and $|\nabla_Fd|=|N-(N\cdot n_F)n_F|\le C\hG$. For $y\in V$ the single crossing of $\ell_y$ is at $s=\eta_h(y)$ and goes from outside to inside, so $\Oh\cap\ell_y=\{\Psi(y,s):\eta_h(y)<s<\frac r2\}$. The two sets in the display are open in $\mathcal T_{r/2}$, disjoint from $\Gh$, have $\Gh$ as relative boundary, and agree on the dense set $\bigcup_{y\in V}\ell_y\setminus\Gh$; two such sets coincide ($U_1\setminus\overline{U_2}$ is open and misses a dense set). For $\Gamma_i^-$: it does not meet $\partial\Oh$, and points of $\mathcal T_{r/2}$ close to it have $d<\eta_h\circ\pi$, so no point of $\Gamma_i^-$ lies in the open set $\Oh$. The same argument gives $\Gamma_i^+\subset\Oh$.
\end{proof}

It remains to derive (OR). The idea is that each wall face carries a tetrahedron of $\Th$ whose inradius is comparable to the face, and the incentre of that tetrahedron lies on the fluid side of the face at a distance much larger than the sagitta; so the face cannot be oriented against $\Gamma$ unless the fluid region extends to the obstacle side of $\Gamma$.

\begin{theorem}[The volume mesh implies the orientation]\label{geom:mesh}
Assume (G1), (G3), (M0$^3$) and (OB). Then (OR) holds, and hence the conclusions of Theorem~\ref{geom:orient} and Corollary~\ref{geom:graph}.
\end{theorem}

\begin{proof}
\emph{(OB) is global on each component.} $\Gamma_i^-$ is connected and disjoint from $\partial\Oh$, so it lies entirely inside or entirely outside $\Oh$; by (OB), $\Gamma_i^-\cap\Oh=\emptyset$.

\emph{Each face sees the side it faces.} Let $T\in\Th$ be the tetrahedron with face $F$, $c_T$ its incentre and $\varrho_T\ge c_s\diam T\ge c_sc_0\hG$ its inradius. Then $c_T=t-\varrho_Tn_F$, with $t\in F$ the point where the insphere touches $F$, and $c_T\in\mathcal T_{r/4}$. By Lemma~\ref{geom:face}(b), $n_F=-\sigma_FN(\pi t)+O(\hG)$, and with $\nabla d=N\circ\pi$,
\[
d(c_T)=d(t)-\varrho_T\int_0^1N(\pi(t-\tau\varrho_Tn_F))\cdot n_F\,d\tau=\sigma_F\varrho_T+O(\hG^2)+O(\varrho_T\hG).
\]
So $\sigma_Fd(c_T)\ge c_sc_0\hG-C\hG^2>C_1\hG^2$ for $\hG\le h_\ast$. With $y:=\pi(c_T)$, the normal segment from $c_T$ to $\Psi(y,\sigma_F\frac r2)$ misses $\Sigma$ and, since $|s|>C_1\hG^2$ along it, misses $\Gh$. It starts in $\operatorname{int}T\subset\Oh$, so $\Psi(y,\sigma_F\frac r2)\in\Oh$.

\emph{Conclusion.} If $\sigma_F=-1$ for some $F$, then $\Psi(y,-\frac r2)\in\Gamma_i^-\cap\Oh=\emptyset$, a contradiction.
\end{proof}

\begin{remark}\label{geom:rem}
(1) Only the tetrahedra touching $\Gh$ are used, through $\varrho_T\ge c\diam F$. (2) Without (OB) the same proof shows that on each $\Gamma_i$ either all $\sigma_F=+1$, or all $\sigma_F=-1$ (a mirror statement), or both occur, and then every generic crossing number is even; Example~\ref{ex:filled} is of the third kind. (3) The manifold condition in (G1) is not needed for Theorem~\ref{geom:orient}: an edge in three or more faces, or a vertex whose star is not a disc, would produce $k(y)\ge2$. (4) The structure of the proof (local homeomorphism plus properness gives a covering; one sheet from a single fibre) is that of Amenta, Choi, Dey and Leekha \cite[Theorem~15]{AmentaChoiDeyLeekha2000}, who assume that adjacent triangles meet at an angle greater than $\pi/2$. Here the fibre count comes from the in/out crossing rule of $\partial\Oh$, and the orientation from the tetrahedra touching $\Gh$; no angle condition between neighbouring faces is needed.
\end{remark}

Both hypotheses are necessary. In the next two examples $\Omega=R\setminus\overline D$ with $D$ the unit ball, $N(y)=y$ and $r=\frac12$.

\begin{example}[(G1) and (OB) do not suffice: an inscribed sliver]\label{ex:sliver}
Let $P_h$ be a convex polyhedron inscribed in the unit sphere satisfying (G1) with $\hG=h$, and let $F$ be a face with $\diam F=h$, so that its circumradius satisfies $r_F\ge h/2$. Consider the square $Q$ of side $h/3$ in the plane of $F$ whose centre lies at distance $h/20$ from $o_F$, and let $q_1,q_2,q_3$ be the points of the spherical cap over $F$ whose orthogonal projections onto the plane of $F$ are three vertices of $Q$. The plane through $q_1,q_2,q_3$ meets the sphere in a circle; choose $q_4$ on the cap, within $h^2$ of the point over the fourth vertex of $Q$, and not on that circle. (The square is placed off-centre, and $q_4$ is perturbed, because four points of the sphere over a square centred on the axis are concyclic, hence coplanar.) All $q_j$ project to within $\frac{\sqrt2}6h+\frac h{20}+h^2<\frac h2\le r_F$ of $o_F$, so they lie strictly beyond the plane of $F$. Let $P':=\conv\{q_j\}$, a genuine tetrahedron of thickness $O(h^2)$, and put $\Oh:=R\setminus(P_h\cup P')$, $\Gh:=\partial P_h\sqcup\partial P'$.
\begin{itemize}
\item $P'\cap P_h=\emptyset$, since $P_h$ lies in the closed half-space bounded by the plane of $F$ and $P'$ in the open complementary one.
\item The faces of $P'$ are, up to a relative perturbation of order $h$, right isosceles triangles with legs $h/3$; their minimal angles tend to $45^\circ$ and their diameters to $\frac{\sqrt2}3h$. So (G1) holds with $c_0\le0.47$, and $\Gh$ is embedded. Two faces of $P'$ face away from the centre and two face towards it.
\item (OB) holds: $\Gamma^-=\{|x|=\frac34\}\subset P_h$ for small $h$.
\item $\pi|_{\Gh}$ is not injective: a generic point of $\pi(P')$ has three preimages, one on $\partial P_h$ and one on each side of $P'$. The faces of $P'$ facing the centre have $\sigma=-1$; the thin lens between $F$ and $P'$ belongs to $\Oh$.
\end{itemize}
By Theorem~\ref{geom:mesh}, this $\Oh$ admits no shape-regular conforming tetrahedral mesh with the faces of $\Gh$ as boundary faces, once $h$ is small: the lens has thickness $O(h^2)$ and is bounded by faces of diameter at least $c_0h$. So a volume hypothesis such as (M0$^3$) is needed if (OR) or the lift itself is not assumed. We have not constructed an example with connected $\Gh$.
\end{example}

\begin{example}[(G1) and (M0$^3$) do not suffice: a filled obstacle]\label{ex:filled}
Let $P'$ be as in Example~\ref{ex:sliver}, $\Oh:=R\setminus P'$ and $\Gh:=\partial P'$. Then $\Gh$ is connected and homeomorphic to $\Gamma$, (G1) and (G3) hold, $\partial\Oh=\partial R\sqcup\Gh$ and $\Oh$ is connected; but $\pi|_{\Gh}$ covers a patch of diameter $O(h)$ twice and misses the rest of $\Gamma$, and (OB) fails. A shape-regular mesh exists: with $o$ the centre of $Q$ and $p_\pm:=o\pm hN(o)$, the bipyramid $\conv(P'\cup\{p_+,p_-\})\setminus\operatorname{int}P'$ is the union of four shape-regular tetrahedra $\conv(G,p_\pm)$, $G$ the faces of $P'$. Completing this to a shape-regular conforming mesh of the box minus a convex polyhedron of fixed shape at scale $h$ is routine; we do not write it out. So even ``$\Gh$ is a triangulation of $\Gamma$'' does not give the lift; a global condition tying $\Oh$ to $\Omega$, such as (OB), is needed.
\end{example}

\subsection{Consequences}
From now on (G1), (G3), (M0$^3$) and (OB) hold. By Theorem~\ref{geom:mesh}, $\sigma_F=+1$ for every face, $\pi|_{\Gh}$ is a homeomorphism and Corollary~\ref{geom:graph} applies. Since $\Oh$ is connected by (M0$^3$), $\Oh\setminus\overline{\mathcal T_{r/2}}=\Omega\setminus\overline{\mathcal T_{r/2}}$: a component $C$ of $\R^3\setminus(\Sigma\cup\overline{\mathcal T_{r/2}})$ misses $\partial\Oh\cup\partial\Omega$, so it lies in or outside each set; if $\overline C$ meets $\partial\mathcal T_{r/2}$, it lies in both or in neither by Corollary~\ref{geom:graph}, and if not, then $\partial C\subset\Sigma$ and $C$ would be a relatively closed and open proper subset of the connected sets $\Oh$ and $\Omega$, so it lies in neither.

\begin{lemma}[One boundary face per tetrahedron]\label{geom:one}
For $\hG\le h_\ast$, no tetrahedron of $\Th$ has two faces in $\Gh$. The same holds for the macro-tetrahedra of an Alfeld refinement.
\end{lemma}

\begin{proof}
Two such faces would share an edge $E$ of $\Gh$, and the dihedral angle of the tetrahedron at $E$ would equal the dihedral angle of $\Oh$ at $E$, which is $\pi+O(\hG)$ because the normals of the two faces agree with $n$ up to $O(\hG)$ (Theorem~\ref{geom:orient}). Shape regularity bounds the dihedral angles of tetrahedra by $\pi-\delta$.
\end{proof}

\begin{lemma}[Uniform Korn, trace and Friedrichs inequalities]\label{geom:korn}
There is a bi-Lipschitz homeomorphism $\Phi_h:\overline\Omega\to\overline\Oh$, the identity outside $\mathcal T_{r/4}$, mapping $\Gamma$ onto $\Gh$, with
\[
\norm{D\Phi_h-I}_{L^\infty}+\norm{D\Phi_h^{-1}-I}_{L^\infty}\le C\hG .
\]
Consequently there are $C_{\rm tr},C_F,\kappa$, independent of $h$, such that every $v\in H^1(\Oh)^3$ with $v|_\Sigma=0$ satisfies
\[
\norm v_{L^2(\Gh)}\le C_{\rm tr}\norm{\nabla v},\qquad\norm v\le C_F\norm{\nabla v},\qquad\norm{\nabla v}^2\le\kappa\,a_h(v,v).
\]
\end{lemma}

\begin{proof}
Let $\chi\in C^\infty(\R)$ with $\chi=1$ near $0$, $\chi=0$ on $[r/4,\infty)$, and put $\Phi_h(\Psi(y,s)):=\Psi(y,s+\chi(s)\eta_h(y))$ for $0\le s<r/2$, and $\Phi_h:=\mathrm{id}$ elsewhere on $\overline\Omega$. Since $|\chi'\eta_h|\le C\hG^2<1$, $s\mapsto s+\chi(s)\eta_h(y)$ is increasing and maps $[0,r/2)$ onto $[\eta_h(y),r/2)$; by Corollary~\ref{geom:graph} and the preceding paragraph, $\Phi_h$ maps $\overline\Omega$ onto $\overline\Oh$ and $\Gamma$ onto $\Gh$. In $(y,s)$ coordinates, $D\hat\Phi_h=I+\begin{psmallmatrix}0&0\\\chi\nabla_\Gamma\eta_h&\chi'\eta_h\end{psmallmatrix}$, and $D\Psi(y,s')D\Psi(y,s)^{-1}=I+O(|s-s'|)$; with Corollary~\ref{geom:graph} this gives the bound. For $w:=v\circ\Phi_h$, $D(w)=D(v)\circ\Phi_h+E$ with $|E|\le C\hG|\nabla v\circ\Phi_h|$. On the fixed Lipschitz domain $\Omega$, $w$ vanishes on $\Sigma$, a set of positive surface measure, so Korn's first inequality, the trace inequality and the Friedrichs inequality hold; changing variables back and absorbing $C\hG\norm{\nabla v}$ for $\hG\le h_\ast$ gives the claim.
\end{proof}

The error analysis compares $u_h$ with an extension of $u$ to $\Oh$, which overlaps the obstacle in a thin layer. In two dimensions a divergence-free extension comes from extending the stream function. In three dimensions we extend the flux $2$-form $\iota_u\mathrm{vol}$ instead, by a Hestenes-type reflection across $\Gamma$ in normal coordinates; closedness of the form, which is $\div u=0$, survives because pull-backs commute with exterior differentiation.

\begin{lemma}[Divergence-free extension]\label{geom:ext}
Let $u\in C^{k+1}(\overline\Omega)$ be divergence free and $p\in C^k(\overline\Omega)$, with $\Gamma\in C^{k+3}$. There are $\ut\in C^{k+1}(\mathcal O)$ and $\pt\in C^k(\mathcal O)$ on $\mathcal O:=\Omega\cup\mathcal T_{r/2}$, with $\ut=u$, $\pt=p$ on $\Omega$ and $\div\ut=0$ on $\mathcal O$. Moreover $\norm{\ut}_{L^\infty(\Gh)}\le C\hG^2$, and $|(F,v)|\le C\hG^2\norm{\nabla v}$ for $v\in\VR$.
\end{lemma}

\begin{proof}
Let $\omega:=\iota_u\mathrm{vol}$, a closed $C^{k+1}$ $2$-form on $\overline\Omega$ ($d\omega=(\div u)\mathrm{vol}$). On a collar $\Gamma_i\times[0,r)$ the pull-back $\beta:=\Psi^\ast\omega$ is closed and of class $C^{k+1}$ up to $s=0$, since $\Psi\in C^{k+2}$. In local coordinates $(\sigma_1,\sigma_2)$ on $\Gamma_i$, $\beta=a\,d\sigma_1\wedge d\sigma_2+b_1\,d\sigma_2\wedge ds+b_2\,ds\wedge d\sigma_1$. For $\lambda>0$ let $R_\lambda(y,s):=(y,-\lambda s)$; then $R_\lambda^\ast\beta$ has coefficients $a(\sigma,-\lambda s)$ and $-\lambda b_\alpha(\sigma,-\lambda s)$. Choose distinct $\lambda_j\in(0,1]$, $j=1,\dots,k+3$, and $c_j$ with $\sum_jc_j(-\lambda_j)^m=1$ for $m=0,\dots,k+2$ (a Vandermonde system), and set $\tilde\beta:=\sum_jc_jR_{\lambda_j}^\ast\beta$ for $s\in(-r/2,0)$ \cite{Hestenes1941}. The coefficients of $\tilde\beta$ and $\beta$ and all their derivatives of order at most $k+1$ agree at $s=0$ (for $a$ this uses $m\le k+1$, for $b_\alpha$ the shifted conditions $m+1\le k+2$), so the glued form is $C^{k+1}$ on $\Gamma_i\times(-r/2,r)$. It is closed on each side, because pull-backs commute with $d$, hence closed everywhere. Pushing forward by $\Psi^{-1}$ and taking the vector field $\ut$ with $\iota_{\ut}\mathrm{vol}$ equal to the result gives $\ut\in C^{k+1}$ with $\div\ut=0$. For $\pt$ use the scalar Hestenes reflection.

Since $\ut|_\Gamma=0$ and $|d|\le C\hG^2$ on $\Gh$, $\norm{\ut}_{L^\infty(\Gh)}\le C\hG^2$. $F$ vanishes on $\Omega$ and is supported in $S_h:=\overline{\Oh\setminus\Omega}=\{\eta_h\circ\pi\le d\le0\}$, a layer of width at most $C\hG^2$ in the normal direction. A one-dimensional Poincar\'e inequality along normal segments, with the Jacobian $J$ of (T), gives $\norm v_{L^2(S_h)}\le C\hG(\norm v_{\Gh}+\hG\norm{\nabla v})$, and $|(F,v)|\le\norm F_\infty|S_h|^{1/2}\norm v_{L^2(S_h)}$ with $|S_h|\le C\hG^2$; conclude with Lemma~\ref{geom:korn}.
\end{proof}
The extension needs neither a vector potential nor the zero-flux condition on $\Gamma$, and it works verbatim in two dimensions. Only $\ut\in C^2$ near $\Gamma$ is used in the geometric lemmas.

\begin{lemma}[Wall structure]\label{geom:wall}
On $\Gamma$: $\nabla u=\dn u\otimes n$, $\dn u\cdot n=0$ and $(\nabla u)^{\mathsf T}n=0$. In particular $W:=\partial_Nu|_\Gamma$, the wall shear, is a tangential field.
\end{lemma}

\begin{proof}
The tangential derivatives of $u$ vanish on $\Gamma$, and $0=\div u=n\cdot\dn u$.
\end{proof}

\section{Stability}\label{sec:stability}

\begin{lemma}[Inverse traces]\label{st:inverse}
There is $C_I$, depending only on $k$ and shape regularity, such that for $v\in\Vh$
\[
\sum_{F\in\FG}\diam F\,\norm{\dn v}^2_{L^2(F)}\le C_I^2\norm{\nabla v}^2,\qquad\text{hence}\qquad\norm{\dn v}_{\Gh}\le C_I(\rho/h)^{1/2}\norm{\nabla v},
\]
and for $q\in\Pih$, $\norm q^2_{\Gh}\le C_I^2\sum_{F}(\diam F)^{-1}\norm q^2_{T_F}\le C_I^2(\rho/h)\norm q^2$, where $T_F$ is the tetrahedron with face $F$.
\end{lemma}

\begin{proof}
Scaling on each $T_F$; the $T_F$ are distinct by Lemma~\ref{geom:one}, and $(\diam F)^{-1}\le\rho/h$.
\end{proof}

\begin{lemma}[Coercivity and continuity]\label{st:coercive}
Let $\mu_0:=4\kappa C_I^2\rho$. The condition $\gamma\ge\gamma_0:=4\kappa C_I^2/c_0$ implies $\mu\ge\mu_0$. If $\mu\ge\mu_0$, then $N_h(v,v)\ge c_1\tnorm v^2$ for all $v\in\VR$ and $|N_h(w,v)|\le M\tnorm w\tnorm v$, with $c_1=1/\max(2\kappa(1+C_I^2),2)$ and $M=3+(\kappa C_I^2)^{-1/2}$.
\end{lemma}

\begin{proof}
As in two dimensions \cite{PadhiLong2D}: $2|\ip{\dn v}v|\le\frac12a_h(v,v)+2\kappa\rho C_I^2h^{-1}\norm v^2_{\Gh}$ by Lemma~\ref{st:inverse}, Young's inequality and Lemma~\ref{geom:korn}; add $\sum_F\diam F\norm{\dn v}^2_F\le C_I^2\norm{\nabla v}^2$. For continuity, each $\dn$-term is at most $(\rho/\mu)^{1/2}\tnorm w\tnorm v$.
\end{proof}

\subsection{The continuous inf-sup constant}
Let $C_B(U)$ denote the Bogovski\u{\i} constant of a bounded connected Lipschitz domain $U$: the smallest $C$ such that every $q\in L^2_0(U)$ is the divergence of some $v\in H^1_0(U)^3$ with $\norm{\nabla v}\le C\norm q$. By the closed range theorem, $C_B(U)^{-1}$ is the continuous inf-sup constant of $U$.

\begin{lemma}[Uniform continuous inf-sup constant]\label{st:bog}
For $\hG\le h_\ast$, $C_B(\Oh)\le2C_B(\Omega)$.
\end{lemma}

This is an instance of the continuity of the inf-sup constant under Lipschitz-close perturbations, due to Bernardi, Costabel, Dauge and Girault \cite[Theorem~4.4]{BCDG2016}: if $F$ is a diffeomorphism with $\norm{DF-I}_\infty,\norm{DF^{-1}-I}_\infty\le\varepsilon$, the inf-sup constants of $U$ and $F(U)$ differ by at most $c(d)\varepsilon$. Our map $\Phi_h$ of Lemma~\ref{geom:korn} is only bi-Lipschitz and piecewise smooth, so we include the short argument, which uses nothing more.

\begin{proof}
Let $\Phi:=\Phi_h$, $\varepsilon:=\norm{D\Phi-I}_\infty+\norm{D\Phi^{-1}-I}_\infty\le C\hG$ and $J:=\det D\Phi$, so $|J-1|\le C\varepsilon$. Given $q\in L^2_0(\Oh)$, put $\hat q:=J\,(q\circ\Phi)$. Then $\int_\Omega\hat q=\int_{\Oh}q=0$ and $(1-C\varepsilon)\norm q\le\norm{\hat q}\le(1+C\varepsilon)\norm q$. Let $\hat w\in H^1_0(\Omega)^3$ with $\div\hat w=\hat q$ and $\norm{\nabla\hat w}\le C_B(\Omega)\norm{\hat q}$, and $w:=\hat w\circ\Phi^{-1}\in H^1_0(\Oh)^3$. By the chain rule $(\nabla w)\circ\Phi=\nabla\hat w\,(D\Phi)^{-1}$, so
\[
(\div w)\circ\Phi=\div\hat w+\operatorname{tr}\bigl(\nabla\hat w\,((D\Phi)^{-1}-I)\bigr),
\]
and, changing variables,
\begin{align*}
(q,\div w)_{\Oh}&=(\hat q,\div\hat w)_\Omega+\int_\Omega\hat q\,\operatorname{tr}\bigl(\nabla\hat w((D\Phi)^{-1}-I)\bigr)\\
&\ge\norm{\hat q}^2-\sqrt3\,\varepsilon\norm{\hat q}\norm{\nabla\hat w}\ge\norm{\hat q}^2\bigl(1-\sqrt3\,\varepsilon C_B(\Omega)\bigr).
\end{align*}
Also $\norm{\nabla w}\le(1+C\varepsilon)\norm{\nabla\hat w}\le(1+C\varepsilon)C_B(\Omega)\norm{\hat q}$. Hence $\sup_w(q,\div w)/\norm{\nabla w}\ge(1-C\varepsilon)\norm q/C_B(\Omega)$ for every $q\in L^2_0(\Oh)$, which is the inf-sup condition on $\Oh$ with constant $(1-C\varepsilon)/C_B(\Omega)$. For $\hG\le h_\ast$ this is at least $1/(2C_B(\Omega))$.
\end{proof}

\begin{remark}\label{st:rem:galdi}
A second route covers $\Oh$ by finitely many pieces that are star-shaped with respect to balls, uniformly in $h$, and applies the corresponding estimates for Bogovski\u{\i}'s operator \cite{Galdi,CostabelMcIntosh2010}. We carried this out for the model case and for general obstacles, with explicit cones and graph charts, but Lemma~\ref{st:bog} makes it unnecessary.
\end{remark}

\subsection{The discrete inf-sup constant on Alfeld refinements}
For a macro-tetrahedron $K$ with Alfeld split $K^r$ let $\mathring P^c_k(K^r)$ be the continuous piecewise-$P_k$ vector fields on $K^r$ vanishing on $\partial K$, and $\mathring P_{k-1}(K^r)$ the piecewise-$P_{k-1}$ functions on $K^r$ with mean zero on $K$. We use the following local result.

\begin{fact}[{Guzm\'an and Neilan \cite[Theorem~3.1]{GuzmanNeilan2018}}]\label{st:gn}
Let $k\ge1$ and $d=3$. For every macro-tetrahedron $K$ and every $p\in\mathring P_{k-1}(K^r)$ there is $v\in\mathring P^c_k(K^r)$ with $\div v=p$ on $K$ and $\norm v_{H^1(K)}\le C_{\rm loc}\norm p_{L^2(K)}$, where $C_{\rm loc}$ depends only on $k$ and the shape regularity of $K^r$.
\end{fact}

We read the wording of this theorem, and the statements of \cite[Proposition~4.1, Proposition~6.1, Corollary~6.2]{GuzmanNeilan2018}, through a text-extraction tool rather than directly from the typeset PDF; the wording should be rechecked against the published version before submission. %% VERIFY: GN18 Thm 3.1 wording matches arXiv:1710.08044 (version dated 7 May 2018; k>=1, C depends only on k and shape regularity of K^r), as do Prop 4.1, Prop 6.1, Cor 6.2 (k>=d); read again via a summarising fetch tool, SIAM typeset version not seen (notes/v7/cites/FINDINGS.txt [3]).
The global stability result of \cite{GuzmanNeilan2018} (Proposition~6.1 and Corollary~6.2, $k\ge d$) combines Fact~\ref{st:gn} with the stability of the Bernardi--Raugel pair \cite[Proposition~4.1]{GuzmanNeilan2018}, \cite{BernardiRaugel1985}, for which no dependence of the constant on the domain is stated there. For a family of domains $\Oh$ that dependence is exactly what matters, and the following lemma supplies it. Zhang \cite{Zhang2005} proves stability on Alfeld refinements for $d=3$, $k\ge3$; we do not use his proof.

\begin{lemma}[(IS3) on Alfeld refinements, uniformly in $h$]\label{st:IS3}
Let $\Th$ be the Alfeld refinement of a conforming macro mesh $\Th^M$ of $\Oh$ with shape-regularity constant $\sigma$, and $k\ge3$. For every $q\in\Pih\cap L^2_0(\Oh)$ there is $v\in\VO$ with $\div v=q$ and
\[
\norm{\nabla v}\le\bigl[C_{\rm loc}+C_\Pi\,C_B(\Oh)\,(1+\sqrt3\,C_{\rm loc})\bigr]\norm q,
\]
where $C_\Pi=C_\Pi(\sigma)$ is the constant of the Fortin operator in step~1. By Lemma~\ref{st:bog}, (IS3) holds with $\beta^{-1}\le C_{\rm loc}+2C_\Pi C_B(\Omega)(1+\sqrt3C_{\rm loc})$, independent of $h$.
\end{lemma}

\begin{proof}
\begin{enumerate}
\item \emph{A local Fortin operator for the Bernardi--Raugel pair.} Let $I_{SZ}$ be the Scott--Zhang $P_1$ interpolant on $\Th^M$ that preserves homogeneous boundary values \cite{ScottZhang1990}. For $u\in H^1_0(\Oh)^3$, $|I_{SZ}u|_{1,K}\le C|u|_{1,\omega_K}$ and $\norm{u-I_{SZ}u}_K\le Ch_K|u|_{1,\omega_K}$, with $C=C(\sigma)$ and $\omega_K$ the patch of $K$. For every interior face $F$ of $\Th^M$ let $b_F$ be its cubic face bubble; $b_Fn_F$ is piecewise $P_3$ on the Alfeld split and continuous, so it lies in $\VO$ because $k\ge3$. Set
\[
\Pi u:=I_{SZ}u+\sum_{F\ \text{interior}}\alpha_Fb_Fn_F,\qquad\alpha_F:=\frac{\int_F(u-I_{SZ}u)\cdot n_F}{\int_Fb_F}.
\]
Then $\int_F\Pi u\cdot n_F=\int_Fu\cdot n_F$ on interior faces, and both vanish on boundary faces. With $|\alpha_F|\le C|F|^{-1/2}\norm{u-I_{SZ}u}_{L^2(F)}$ and the scaled trace inequality, $|\Pi u|_{1,K}\le C(\sigma)|u|_{1,\omega'_K}$, hence $\norm{\nabla\Pi u}\le C_\Pi\norm{\nabla u}$ by finite overlap. No constant of the domain enters: every estimate is on a patch, and $u\in H^1_0$ is used only through the preserved zero trace.
\item \emph{Macro means.} Let $\bar q$ be the $L^2$ projection of $q$ onto macro-piecewise constants, so $\bar q\in L^2_0(\Oh)$ and $\norm{\bar q}\le\norm q$. Take $u\in H^1_0(\Oh)^3$ with $\div u=\bar q$ and $\norm{\nabla u}\le C_B(\Oh)\norm{\bar q}$, and $v_1:=\Pi u\in\VO$. By the divergence theorem and the flux property, $\int_K\div v_1=\int_Kq$ for every $K\in\Th^M$, and $\norm{\nabla v_1}\le C_\Pi C_B(\Oh)\norm q$.
\item \emph{Local fluctuations.} $r:=q-\div v_1$ is piecewise $P_{k-1}$ on each $K^r$ and has zero mean on every $K$. Fact~\ref{st:gn} gives $w_K\in\mathring P^c_k(K^r)$ with $\div w_K=r|_K$ and $\norm{\nabla w_K}_K\le C_{\rm loc}\norm r_K$; extended by zero, $w:=\sum_Kw_K\in\VO$.
\item \emph{Sum.} $v:=v_1+w$ satisfies $\div v=q$, and with $\norm r\le\norm q+\sqrt3\norm{\nabla v_1}$ the bound follows.\qedhere
\end{enumerate}
\end{proof}

\begin{remark}[Where the domain could enter]\label{st:rem:chenliu}
Only through $C_B(\Oh)$ in step~2. The local problem in step~3 is posed on one macro-element with zero data on its whole boundary, so boundary vertices of $\Oh$ play no role, and on Alfeld refinements there is no analogue of the two-dimensional singular-vertex condition. Subdivided wall faces do produce one (Proposition~\ref{st:wf}). The same combination (uniform continuous inf-sup via \cite{BCDG2016}, Bernardi--Raugel face fluxes, and \cite[Theorem~3.1]{GuzmanNeilan2018}) is used by Chen and Liu \cite[Lemma~3.9]{ChenLiu2026} for polyhedral approximations of curved domains.
\end{remark}

\begin{remark}[(IS3) on other meshes]\label{st:rem:status}
For the Scott--Vogelius pair with discontinuous $P_{k-1}$ pressures in three dimensions, the known stability results are for split meshes. On Alfeld refinements they hold for $k\ge3$ \cite{Zhang2005,GuzmanNeilan2018}; see \cite[Table~2.1]{FarrellMitchellScott2024}. On Freudenthal meshes stability holds for $k\ge6$ \cite{Zhang2011}; such meshes cannot be fitted to an inscribed surface. According to \cite{FarrellMitchellScott2024}, Neilan \cite[Proposition~6.5]{Neilan2015} extends Zhang's $k\ge6$ result to more general meshes satisfying a condition that he states as a conjecture. We know of no unconditional result on general unsplit tetrahedral meshes, and on such meshes (IS3) remains an assumption. %% VERIFY: Neilan 2015 read only through FMS24 (notes/v6/cites/FINDINGS.txt [5]).

Worsey--Farin splits are different. The stable Worsey--Farin Stokes pairs of Guzm\'an, Lischke and Neilan \cite{GuzmanLischkeNeilan2022} and of Fabien, Guzm\'an, Neilan and Zytoon \cite{FabienGuzmanNeilanZytoon2022} do \emph{not} use the full space of discontinuous piecewise-$P_{k-1}$ pressures: the pressure space carries conditions at the singular edges of the split. With boundary conditions, \cite{FabienGuzmanNeilanZytoon2022} impose $\theta_e(q)=q^{(1)}-q^{(2)}=0$ on the boundary singular edges (and an alternating-sum condition on interior ones). %% VERIFY: FGNZ condition read by the referee from arXiv:2105.09214; GLN22 body not re-read (arXiv rate limit), only its abstract.
So (IS3) \emph{as formulated} fails on Worsey--Farin splits (Proposition~\ref{st:wf}); this is known, and it is the three-dimensional analogue of the two-dimensional singular boundary vertex. None of our upper bounds, which all use (IS3) for the full space $\Pih$, applies to Worsey--Farin meshes as they stand; they would have to be redone with a constrained pressure space, together with a uniform inf-sup constant for that pair on the perturbed domains $\Oh$ (as Lemma~\ref{st:IS3} provides for Alfeld) and a lift lemma onto the constrained trace space of Proposition~\ref{st:wf}.
\end{remark}

\begin{proposition}[(IS3) fails on subdivided wall faces]\label{st:wf}
Let a wall face $F\in\FG$ of a macro-element be subdivided into coplanar sub-faces, and let $E$ be an edge of that subdivision interior to $F$. For $v\in\VO$, $\div v$ is continuous across the interior face that meets $F$ along $E$, at every point of $E$. Consequently $\div\VO\subsetneq\Pih\cap L^2_0$ for discontinuous $P_{k-1}$ pressures, and (IS3) is false. On one barycentric Worsey--Farin macro-element (interior point at the barycentre, face points at the face barycentres, $12$ sub-tetrahedra) with $v=0$ on the whole boundary, $\operatorname{codim}\div\VO=12k-4$ exactly for $k=2,3,4$ (ranks $27/47$, $87/119$, $195/239$), which is $4\times(3k-1)$: $k$ conditions on each of the three edges from the face point to the vertices of a face, one of them redundant at the face point. \status{computer-assisted, exact} for the count.
\end{proposition}
\begin{proof}
Let $s_1,s_2$ be the sub-tetrahedra sharing the interior face through $E$, both with a face in the plane of $F$, and $x\in E$. Then $\nabla v|_{s_j}(x)\,t=0$ for $t$ in the plane of $F$, since $v=0$ there, and $\nabla v|_{s_1}(x)\,d=\nabla v|_{s_2}(x)\,d$ for a direction $d$ in the common face transversal to $F$, by continuity across that face. These directions span $\R^3$, so $\nabla v|_{s_1}(x)=\nabla v|_{s_2}(x)$ and $\div v$ is continuous at $x$. A discontinuous piecewise-$P_{k-1}$ function with a jump along $E$ is therefore not in $\div\VO$.
\end{proof}
\cert{notes/v7/td/wf\_witness.py 2|3|4}{wf\_ranks.log}
The velocities of $\GS$ and $\CNS$ with the full space $\Pih$ are still pointwise divergence-free when they exist, but the pressure is not unique. Moreover the trace space of $\Zh$ on a subdivided face is a proper subspace of the continuous piecewise-$P_k$ zero-flux traces: along $E$ the jump of $\div_Fv_F-(d_F\cdot\nabla_F)(v\cdot n_F)/d_n$ must vanish, where $d=d_F+d_nn_F$ is the direction of the interior face. So Lemma~\ref{lk:lift} does not transfer either.

\begin{lemma}[Divergence range]\label{st:divrange}
Under (IS3), $\div\VR=\Pih$ and $\Wh\ne\emptyset$, and the velocities of $\GS$ and $\CNS$ are pointwise divergence-free.
\end{lemma}

\begin{proof}
$\div\VO=\Pih\cap L^2_0$ is (IS3). For $F\in\FG$, the cubic face bubble $b_F$ of the tetrahedron $T_F$ gives $b_Fn_F\in\VR$ ($k\ge3$), and $\int_{\Oh}\div(b_Fn_F)=\int_Fb_F\ne0$; this adds the constants. The rest is as in two dimensions \cite{PadhiLong2D}: for any $\mathcal G\in\Vh$ with trace $\gI$ on $\Sigma$ there is $v\in\VR$ with $\div v=\div\mathcal G$, so $\mathcal G-v\in\Wh$, and the continuity row forces $\div u_h=0$.
\end{proof}

\section{Error analysis}\label{sec:error}

Throughout this section (IS3) holds; on Alfeld refinements with $k\ge3$ it is Lemma~\ref{st:IS3}.

The plan follows the two-dimensional analysis of \cite{PadhiLong2D}. Inserting the extended exact solution into the discrete equations leaves two residuals: a \emph{geometric} one, $\Gc$, of size $\hG^{3/2}$ in the energy dual, caused by the fact that $\ut$ does not vanish on $\Gh$ and that $\Gh$ is not $\Gamma$; and, for $\GS$ only, a \emph{pressure} residual $\Rc(v)=-\ip{\pt-\pbar}{v\cdot n}$, the traction that the printed method omits. The geometric residual gives the rate $\hG^{3/2}$. The pressure residual is of order one, and the question is how the discrete solution responds to it. The answer is the leak: the response is, to leading order, $(h/\mu)$ times a fixed divergence-free field $v_\varphi$ with normal trace $-(p-\pbar)$ on the wall (Lemma~\ref{er:vphi}). A cancellation that uses $\div v=0$ on each wall face (Lemma~\ref{er:cancel}) shows that this field has bounded energy, which is why the $H^1$ rate of $\GS$ is $1$ and not $\frac12$.

\subsection{Consistency}

\begin{lemma}[Error equation]\label{er:eq}
For $v\in\VR$,
\[
N_h(\ut,v)-(\pt,\div v)+\ip{\pt}{v\cdot n}=(\ft,v)+\Gc(v),\qquad
\Gc(v):=\ip{(\nabla\ut)^{\mathsf T}n}{v}-\ip{\ut}{\dn v}+\frac\mu h\ip{\ut}{v}-(F,v).
\]
\end{lemma}

\begin{proof}
Integrate $a_h(\ut,v)=(D(\ut),\nabla v)$ by parts on $\Oh$ ($v=0$ on $\Sigma$), use $\div D(\ut)=\Delta\ut$, $-\Delta\ut=\ft-F-\nabla\pt$, $-(\nabla\pt,v)=(\pt,\div v)-\ip{\pt}{v\cdot n}$, and $D(\ut)n-\dn\ut=(\nabla\ut)^{\mathsf T}n$.
\end{proof}

\begin{proposition}[The transposed-gradient term]\label{er:transposed}
On $F\in\FG$, $(\nabla\ut)^{\mathsf T}n_F=n(x^\ast)\bigl(\dn u(x^\ast)\cdot(n_F-n(x^\ast))\bigr)+O(\hG^2)=O(\hG)$. Hence, for $\gamma\ge\gamma_0$ and $v\in\VR$,
\[
|\ip{(\nabla\ut)^{\mathsf T}n}{v}|\le C\hG\norm v_{L^1(\Gh)}\le C\hG(h/\mu)^{1/2}\tnorm v=C\gamma^{-1/2}\hG^{3/2}\tnorm v .
\]
\end{proposition}

\begin{proof}
By Lemma~\ref{geom:wall}, $(\nabla u)^{\mathsf T}n_F=n\,(\dn u\cdot n_F)=n\,(\dn u\cdot(n_F-n))$ on $\Gamma$. Taylor-expand $\nabla\ut$ from $x^\ast$ to $x$ (distance $O(\hG^2)$) and use $|n_F-n(x^\ast)|\le C\hG$ (Theorem~\ref{geom:orient}). Then use $|\Gh|\le C$ and $\norm v^2_{\Gh}\le(h/\mu)\tnorm v^2$.
\end{proof}

In the $\norm{\nabla v}$-dual norm the transposed term is generically only of order $\hG$ in three dimensions: on the sphere, testing with the interpolant of $\chi\psi n$ gives $\sum_F|F|\psi(x_F^\ast)\,\dn u(x_F^\ast)\cdot(x_c-o_F)+O(\hG^2)$ by Lemma~\ref{geom:sag}(c), and the two-dimensional odd-symmetry cancellation is not available. The energy-dual bound above is all that the error analysis uses.

\begin{lemma}[Geometric consistency]\label{er:G}
For $\gamma\ge\gamma_0$: $|\Gc(v)|\le C(1+\gamma^{1/2})\hG^{3/2}\tnorm v$ for $v\in\VR$, and $|\Gc(v)|\le C\hG^{3/2}\norm{\nabla v}$ for $v\in\VO$.
\end{lemma}

\begin{proof}
The transposed term is Proposition~\ref{er:transposed}. By Lemmas~\ref{geom:ext} and~\ref{st:inverse}, $|\ip{\ut}{\dn v}|\le C\hG^2C_I(c_0\hG)^{-1/2}\norm{\nabla v}$; $\frac\mu h|\ip\ut v|\le(\mu/h)^{1/2}C\hG^2\tnorm v=C\gamma^{1/2}\hG^{3/2}\tnorm v$; and $|(F,v)|\le C\hG^2\norm{\nabla v}$. On $\VO$ the first and third terms vanish.
\end{proof}

\begin{lemma}[Fluxes]\label{er:flux}
For $v\in\Zh$, $\sum_i\int_{\Gamma_{h,i}}v\cdot n=0$, where $\Gamma_{h,i}:=\pi^{-1}(\Gamma_i)\cap\Gh$. If $m\ge2$, the map $v\mapsto(\int_{\Gamma_{h,i}}v\cdot n)_i$ is onto $\{x\in\R^m:\sum_ix_i=0\}$ for small $h$.
\end{lemma}

\begin{proof}
The first statement is $\int_{\partial\Oh}v\cdot n=\int_{\Oh}\div v=0$ with $v=0$ on $\Sigma$. For the second, apply Lemma~\ref{er:vphi} below with $q=\alpha_i$ constant on $\Gamma_i$, $\sum\alpha_i|\Gamma_i|=0$: its field has flux $-|\Gamma_i|\alpha_i+O(\hG^2+h^k)|\alpha|$ through $\Gamma_{h,i}$, by Lemma~\ref{geom:face}(c).
\end{proof}

\begin{corollary}[$\GS$ on $\Zh$]\label{er:GSZ}
For $v\in\Zh$ and every constant $c$, $N_h(\ut-u_h,v)=\Rc(v)+\Gc(v)$ with $\Rc(v):=-\ip{\pt-c}{v\cdot n}$.
\end{corollary}

\begin{remark}[The constant is global]\label{er:rem:global}
By Lemma~\ref{er:flux}, $\int_{\Gamma_{h,i}}v\cdot n$ is free on $\Zh$ for each component separately when $m\ge2$; only the total vanishes. So the constant in $\Rc$ can only be a single global one, and the relevant size of the wall pressure is $G=\norm{p-\pbar}_{L^2(\Gamma)}$ with the global mean $\pbar$, not the sum of per-component variances. Correspondingly, the correction in $\CNS$ uses the global mean $\pbG$; per-component means would make the method inconsistent. The two-dimensional statements and proofs are in \cite{PadhiLong2D}; they are dimension-free.
\end{remark}

\subsection{Approximation and the leak field}

Let
\[
\epsd:=\bigl(1+(\gamma\hG)^{1/2}\bigr)h^k+h^{k+1}\bigl(\hG^{-1/2}+(\mu/h)^{1/2}\bigr).
\]

\begin{lemma}[Approximation in $\Wh$]\label{er:approx}
$E_h:=\inf_{z\in\Wh}\tnorm{\ut-z}\le C\epsd$. For bounded $\gamma$ the second term is at most $Ch^k$ if $\hG\ge h^2$. If the outer data have zero net discrete flux, $\int_\Sigma\gI\cdot\nu=0$, it reduces to $C(1+\gamma^{1/2})\hG^{k+1/2}$. Every $z\in\Wh$ satisfies $\tnorm{\ut-z}\ge(\mu/h)^{1/2}|\int_\Sigma\gI\cdot\nu|/|\Gh|^{1/2}$.
\end{lemma}

\begin{proof}
Let $w_h:=I_h\ut$ be the Lagrange interpolant. Boundary tetrahedra have diameter at most $C\hG$, so $\norm{\nabla(\ut-w_h)}\le Ch^k$, $\norm{\ut-w_h}_{L^\infty(\Gh)}\le C\hG^{k+1}$ and $\norm{\dn(\ut-w_h)}_{L^\infty(\Gh)}\le C\hG^k$, whence $\tnorm{\ut-w_h}\le C(1+(\gamma\hG)^{1/2})h^k$ and $\norm{\div w_h}\le Ch^k$.
\emph{Net flux.} Since $\int_{\Gh}\ut\cdot n=-\int_\Sigma g\cdot\nu=0$, $m:=\int_{\Oh}\div w_h=\int_\Sigma(\gI-g)\cdot\nu+\int_{\Gh}(w_h-\ut)\cdot n$, and $|m|\le Ch^{k+1}$ by pointwise interpolation estimates. (Lagrange interpolation on triangles has no extra quadrature exactness in general, unlike the one-dimensional Newton--Cotes case of \cite{PadhiLong2D}.)
\emph{Bubble.} $b:=\sum_{F\in\FG}b_Fn_F$ satisfies $\int_{\Gh}b\cdot n=c_b|\Gh|$ with $c_b>0$, $\norm{\nabla b}+\norm{\div b}\le C\hG^{-1/2}$ (there are $O(\hG^{-2})$ bubbles, each with $H^1$-seminorm squared $O(\hG)$), $\norm b_{\Gh}\le C$ and $\sum_F\diam F\norm{\dn b}^2_F\le C\hG^{-1}$. With $a:=m/(c_b|\Gh|)$, $\tnorm{ab}\le C|a|(\hG^{-1/2}+(\mu/h)^{1/2})$.
\emph{Correction.} $q:=\div(w_h-ab)\in\Pih\cap L^2_0$ with $\norm q\le C\epsd$; (IS3) gives $v\in\VO$ with $\div v=q$ and $\tnorm v\le(1+C_I^2)^{1/2}\beta^{-1}\norm q$. Then $z:=w_h-ab-v\in\Wh$. The lower bound holds because $\int_{\Gh}z\cdot n=-\int_\Sigma\gI\cdot\nu$ for $z\in\Wh$, while $\int_{\Gh}\ut\cdot n=0$.
\end{proof}

\begin{lemma}[A divergence-free field with normal trace $q$]\label{er:w}
Let $q\in C^{k+1}(\Gamma)$ with $\int_\Gamma q=0$. There is $w\in H^{k+1}(\Oh\cup\Omega)^3$, of class $C^{k+1}$ on $\mathcal T_{r/16}$, with $\div w=0$, $w=0$ near $\Sigma$, and $w|_\Gamma=qN=-qn$.
\end{lemma}

\begin{proof}
Let $\bar q_i$ be the mean of $q$ on $\Gamma_i$ and $\phi_i\in C^{k+2,\alpha}(\Gamma_i)$ a solution of $\Delta_\Gamma\phi_i=q-\bar q_i$ (Schauder theory on the closed $C^{k+3}$ surface). Let $\chi\in C^\infty(\R)$ with $\chi=1$ on $(-\infty,r/8]$ and $\chi=0$ on $[r/4,\infty)$. On $\Gamma\times(-r,r)$ with volume $dA\,ds$ put $V:=(-\chi'(s)\nabla_\Gamma\phi_i,\ \chi(s)q)$, whose product divergence is $-\chi'\Delta_\Gamma\phi_i+\chi'q=\chi'\bar q_i$. Let $w_0$ be its Piola push-forward by $\Psi$, i.e.\ $\iota_{w_0}\mathrm{vol}=(\Psi^{-1})^\ast\iota_V(dA\wedge ds)$. Since $\Psi^\ast\mathrm{vol}=J\,dA\wedge ds$, $\div w_0=J^{-1}\chi'(s)\bar q_i$; at $s=0$, $D\Psi$ is the identity on $T\Gamma$ plus $\partial_s\mapsto N$, $J=1$ and $\chi'=0$, so $w_0=qN$ on $\Gamma$. Regularity: $\nabla_\Gamma\phi_i\in C^{k+1,\alpha}$ and $D\Psi,J\in C^{k+1}$, so $w_0\in C^{k+1}$.
$g:=\div w_0\in C^k$ is supported in $\{r/8\le d\le r/4\}\subset\Omega$, and $\int g=\sum_i\bar q_i|\Gamma_i|(\chi(r/4)-\chi(0))=-\int_\Gamma q=0$. Let $\mathcal W:=\{x\in\Omega:d(x)>r/16\}$. It is open and connected: the map $\Psi(y,s)\mapsto\Psi(y,\max(s,r/8))$ on $\Omega\cap\mathcal T_r$, extended by the identity, is a continuous map of the connected set $\Omega$ onto $\{x\in\Omega:d(x)\ge r/8\}$, and $\mathcal W$ is the union of this set with the connected collars $\Psi(\Gamma_i\times(r/16,r/8])$, each of which meets it. Join the $m$ connected sets $\Psi(\Gamma_i\times[r/8,r/4])$, whose union contains $\operatorname{supp}g$, by paths in $\mathcal W$; the union $K$ is compact and connected. Take $\psi\in C^\infty_c(\mathcal W)$ with $0\le\psi\le1$ and $\psi=1$ near $K$, a regular value $t\in(0,1)$ of $\psi$ (Sard), and let $U$ be the component of $\{\psi>t\}$ containing $K$: a bounded connected domain with $C^\infty$ boundary and $\overline U\subset\mathcal W$, so $\overline U$ has positive distance from $\Sigma$ and $d>r/16$ on $\overline U$. A Bogovski\u{\i} field $w_1\in H^{k+1}_0(U)^3$ with $\div w_1=-g$ exists \cite{CostabelMcIntosh2010,Galdi}. Put $w:=w_0+w_1$ (extended by zero); $w=w_0$ on $\mathcal T_{r/16}$. If $m=1$ or $q$ has zero mean on each component, $w_1=0$. %% CHECKED (v7): Costabel--McIntosh (arXiv:0808.2614 = Math. Z. 265) state in the introduction solvability of div u=f in W^{m,p}_0(Omega)^n for f in W^{m-1,p}_0, int f=0, on all bounded Lipschitz domains; their Cor. 3.4 / Lemma 4.3 / Thm 4.6 give H^s -> H^{s+1} for all s (notes/v7/cites/FINDINGS.txt [4]). No Galdi locator is cited here.
\end{proof}
For the model case and $q=(p-\pbar)|_\Gamma$ this is the explicit field $w=\chi(r)r^{-2}q(\omega)e_r-\chi'(r)r^{-1}\nabla_{S^2}\phi_0(\omega)$ in spherical coordinates, with $\Delta_{S^2}\phi_0=q$.

From now on $q:=(p-\pbar)|_\Gamma$, which requires $p|_\Gamma\in C^{k+1}$, and $w$ is the field of Lemma~\ref{er:w}.

\begin{lemma}[Discrete leak field]\label{er:vphi}
Let $b$ and $c_b$ be as in the proof of Lemma~\ref{er:approx}, and
\[
v_\varphi:=I_hw-ab-v_0,\qquad a:=\frac{\int_{\Gh}I_hw\cdot n}{c_b|\Gh|},
\]
with $v_0\in\VO$ given by (IS3) with $\div v_0=\div(I_hw-ab)$. Then $v_\varphi\in\Zh$ and
\begin{gather*}
\norm{v_\varphi-w}_{L^\infty(\Gh)}\le C\hG^{k+1},\qquad\norm{\nabla(v_\varphi-w)}\le Ch^k,\\
\norm{\dn(v_\varphi-w)}_{\Gh}\le Ch^k\hG^{-1/2},\qquad\sum_F\diam F\norm{\dn(v_\varphi-w)}^2_F\le Ch^{2k}.
\end{gather*}
\end{lemma}

\begin{proof}
Since $\div w=0$ on $\Oh$ and $w=0$ near $\Sigma$, $\int_{\Gh}w\cdot n=0$. So $|a|\le C\norm{I_hw-w}_{L^\infty(\Gh)}\le C\hG^{k+1}$ (boundary tetrahedra have diameter $\le C\hG$ and $w\in C^{k+1}$ there), and $\div(I_hw-ab)\in\Pih\cap L^2_0$ with norm at most $\sqrt3\norm{\nabla(I_hw-w)}+\sqrt3|a|\norm{\nabla b}\le Ch^k$. Hence $\norm{\nabla v_0}\le Ch^k$; as $v_0=0$ on $\Gh$, Lemma~\ref{st:inverse} gives $\norm{\dn v_0}_{\Gh}\le Ch^k\hG^{-1/2}$ and $\sum_F\diam F\norm{\dn v_0}^2_F\le Ch^{2k}$. The terms $I_hw-w$ and $ab$ contribute $O(\hG^{k+1})$ pointwise on $\Gh$, $O(h^k)$ in $H^1$ and $O(\hG^k)$ pointwise for $\dn$.
\end{proof}

These bounds replace the $W^{1,\infty}$ bound for the curl of a $C^1$ stream-function interpolant used in two dimensions \cite{PadhiLong2D}, in every use below.

\begin{lemma}[Normal-load cancellation]\label{er:cancel}
Let $S\in W^{1,\infty}$ on a neighbourhood of $\Gh$ with $|S-(S\cdot n_F)n_F|\le C\hG$ on every $F\in\FG$. (For instance $S=w$, since $w=qN+O(d)$ and $|n_F-n|\le C\hG$.) Then for $v\in\Zh$
\[
|\ip{S}{\dn v}|\le C\bigl(\norm v_{\Gh}+\hG^{1/2}\norm{\nabla v}\bigr)\le C\norm{\nabla v},
\]
with $C$ depending on $\norm S_{W^{1,\infty}}$, $k$ and the mesh constants.
\end{lemma}

\begin{proof}
\emph{Tangential part.} $|\ip{S-(S\cdot n_F)n_F}{\dn v}|\le C\hG\norm{\dn v}_{\Gh}\le C\hG^{1/2}\norm{\nabla v}$ by Lemma~\ref{st:inverse}.
\emph{Normal part.} On $F$, $v$ is a polynomial on $T_F$ with $\div v=0$, and $n_F$ is constant; so $\dn v\cdot n_F=-\div_Fv_F$, where $v_F$ is the part of $v$ tangential to $F$ and $\div_F$ the in-plane divergence. With $\sigma_F:=S\cdot n_F$ and $\nu_{\partial F}$ the in-plane outward conormal,
\[
\int_F\sigma_F\,\dn v\cdot n_F=\int_F\nabla_F\sigma_F\cdot v_F-\int_{\partial F}\sigma_F\,v\cdot\nu_{\partial F},
\]
and the first term is at most $C\norm v_{L^1(F)}$.
\emph{Edge terms.} Each edge $E$ of $\Gh$ lies in exactly two faces $F,F'$, and $v$ is continuous across $E$, so the two contributions combine to $\int_Ev\cdot(\sigma_F\nu_{E,F}+\sigma_{F'}\nu_{E,F'})$. With $t_E$ a unit tangent, $\nu_{E,F}=t_E\times n_F$ and $\nu_{E,F'}=-t_E\times n_{F'}$, so $\nu_{E,F}+\nu_{E,F'}=t_E\times(n_F-n_{F'})=O(\hG)$; and $\sigma_F-\sigma_{F'}=S\cdot(n_F-n_{F'})=O(\hG)$. So the coefficient is $O(\hG)$.
\emph{Summing.} $\norm v_{L^1(\partial F)}\le C(\diam F)^{-1}\norm v_{L^1(F)}$ for $v|_F\in P_k$, so $\sum_E\hG\norm v_{L^1(E)}\le C\norm v_{L^1(\Gh)}\le C\norm v_{\Gh}$. Finish with Lemma~\ref{geom:korn}.
\end{proof}

If boundary faces of $\Th$ subdivide the faces of $\Gh$ (for example on Worsey--Farin splits), the extra edges lie in a plane, where $n_F=n_{F'}$ and the conormals are opposite, so their contributions cancel exactly. On such meshes, however, (IS3) as formulated fails (Proposition~\ref{st:wf}), so the results of this section do not apply to them as they stand.

\subsection{The printed method}
We now quantify the leak. The energy norm contains the penalty-weighted wall term $(\mu/h)\norm{\cdot}_{\Gh}^2$, so a normal slip of size $(h/\mu)G$ costs energy of order $(h/\mu)^{1/2}G$; the $H^1$ norm does not see the wall term directly, and the slip costs only order $h/\mu$ there. Theorems~\ref{er:A}--\ref{er:B} make both statements precise, with matching lower bounds, and Corollary~\ref{er:Bp} identifies the leading term.

Throughout this subsection $\gamma\ge\gamma_0$. By Theorem~\ref{geom:orient}, Lemma~\ref{geom:face}(c) and $\pt(x)=p(x^\ast)+O(\hG^2)$ on $\Gh$,
\begin{equation}\label{er:riemann}
\norm{\pt-\pbar}_{\Gh}\le G(1+C\hG^2)+C\hG^2,\qquad \norm{v_\varphi}_{\Gh}=G(1+O(\hG^2))+O(\hG^{k+1}),\qquad \tnorm{v_\varphi}^2=\frac\mu h\norm{v_\varphi}^2_{\Gh}+O(1).
\end{equation}

\begin{theorem}[$\GS$: energy norm, exact rate $1/2$]\label{er:A}
$\tnorm{\ut-u_h}\le C\bigl[(\hG/\gamma)^{1/2}G+(1+\gamma^{1/2})\hG^{3/2}+\epsd\bigr]$. If $G>0$ and $h\le h_0(G)$, with $h_0$ independent of $\gamma\ge\gamma_0$,
\[
\tnorm{\ut-u_h}\ge(2M)^{-1}(\hG/\gamma)^{1/2}G-C(1+\gamma^{1/2})\hG^{3/2}.
\]
Here $(\hG/\gamma)^{1/2}=(h/\mu)^{1/2}$.
\end{theorem}

\begin{proof}
\emph{Upper bound.} For $z\in\Wh$, $e_h:=z-u_h\in\Zh$, and by Corollary~\ref{er:GSZ} with $c=\pbar$ and Lemma~\ref{st:coercive}, $c_1\tnorm{e_h}^2\le N_h(z-\ut,e_h)+\Rc(e_h)+\Gc(e_h)$. By \eqref{er:riemann}, $|\Rc(e_h)|\le(G+C\hG^2)(h/\mu)^{1/2}\tnorm{e_h}$; apply Lemmas~\ref{er:G} and~\ref{er:approx}.
\emph{Lower bound.} On $F$, $v_\varphi\cdot n_F=w(x^\ast)\cdot n_F+O(\hG^2)=-q(x^\ast)\,n(x^\ast)\cdot n_F+O(\hG^2)$ and $n(x^\ast)\cdot n_F=1+O(\hG^2)$; with Lemma~\ref{geom:face}(c), $\Rc(v_\varphi)=G^2+O(\hG^2+h^k)$. By continuity $M\tnorm{\ut-u_h}\tnorm{v_\varphi}\ge|\Rc(v_\varphi)|-|\Gc(v_\varphi)|$, and $|\Gc(v_\varphi)|\le C(1+\gamma^{1/2})\hG^{3/2}\tnorm{v_\varphi}$. By \eqref{er:riemann}, $|\Rc(v_\varphi)|/\tnorm{v_\varphi}\ge\frac12(h/\mu)^{1/2}G$ for $h\le h_0(G)$.
\end{proof}

\begin{theorem}[$\GS$: $H^1$ rate $1$]\label{er:C}
With $C_p:=C(1+\norm{p-\pbar}_{W^{1,\infty}(\Gamma)})$,
\[
\norm{\nabla(\ut-u_h)}\le C_p\,\frac h\mu+C\bigl[(1+\gamma^{1/2})\hG^{3/2}+\epsd\bigr].
\]
\end{theorem}

\begin{proof}
\emph{Split.} For $z\in\Wh$ write $z-u_h=\delta_R+\delta_G$ with $\delta_R,\delta_G\in\Zh$ solving $N_h(\delta_R,v)=\Rc(v)$ and $N_h(\delta_G,v)=\Gc(v)-N_h(\ut-z,v)$ for all $v\in\Zh$. Then $\norm{\nabla\delta_G}\le\tnorm{\delta_G}\le c_1^{-1}(C(1+\gamma^{1/2})\hG^{3/2}+ME_h)$.
\emph{Ansatz.} Put $r_h:=\delta_R-(h/\mu)v_\varphi\in\Zh$. Then $N_h(r_h,v)=m(v)+\ell(v)$ with
\[
m(v):=-\ip{(\pt-\pbar)n+v_\varphi}{v},\qquad\ell(v):=-\frac h\mu\bigl[a_h(v_\varphi,v)-\ip{\dn v_\varphi}v-\ip{v_\varphi}{\dn v}\bigr].
\]
\emph{Bound on $m$.} On $F$, $(\pt-\pbar)n_F+v_\varphi=q(x^\ast)(n_F-n(x^\ast))+O(\hG^2)=O(\hG)$, with nonzero face means in general (Lemma~\ref{geom:sag}(c)). So we use only the energy-dual bound $|m(v)|\le C_p\hG\norm v_{L^1(\Gh)}\le C_p\hG(h/\mu)^{1/2}\tnorm v$.
\emph{Bound on $\ell$.} $|a_h(v_\varphi,v)|\le C\norm{\nabla v}$; $|\ip{\dn v_\varphi}v|\le(C+Ch^k\hG^{-1/2})\norm v_{\Gh}\le(C_p+Ch^k\hG^{-1/2})\norm{\nabla v}$; and by Lemma~\ref{er:cancel} with $S=w$ and Lemma~\ref{er:vphi}, $|\ip{v_\varphi}{\dn v}|\le|\ip w{\dn v}|+C\hG^{k+1}\norm{\dn v}_{L^1(\Gh)}\le C_p\norm{\nabla v}$. So $|\ell(v)|\le[C_p(h/\mu)+C(h/\mu)h^k\hG^{-1/2}]\norm{\nabla v}$, and $(h/\mu)h^k\hG^{-1/2}=\gamma^{-1}h^k\hG^{1/2}\le Ch^k$.
\emph{Close.} Testing with $r_h$, $c_1\tnorm{r_h}\le C_p(\hG(h/\mu)^{1/2}+h/\mu)+Ch^k$, and $\hG(h/\mu)^{1/2}\le\frac12(\hG^2+h/\mu)$ gives $\tnorm{r_h}\le C_p(h/\mu+\hG^2)+C\epsd$. Hence $\norm{\nabla\delta_R}\le(h/\mu)\norm{\nabla v_\varphi}+\tnorm{r_h}\le C_p(h/\mu+\hG^2)+C\epsd$, and $\hG^2\le\hG^{3/2}$.
\end{proof}

The rate is $1$ and not $1/2$ because $\div v=0$ turns the normal part of the missing traction into a tangential derivative (Lemma~\ref{er:cancel}), as in two dimensions.

\begin{theorem}[$\GS$: boundary slip and $H^1$ lower bound]\label{er:B}
Let $G>0$ and $X:=C_p(\hG/\gamma)+(1+\gamma^{1/2})\hG^{3/2}+\epsd$. Suppose that, for a fixed small $\delta$,
\[
X\le\delta\,G\min(G,\hG^{1/2}),\qquad\hG/\gamma\le\delta G,\qquad\hG\le h_0(G).
\]
All three hold for fixed $\gamma$ and bounded $\rho$ as $\hG\to0$. Then
\begin{enumerate}[label=(\roman*)]
\item $c(\hG/\gamma)G\le\norm{\ut-u_h}_{L^2(\Gh)}\le C(\hG/\gamma)^{1/2}\tnorm{\ut-u_h}$;
\item $\norm{\nabla(\ut-u_h)}\ge c(\hG/\gamma)G-Ch^{k+1/2}$.
\end{enumerate}
\end{theorem}

\begin{proof}
Let $e_u:=\ut-u_h$ and $v:=v_\varphi$. Corollary~\ref{er:GSZ} gives
\[
\frac\mu h\ip{e_u}v=\Rc(v)+\Gc(v)-a_h(e_u,v)+\ip{\dn e_u}v+\ip{e_u}{\dn v}.
\]
(1) $\Rc(v)=G^2+O(\hG^2+h^k)$. (2) $|\Gc(v)|\le C(\hG+\gamma\hG^3+\hG^{3/2})$: the transposed term is $O(\hG)$ pointwise (Proposition~\ref{er:transposed}); in the penalty term, $\ut=d\,W(x^\ast)+O(d^2)$ is tangential to $\Gamma$ while $v=-qn(x^\ast)+O(\hG^2)$ is normal, so $\ut\cdot v=O(\hG^4)$. By the hypotheses, $|\Gc(v)|\le\delta G^2$. (3) $|a_h(e_u,v)|\le C\norm{\nabla e_u}\le CX\le C\delta G^2$ by Theorem~\ref{er:C}. (4) By Theorem~\ref{er:A}, $|\ip{e_u}{\dn v}|\le C_p\norm{e_u}_{\Gh}\le C_p(h/\mu)^{1/2}\tnorm{e_u}\le C_p[(h/\mu)G+X]\le2C_p\delta G^2$. (5) With Lemma~\ref{st:inverse}, $|\ip{\dn e_u}v|\le\bigl(\norm{\dn(\ut-z)}_{\Gh}+C\hG^{-1/2}\norm{\nabla(z-u_h)}\bigr)\norm v_{\Gh}\le C\hG^{-1/2}XG\le C\delta G^2$. For $\delta$ small, $|\ip{e_u}v|\ge\frac12(h/\mu)G^2$; divide by $\norm v_{\Gh}\le G+C\hG^2$. For (ii) let $E\in H^1(\Oh)^3$ lift $(g-\gI)|_\Sigma$, vanishing near $\Gh$, with $\norm E_{H^1}\le Ch^{k+1/2}$; then $\norm{e_u}_{\Gh}=\norm{e_u-E}_{\Gh}\le C_{\rm tr}(\norm{\nabla e_u}+\norm{\nabla E})$.
\end{proof}

\begin{corollary}[Leak law]\label{er:Bp}
Let $G>0$ and $\hG\to0$ with $h/\mu\to0$, $\gamma\hG\to0$ and $\epsd/(\hG/\gamma)^{1/2}\to0$. Then
\[
\ut-u_h=\frac h\mu\,v_\varphi+\vartheta_h,\qquad\tnorm{\vartheta_h}=o\bigl((h/\mu)^{1/2}\bigr),\qquad\norm{\vartheta_h}_{\Gh}=o(h/\mu),
\]
and consequently
\[
\frac{\norm{\ut-u_h}_{L^2(\Gh)}}{(h/\mu)G}\to1,\qquad\frac{\tnorm{\ut-u_h}}{(h/\mu)^{1/2}G}\to1 .
\]
To leading order the printed method replaces no-slip by the leak condition $(\mu/h)\,u_h\cdot n\approx p-\pbar$, with the global mean $\pbar$.
\end{corollary}

\begin{proof}
Write $\ut-u_h=(\ut-z)+\delta_G+\delta_R$ and $\delta_R=(h/\mu)v_\varphi+r_h$ as above, so $\vartheta_h=(\ut-z)+\delta_G+r_h$. Relative to $(h/\mu)^{1/2}G$: $\tnorm{r_h}\le C_p(h/\mu+\hG^2)+C\epsd$ contributes $O((h/\mu)^{1/2}+(\gamma\hG^3)^{1/2}+\epsd(\gamma/\hG)^{1/2})\to0$, and $\tnorm{\delta_G}+\tnorm{\ut-z}\le C[(1+\gamma^{1/2})\hG^{3/2}+\epsd]$ contributes $O((\gamma^{1/2}+\gamma)\hG+\epsd(\gamma/\hG)^{1/2})\to0$. The $\Gh$ norms carry an extra factor $(h/\mu)^{1/2}$. The limits follow from \eqref{er:riemann}.
\end{proof}

For fixed $\gamma$ and bounded $\rho$, Theorems~\ref{er:C} and~\ref{er:B} give $\norm{\nabla(\ut-u_h)}\asymp h/\mu$. Section~\ref{sec:leak} identifies the constant: $\norm{\nabla(\ut-u_h)}\mu/(hG)\to\norm{\nabla U}/G$, with $U$ the Stokes flow driven by the leak as Dirichlet data (Corollary~\ref{lk:gen}), and the leak limit holds for every $\gamma\ge\gamma_0$ (Theorem~\ref{lk:up}).

\subsection{The corrected method}
The corrected method removes the pressure residual. The only subtlety is the constant pressure mode: the pressure part of the traction must be applied with the wall mean removed, otherwise constants in $\Pih$ are invisible to the momentum equation.

\begin{proposition}[The plain traction is singular]\label{er:singular}
If $\ip{p_h-\pbG(p_h)}{v\cdot n}$ is replaced by $\ip{p_h}{v\cdot n}$, the square system ($\VR\times\Pih$ tested by $\VR\times\Pih$) has the kernel vector $(0,1)$ and is singular.
\end{proposition}

\begin{proof}
For $v\in\VR$, $-(1,\div v)+\ip1{v\cdot n}=-\int_{\partial\Oh}v\cdot n+\int_{\Gh}v\cdot n=0$, since $v=0$ on $\Sigma$.
\end{proof}
The $L^2_0$ variant, and the relation to \cite{FrachonNilssonZahedi2024}, are discussed in \cite{PadhiLong2D}; the arguments there are dimension-free.

With the mean removed, the discrete pressure appears in the momentum equation through the wall term, which is not controlled by the inf-sup condition alone. The proof therefore bounds the pressure by (IS3) in terms of the velocity error, and absorbs the wall term for penalties above a threshold $\gamma_2$ that depends on the inf-sup constant.

\begin{theorem}[$\CNS$: well posed, with Scott's rate]\label{er:D}
There is $\gamma_2\asymp\max(\gamma_0,\beta^{-2})$ such that for $\gamma\ge\gamma_2$ the method $\CNS$ has a unique solution and
\[
\tnorm{\ut-u_h}\le C\bigl[(1+\gamma^{1/2})\hG^{3/2}+\epsd\bigr].
\]
Consequently, for $\gamma\in[\gamma_2,\gamma_{\max}]$ and $\hG\ge h^2$,
\[
\norm{\ut-u_h}_{H^1(\Oh)}\le C\bigl(\hG^{3/2}+h^k\bigr),
\]
and with outer data of zero discrete net flux the proviso $\hG\ge h^2$ is not needed. On Alfeld refinements this holds for every $k\ge3$.
\end{theorem}

\begin{proof}
\emph{Consistency.} Let $B^\ast(q,v):=-(q,\div v)+\ip{q-\pbG(q)}{v\cdot n}$. For a constant $c$, $B^\ast(q+c,v)=B^\ast(q,v)-c\int_{\Gh}v\cdot n$. With $p^\ast:=\pt-\pbG(\pt)$, Lemma~\ref{er:eq} gives $N_h(\ut,v)+B^\ast(p^\ast,v)=(\ft,v)+\Gc(v)$, so with $e_u:=\ut-u_h$, $e_p:=p^\ast-p_h$,
\begin{equation}\label{er:cnserr}
N_h(e_u,v)+B^\ast(e_p,v)=\Gc(v)\qquad\forall v\in\VR .
\end{equation}
\emph{Split.} $e_p=\eta+\xi$ with $\eta:=p^\ast-\pi_hp^\ast$ ($\pi_h$ the $L^2$ projection onto $\Pih$) and $\xi\in\Pih$; $(\eta,\div v)=0$ for $v\in\VR$, and $\norm\eta_{L^\infty(\Gh)}\le C\hG^k$.
\emph{Pressure.} For $v\in\VO$, $(\xi-\bar\xi,\div v)=a_h(e_u,v)-\ip{e_u}{\dn v}-\Gc(v)$, where $\bar\xi$ is the mean over $\Oh$. With $|\ip{e_u}{\dn v}|\le(\rho/\mu)^{1/2}C_I\tnorm{e_u}\norm{\nabla v}$ and Lemma~\ref{er:G}, (IS3) gives $\beta\norm{\xi-\bar\xi}\le C(\tnorm{e_u}+\hG^{3/2})$.
\emph{Velocity.} For $z\in\Wh$, $e_h:=z-u_h\in\Zh$. Testing \eqref{er:cnserr} with $e_h$, and using $(e_p,\div e_h)=0$ and $\int_{\Gh}e_h\cdot n=0$,
\[
c_1\tnorm{e_h}^2\le ME_h\tnorm{e_h}+|\Gc(e_h)|+|\ip\eta{e_h\cdot n}|+|\ip{\xi-\bar\xi}{e_h\cdot n}|,
\]
with $|\ip\eta{e_h\cdot n}|\le C\hG^k(h/\mu)^{1/2}\tnorm{e_h}$ and, by Lemma~\ref{st:inverse} and $\rho/\mu\le(c_0\gamma)^{-1}$, $|\ip{\xi-\bar\xi}{e_h\cdot n}|\le C_I(c_0\gamma)^{-1/2}\norm{\xi-\bar\xi}\tnorm{e_h}$.
\emph{Absorb.} Insert the pressure bound and $\tnorm{e_u}\le E_h+\tnorm{e_h}$; for $\gamma\ge\gamma_2$ with $C_I(c_0\gamma_2)^{-1/2}C\beta^{-1}\le c_1/2$ the pressure term is absorbed.
\emph{Uniqueness.} For zero data, $z=0$, so $u_h=0$; then $p_h$ is constant, and \eqref{er:cnserr} reads $-c\int_{\Gh}v\cdot n=0$ for all $v$, so $c=0$ by the face bubble of Lemma~\ref{st:divrange}. The system is square.
\emph{$H^1$ bound.} For bounded $\gamma$ the flux term of $\epsd$ is $Ch^{k+1}\hG^{-1/2}\le Ch^k$ iff $\hG\ge h^2$; with zero net outer flux it is $C(1+\gamma^{1/2})\hG^{k+1/2}$ (Lemma~\ref{er:approx}). The $L^2$ part uses the lift $E$ of Theorem~\ref{er:B} and Lemma~\ref{geom:korn}. On Alfeld refinements (IS3) is Lemma~\ref{st:IS3}.
\end{proof}

\begin{remark}[Flux-corrected outer data]\label{er:rem:flux}
In the model case one may replace $\gI$ by $\tilde g_I:=\gI-\frac{m_R}{3|R|}x$ with $m_R:=\int_{\partial R}\gI\cdot\nu$; since $\int_{\partial R}x\cdot\nu=3|R|$, this has zero net flux. In two dimensions the corresponding $H^1$ bound holds uniformly in $\gamma$ \cite{PadhiLong2D}. In three dimensions Corollary~\ref{lk:gen} gives a $\gamma$-uniform $H^1$ bound for $\GS$ with $\gI$ itself (the flux term appears as $h^{k+1}\hG^{-1/2}$); for $\CNS$ we have not carried the uniformity over.
\end{remark}

\section{Sharpness of \texorpdfstring{$\hG^{3/2}$}{hG^{3/2}}}\label{sec:sharp}

Theorem~\ref{er:D} gives $\norm{\nabla(\ut-u_h)}\le C[(1+\gamma^{1/2})\hG^{3/2}+\epsd]$ for $\CNS$, and Theorem~\ref{er:A} the same for $\GS$ when $G=0$. We prove the reverse inequality whenever $|\hh|\,\partial_Nu\not\equiv0$ on $\Gamma$. The lower bound does not use (IS3), except to guarantee that $u_h$ exists.

The idea is to test the error equation with discrete divergence-free fields $v$ that vanish on $\Gh$. For such $v$ neither the pressure nor the penalty enters, and the error equation says that $\ut-u_h$ must compensate the Nitsche term $\ip{\ut}{\dn v}$, which is nonzero because $\ut$ does not vanish on $\Gh$. On a face, $\ut\approx d\,\partial_Nu$ with $d$ the quadratic sagitta, so this term is a sum of three edge moments of $\dn v$ per face, weighted by the second fundamental form. The lower bound follows once we have local fields, one per face, that make these moments as large as scaling allows. This is the local condition (M3$^3$) below; the rest of the section verifies it.

\subsection{An abstract lower bound}
Let $Y_h:=\Zh\cap\VO$ and
\[
Y_h^1:=\Bigl\{v\in Y_h:\ \int_F\dn v\,dA=0\ \ \forall F\in\FG\Bigr\}.
\]
For $F\in\FG$ let $T_F$ be the tetrahedron of $\Th$ with face $F$; the $T_F$ are distinct (Lemma~\ref{geom:one}).

\begin{lemma}\label{sh:normal}
If $v\in Y_h$ and $F\in\FG$, then on $F$ (trace from $T_F$) $\nabla v=\dn v\otimes n_F$ and $\dn v\cdot n_F=0$.
\end{lemma}

\begin{proof}
$v|_{T_F}$ is a polynomial vanishing on the plane of $F$, so its tangential derivatives vanish there, and $0=\div v=\dn v\cdot n_F$.
\end{proof}

\begin{lemma}[Error identity on $Y_h$]\label{sh:identity}
Let $u_h$ be the velocity of $\CNS$ ($\gamma\ge\gamma_2$) or of $\GS$ ($\gamma\ge\gamma_0$, any $G$), and $e_u:=\ut-u_h$. For $v\in Y_h$, $a_h(e_u,v)-\ip{e_u}{\dn v}=-\ip{\ut}{\dn v}-(F,v)$.
\end{lemma}

\begin{proof}
By \eqref{er:cnserr}, resp.\ Corollary~\ref{er:GSZ}, $N_h(e_u,v)=\Gc(v)$ on $Y_h$: the pressure terms vanish because $\div v=0$ and $v=0$ on $\Gh$. Since $v=0$ on $\Gh$, $N_h(e_u,v)=a_h(e_u,v)-\ip{e_u}{\dn v}$, and the transposed and penalty parts of $\Gc(v)$ vanish.
\end{proof}

\begin{lemma}[Trace control on $Y_h^1$]\label{sh:poinc}
Let $C_{PT}$ be the constant in the scaled Poincar\'e trace inequality $\norm{w-\bar w_T}_{L^2(F)}\le C_{PT}(\diam T)^{1/2}\norm{\nabla w}_T$ for a face $F$ of a tetrahedron $T$, where $\bar w_T$ is the mean of $w$ over $T$. Then $|\ip w{\dn v}|\le C_{PT}C_I\norm{\nabla w}\norm{\nabla v}$ for $w\in H^1(\Oh)^3$ and $v\in Y_h^1$.
\end{lemma}

\begin{proof}
$\int_Fw\cdot\dn v=\int_F(w-\bar w_{T_F})\cdot\dn v$; use Cauchy--Schwarz on each $F$, the local form of Lemma~\ref{st:inverse}, and Cauchy--Schwarz over the distinct $T_F$.
\end{proof}

\begin{corollary}[Abstract lower bound]\label{sh:abstract}
With $\Lambda_0:=2+C_{PT}C_I$,
\[
\Lambda_0\norm{\nabla e_u}\ge\sup_{0\ne v\in Y_h^1}\frac{|\ip\ut{\dn v}+(F,v)|}{\norm{\nabla v}} .
\]
\end{corollary}

\begin{proof}
$|a_h(e_u,v)|\le2\norm{\nabla e_u}\norm{\nabla v}$, and Lemmas~\ref{sh:identity} and~\ref{sh:poinc}.
\end{proof}

Neither the pressure nor the penalty enters: the identity is the discrete momentum equation on $Y_h$, which $\ut$ violates through Nitsche's symmetric term acting on the nonzero trace of $\ut$ on $\Gh$.

\subsection{The leading term: three edge moments per face}
Let $W:=\partial_Nu|_\Gamma$, a tangential field (Lemma~\ref{geom:wall}), and $W_F:=W(y_F)$. For $F\in\FG$ put $w^F_{ij}:=\hh_F(e_{ij},e_{ij})$, $Q_F:=\frac12\sum_{i<j}w^F_{ij}\mu_i\mu_j$ (so $d=Q_F+O(\ell^3)$ on $F$ by Lemma~\ref{geom:face}(a)), and for $v\in Y_h$
\[
\M_F(v):=\int_FQ_F\,\dn v\,dA=\frac12\sum_{i<j}w^F_{ij}\int_F\mu_i\mu_j\,\dn v\,dA,\qquad \ell_W(v):=\sum_{F\in\FG}W_F\cdot\M_F(v).
\]
By Lemma~\ref{sh:normal}, $\M_F(v)$ is parallel to $F$. In the model case $w^F_{ij}=-\ell_{ij}^2$ (Lemma~\ref{geom:sag}).

\begin{lemma}[Leading term]\label{sh:lead}
For $v\in Y_h$, $|\ip\ut{\dn v}-\ell_W(v)|\le C\hG^{5/2}\norm{\nabla v}$ and $|(F,v)|\le C\hG^2\norm{\nabla v}$.
\end{lemma}

\begin{proof}
On $\Gh$, $x=x^\ast+d(x)N(x^\ast)$ and $\ut(x^\ast)=0$, so $\ut(x)=d(x)W(x^\ast)+O(d^2)$. With $|d-Q_F|\le C\hG^3$, $|W(x^\ast)-W_F|\le C\hG$ and $|Q_F|\le C\hG^2$, $\ut=Q_FW_F+R$ on $F$ with $|R|\le C\hG^3$. Hence $\int_F\ut\cdot\dn v=W_F\cdot\M_F(v)+O(\hG^3)\norm{\dn v}_{L^1(F)}$, and $\sum_F\norm{\dn v}_{L^1(F)}\le\sum_F|F|^{1/2}C_I(\diam F)^{-1/2}\norm{\nabla v}_{T_F}\le C\hG^{-1/2}\norm{\nabla v}$ since $\#\FG\le C\hG^{-2}$. The bound on $(F,v)$ is Lemma~\ref{geom:ext}.
\end{proof}

\subsection{The local condition and the lower bound}
\begin{assumption}[(M3$^3$)]\label{sh:M3}
There are $\kappa_0>0$ and $K_{\rm ov},C_s,C_1$ such that for every $F\in\FG$ and every unit $t\in T_{y_F}\Gamma$ there is $v_{F,t}\in Y_h^1$ with
\begin{enumerate}[label=(\roman*)]
\item $\norm{\nabla v_{F,t}}=1$;
\item $\operatorname{supp}v_{F,t}\subset\omega_F$, a union of tetrahedra within distance $C_s\hG$ of $F$, each tetrahedron lying in at most $K_{\rm ov}$ of the sets $\omega_F$;
\item $t\cdot\sum_{F'\subset\omega_F\cap\Gh}\M_{F'}(v_{F,t})\ \ge\ \kappa_0\,|\hh_{y_F}|\,(\diam F)^{5/2}-C_1(\diam F)^{7/2}$.
\end{enumerate}
\end{assumption}
Condition (iii) is scale invariant: under $x\mapsto\lambda x$ (with $v$ renormalised) $\norm{\nabla v}$ scales like $\lambda^{1/2}$, $w^F$ like $\lambda^2$ and $\M_F$ like $\lambda^3$. Since (iii) is linear in $v$, only the sign of $v$ is at our disposal; statements about $-\M_F$ are equivalent. The vector $t$ is tangent to $\Gamma$ at $y_F$, not to $F$; this is the form in which the witnesses are assembled.

Given the local witnesses, the lower bound is obtained by adding them up, face by face, with weights $|\hh|\,|W|$ and directions $W/|W|$; the overlap of their supports is bounded, so their energies add.

\begin{theorem}[$H^1$ lower bound]\label{sh:lb}
Assume the setting of Section~\ref{sec:setting}, (M3$^3$), and $|\hh|\,W\not\equiv0$. Let $u_h$ be the velocity of $\CNS$ ($\gamma\ge\gamma_2$) or of $\GS$ ($\gamma\ge\gamma_0$, any $G$). Then for $\hG\le h_1$
\[
\norm{\nabla(\ut-u_h)}_{\Oh}\ \ge\ c\,\hG^{3/2}\,\bigl\|\,|\hh|\,W\bigr\|_{L^2(\Gamma)}-C\hG^2,
\]
with $c=c(\kappa_0,K_{\rm ov},\Lambda_0,c_0,\theta_0)$ and $C$ independent of $\gamma,\mu,\rho,h$, the pressure and $\epsd$. Here $h_1$ depends also on $C_s$, $C_1$, $\norm W_{W^{1,\infty}}$ and $\norm{|\hh|W}_{L^2}$. In the model case $|\hh|=1$.
\end{theorem}

\begin{proof}
For each $F$ with $W_F\ne0$ let $t_F:=W_F/|W_F|\in T_{y_F}\Gamma$ and $v_F:=|\hh_{y_F}|\,|W_F|\,v_{F,t_F}$ (else $v_F:=0$), and $v:=\sum_Fv_F\in Y_h^1$. By linearity,
\[
\ell_W(v)=\sum_F\Bigl[W_F\cdot\sum_{F'}\M_{F'}(v_F)+\sum_{F'}(W_{F'}-W_F)\cdot\M_{F'}(v_F)\Bigr],
\]
the inner sums running over the (boundedly many) faces $F'\subset\omega_F$. Since $|\M_{F'}(v)|\le C\hG^2|F'|^{1/2}\norm{\dn v}_{L^2(F')}\le C\hG^{5/2}\norm{\nabla v}$ and $|W_{F'}-W_F|\le CC_s\hG\norm W_{W^{1,\infty}}$, (M3$^3$)(iii) gives
\[
\ell_W(v)\ge\kappa_0c_0^{5/2}\hG^{5/2}\sum_F|\hh_{y_F}|^2|W_F|^2-C\hG^{7/2}\sum_F|\hh_{y_F}||W_F|\bigl(|W_F|+\norm W_{W^{1,\infty}}\bigr),
\]
and $\norm{\nabla v}^2\le K_{\rm ov}\sum_F|\hh_{y_F}|^2|W_F|^2$. The projections $\pi(F)$ tile $\Gamma$ with area Jacobian $1+O(\hG^2)$ (Theorem~\ref{geom:orient}, Lemma~\ref{geom:face}(c)), and $|\hh|,|W|$ are Lipschitz, so $S:=\sum_F|F||\hh_{y_F}|^2|W_F|^2=\norm{|\hh|W}^2_{L^2(\Gamma)}+O(\hG)$; with $c_A\hG^2\le|F|\le\hG^2$ ($c_A$ from $c_0,\theta_0$) we get $\hG^{-2}S\le\sum_F|\hh_{y_F}|^2|W_F|^2\le(c_A\hG^2)^{-1}S$, and the error sum is at most $C\#\FG\le C\hG^{-2}$. Hence
\[
\frac{\ell_W(v)}{\norm{\nabla v}}\ge\kappa_0c_0^{5/2}\Bigl(\frac{c_A}{K_{\rm ov}}\Bigr)^{1/2}\hG^{3/2}S^{1/2}-C\frac{\hG^{5/2}}{S^{1/2}},
\]
and for $\hG\le h_1$ the second term is at most half the first and $S\ge\frac34\norm{|\hh|W}^2$. Conclude with Corollary~\ref{sh:abstract} and Lemma~\ref{sh:lead}.
\end{proof}

On flat parts of $\Gamma$ the faces lie in $\Gamma$ and the chordal bump vanishes, so the weight $|\hh|$ cannot be removed.

\subsection{Witnesses by affine transport}\label{sh:sec:transport}
On a single macro-element, witnesses can be computed once on a reference tetrahedron and transported to every shape by the Piola map, which preserves divergence-free fields and transforms the edge moments by an explicit linear map. So (M3$^3$) reduces to a finite-dimensional surjectivity statement on the reference element, which can be checked in exact arithmetic.

Let $T$ be a tetrahedron with a face $F$ and apex $p_0$, $\hat T=\conv(\hat p_0,\dots,\hat p_3)$ with $\hat p_0=e_3$, $\hat p_1=0$, $\hat p_2=e_1$, $\hat p_3=e_2$ and $\hat F=\hat T\cap\{\hat x_3=0\}$, and $x=A\hat x+a$ the affine map with $\hat F\mapsto F$, $\det A>0$. Let $\nu=-n_F$ be the inward unit normal of $T$ on $F$ and $h_T$ the height of $T$ over $F$. For $\hat v$ vanishing on $\hat F$ let $v:=(\det A)^{-1}A\,\hat v\circ A^{-1}$ (Piola) and $\gh:=\partial_{\hat\nu}\hat v|_{\hat F}$.

\begin{lemma}[Piola transport]\label{sh:piola}
(a) The Piola map preserves continuity, zero boundary values and piecewise-$P_k$ structure on affinely corresponding subdivisions, and $\div v=(\det A)^{-1}(\widehat\div\,\hat v)\circ A^{-1}$; it maps the Alfeld split of $\hat T$ onto that of $T$.
(b) For every polynomial weight $\phi$ in the barycentrics of $F$, $\int_F\phi\,\partial_\nu v\,dA=h_T^{-2}A\int_{\hat F}\phi\,\gh\,d\hat A$.
(c) $\norm{\nabla v}_T\le(\det A)^{-1/2}\norm A\norm{A^{-1}}\norm{\hat\nabla\hat v}_{\hat T}$.
\end{lemma}

\begin{proof}
(a) is standard ($A$ is constant and maps barycentre to barycentre). (b) On $\hat F$, $\hat\nabla\hat v=\gh\otimes e_3$, so $\partial_\nu v=(\det A)^{-1}A\gh\,(A^{-\mathsf T}e_3\cdot\nu)$ with $A^{-\mathsf T}e_3=|A^{-\mathsf T}e_3|\nu$. By Nanson's formula $dA=\det A\,|A^{-\mathsf T}e_3|\,d\hat A$, so the factor is $|A^{-\mathsf T}e_3|^2=|\nabla\lambda_0|^2=h_T^{-2}$. (c) $\nabla v=(\det A)^{-1}A(\hat\nabla\hat v)A^{-1}$.
\end{proof}
Since $\dn=-\partial_\nu$ on $F$, $\M_F(v)=-\frac12\sum w_{ij}\int_F\mu_i\mu_j\partial_\nu v$; the sign is absorbed by $v\mapsto-v$.

\begin{proposition}[Joint criterion]\label{sh:joint}
Let $\hat{\mathcal Y}$ be a space of continuous piecewise-polynomial divergence-free fields on a subdivision of $\hat T$ vanishing on $\partial\hat T$, $\hat{\mathcal Y}^1:=\{\hat v\in\hat{\mathcal Y}:\int_{\hat F}\gh=0\}$, and suppose the joint edge-moment map
\[
\hat E:\hat{\mathcal Y}^1\to(\R^2)^3,\qquad\hat v\mapsto\Bigl(\int_{\hat F}\mu_1\mu_2\gh,\ \int_{\hat F}\mu_1\mu_3\gh,\ \int_{\hat F}\mu_2\mu_3\gh\Bigr)
\]
is onto, with right inverse $\hat R$ (Euclidean norm on $\R^6$ to $\norm{\hat\nabla\cdot}_{L^2(\hat T)}$). Then for every tetrahedron $T$ with face $F$, every real weight vector $w\ne0$ and every unit $t\parallel F$ there is $v$ in the Piola image, supported in $T$, with $\norm{\nabla v}=1$ and
\[
t\cdot\frac12\sum_{i<j}w_{ij}\int_F\mu_i\mu_j\,\partial_\nu v\ \ge\ \frac{|w|\,(\det A)^{1/2}}{2h_T^2\norm A\norm{A^{-1}}^2\norm{\hat R}}\ \ge\ c(\text{shape})\,\frac{|w|\,(\diam F)^{1/2}}{\norm{\hat R}} .
\]
\end{proposition}

\begin{proof}
$s:=A^{-1}t$ is parallel to $\hat F$. Take $\hat v:=\hat R\bigl((w_{ij}s/|w|^2)_{i<j}\bigr)$, so that $\sum_{i<j}w_{ij}\hat E_{ij}(\hat v)=s$ and $\norm{\hat\nabla\hat v}\le\norm{\hat R}\norm{A^{-1}}/|w|$. By Lemma~\ref{sh:piola}(b), $\frac12\sum w_{ij}\int_F\mu_i\mu_j\partial_\nu v=\frac12h_T^{-2}As=\frac12h_T^{-2}t$; normalise with (c). Shape regularity gives $\norm A\lesssim\diam T$, $\norm{A^{-1}}\lesssim(\diam T)^{-1}$, $\det A\gtrsim(\diam T)^3$ and $h_T\lesssim\diam T\lesssim\diam F$.
\end{proof}
No compactness over shapes is used: $\hat{\mathcal Y}^1$ and $\hat E$ are fixed, and the shape enters only through $A$ and $w$.

\begin{theorem}[(M3$^3$) on Alfeld refinements, $k\ge4$]\label{sh:alfeld}
Let $\Th$ be an Alfeld refinement and $k\ge4$. Then (M3$^3$) holds with $\omega_F=T_F$ (the macro-tetrahedron containing $F$), $K_{\rm ov}=1$ and $\kappa_0=c(\text{shape regularity},\theta_0)/\norm{\hat R}$. \status{computer-assisted, exact}
\end{theorem}

\begin{proof}
Let $\hat{\mathcal Y}_k$ be the continuous piecewise-$P_k$ divergence-free fields on the Alfeld split of $\hat T$ vanishing on $\partial\hat T$. An exact rational computation in the Bernstein--B\'ezier basis ($35$ interior domain points, $105$ unknowns, $80$ divergence equations) gives, for $k=4$: the divergence has rank $79$ (onto the mean-free piecewise $P_3$), $\dim\hat{\mathcal Y}_4=26$, and the eight functionals ($\int_{\hat F}\gh$ and the three edge moments, two tangential components each) have rank $8$ on $\hat{\mathcal Y}_4$. So $\hat E$ is onto on $\hat{\mathcal Y}^1_4$ ($\dim=24$), and $\hat{\mathcal Y}_4\subset\hat{\mathcal Y}_k$. By Lemma~\ref{sh:piola}(a),(b) the Piola image on $T_F$, extended by zero, lies in $Y_h^1$: it is continuous and piecewise $P_k$ on the Alfeld refinement, divergence-free, zero on $\partial T_F\supset F$, has $\int_F\dn v=0$, and $T_F\cap\Gh=F$ (Lemma~\ref{geom:one}), so no other face constraint or moment is affected. For unit $t\in T_{y_F}\Gamma$ apply Proposition~\ref{sh:joint} to $t':=P_Ft/|P_Ft|$, $P_F$ the projection onto the plane of $F$; $|P_Ft|\ge\frac12$ because the angle between $T_{y_F}\Gamma$ and $F$ is $O(\hG)$, and $t\cdot\M_F=|P_Ft|\,t'\cdot\M_F$ because $\M_F\parallel F$. Finally $|w^F|\ge c\ell^2|\hh_{y_F}|-C\ell^3$ (Lemma~\ref{geom:face}(d)).
\end{proof}
The rank computation was repeated independently, and the norm $\norm{\hat R}=91.678155469$ of the right inverse is certified (Appendix~\ref{app:cert}); for a regular macro-tetrahedron $\kappa_0=2.11\cdot10^{-3}$ (Proposition~\ref{ct:kappa}).

\begin{proposition}[(M3$^3$) on any mesh, $k\ge9$]\label{sh:bubble}
Let $K=T_F$, $\lambda_0$ its barycentric coordinate vanishing on $F$, $B:=(\lambda_1\lambda_2\lambda_3)^2$, $\tau\parallel F$ a unit vector and $w\ne0$ real weights. With $Q_w:=\sum_{i<j}w_{ij}\mu_i\mu_j$ written in $\lambda_1,\lambda_2,\lambda_3$, $\bar Q_w:=\int_FBQ_w/\int_FB$ and $\psi:=\lambda_0^2B(Q_w-\bar Q_w)(\nu\times\tau)$, the field $v:=\curl\psi\in P_9^3$, extended by zero, lies in $Y_h^1$ and
\[
\frac12\sum_{i<j}w_{ij}\int_F\mu_i\mu_j\partial_\nu v=-|\nabla\lambda_0|^2\,\tau\int_FB(Q_w-\bar Q_w)^2,\qquad\frac1{|F|}\int_FB(Q_w-\bar Q_w)^2\ge\frac{|w|^2}{8316000}.
\]
Hence (M3$^3$) holds for $k\ge9$ on every shape-regular mesh, with $\omega_F=T_F$ and $K_{\rm ov}=1$.
\end{proposition}

\begin{proof}
$\psi$ vanishes to second order on $\partial K$, so $v$ vanishes on $\partial K$, is continuous after extension by zero, and is divergence free. With $f:=\lambda_0^2B(Q_w-\bar Q_w)$, $\curl(f\Phi)=\nabla f\times\Phi+f\curl\Phi$, and on $F$ only the term with $2\lambda_0\nabla\lambda_0$ survives one normal derivative: $\partial_\nu v|_F=2|\nabla\lambda_0|^2B(Q_w-\bar Q_w)\,\nu\times(\nu\times\tau)=-2|\nabla\lambda_0|^2B(Q_w-\bar Q_w)\tau$, which has mean zero. The identity is $\int B(Q_w-\bar Q_w)Q_w=\int B(Q_w-\bar Q_w)^2$. The lower bound is $w^{\mathsf T}\mathcal Vw$, where $\mathcal V=\frac1{1188000}I-\frac1{2772000}(J-I)$ is the exact $B$-weighted covariance matrix of the three edge bubbles on $F$ ($J$ the all-ones matrix), with eigenvalues $1/831600$ (twice) and $1/8316000$. By scaling, $\norm{\nabla v}\le C|w|(\diam K)^{-1/2}$ and the moment is at least $c|w|^2|F|/h_K^2$, so the ratio is at least $c|w|(\diam F)^{1/2}$. Transfer to $t\in T_{y_F}\Gamma$ as in Theorem~\ref{sh:alfeld}.
\end{proof}

\begin{proposition}[$k=3$: one macro-element is not enough]\label{sh:k3single}
On the Alfeld split with $k=3$, $\dim\hat{\mathcal Y}_3=6$, $\dim\hat{\mathcal Y}^1_3=4$, and the image of $\hat E$ on $\hat{\mathcal Y}^1_3$ is the line spanned by $(\hat p_1-\hat p_2,\ \hat p_3-\hat p_1,\ \hat p_2-\hat p_3)$. Consequently, for a field supported in one macro-tetrahedron $T_F$, the physical moment $\M_F(v)$ is a multiple of $\sum_{\rm cyc}w_{ij}e_{ij}$; on the sphere this is a multiple of $R_{\pi/2}(x_c-o_F)$ (Lemma~\ref{geom:sag}(d)), which vanishes on equilateral faces, and it spans one direction only. So (M3$^3$) cannot be verified with $\omega_F=T_F$ for $k=3$. \status{computer-assisted, exact}
\end{proposition}

\subsection{Degree three}
For $k=3$ a single macro-element cannot work (Proposition~\ref{sh:k3single}): its fields see only one combination of the three edge weights, in one direction. The remedy is to use several macro-elements at once. Take a wall vertex $z$ of the face $F$ and the macro edge $[z,q_F]$ from $z$ to the apex $q_F$ of $T_F$; the \emph{edge star} $R(z,q_F)$ is the union of the macro-tetrahedra that contain this edge. Its wall faces all meet at $z$, and discrete divergence-free fields on the star can transport a gradient at $z$ between them. Appendix~\ref{app:k3} describes the traces of such fields exactly, computes the edge moments of an explicit family of ``vertex moves'' in closed form, and shows that, together with two further families of moves, they see every tangential direction, including at saddle and parabolic points of $\Gamma$. The result is the following theorem.

\begin{theorem}[(M3$^3$) for $k=3$; Theorem~\ref{k3:witness}]
On Alfeld refinements there are $\kappa_0=\kappa_0(\theta,\theta_0)>0$, $C_1$ and $h_\ast$ such that for $\hG\le h_\ast$, (M3$^3$) holds with $\omega_F$ an edge star at a vertex of $F$ and $K_{\rm ov},C_s\le C(\theta)$, where $\theta$ is the minimal angle of the macro mesh.
\end{theorem}
The constant $\kappa_0$ comes from compactness; on the sphere an explicit criterion is available (Theorem~\ref{k3:sphere}).

\begin{theorem}[Lower bound for $k\ge3$ on Alfeld refinements]\label{sh:k3main}
Under the assumptions of Section~\ref{sec:setting}, on an Alfeld refinement with $k\ge3$, if $|\hh|\,\partial_Nu\not\equiv0$ and $\hG\le h_1$, then for $\CNS$ ($\gamma\ge\gamma_2$) and for $\GS$ ($\gamma\ge\gamma_0$)
\[
\norm{\nabla(\ut-u_h)}_{\Oh}\ge c\,\hG^{3/2}\bigl\|\,|\hh|\,\partial_Nu\bigr\|_{L^2(\Gamma)}-C\hG^2,
\]
with $c=c(\theta,\theta_0,\Lambda_0)$. It covers saddle points, parabolic points and flat regions.
\end{theorem}

\begin{proof}
Theorem~\ref{k3:witness} and Theorem~\ref{sh:lb}; for $k>3$ use $Y_3\subset Y_k$.
\end{proof}
At a parabolic point the choice of vertex matters: for $\hh=(x+y)^2$ a fixed vertex can give a degenerate star (an exact computation; Appendix~\ref{app:repro}). This affects semidefinite regions such as cylinders.

\begin{corollary}[Sharpness of Scott's rate]\label{sh:sharp}
On Alfeld refinements with $k\ge3$, or on any mesh with $k\ge9$ satisfying (IS3), if $|\hh|\,\partial_Nu\not\equiv0$, $\gamma\in[\gamma_2,\gamma_{\max}]$ and $\epsd\le C\hG^{3/2}$, then for $\CNS$
\[
c\,\hG^{3/2}\le\norm{\nabla(\ut-u_h)}\le\tnorm{\ut-u_h}\le C\hG^{3/2}.
\]
The same holds for $\GS$ when $G=0$.
\end{corollary}

\begin{proof}
Theorems~\ref{er:D}, \ref{er:A} and~\ref{sh:lb} with Theorems~\ref{sh:alfeld}, \ref{sh:k3main} or Proposition~\ref{sh:bubble}.
\end{proof}

\begin{remark}[Worsey--Farin splits]\label{sh:rem:wf}
On Worsey--Farin splits each face of $\Gh$ is subdivided into three coplanar sub-faces of $\Th$ (so (M0$^3$) is used with this modification), and $\Yh^1$ then requires $\int\dn v=0$ on each sub-face; the moments $\M_F$ are still taken over the whole face $F$ of $\Gh$. The lower bound of Theorem~\ref{sh:lb} uses only the discrete momentum equation on $\Yh$, so it holds for every solution of $\GS$ or $\CNS$ that exists; but existence and the upper bounds are not available there, because (IS3) as formulated fails (Proposition~\ref{st:wf}). Proposition~\ref{sh:bubble} still applies for $k\ge9$, and Proposition~\ref{sh:wf} below gives single-macro witnesses for $k\ge4$ on barycentric Worsey--Farin splits (interior point at the barycentre, face points at the face barycentres; this split is affine-invariant, so the Piola arguments of Section~\ref{sh:sec:transport} apply, unlike the incentre-based split of \cite{GuzmanLischkeNeilan2022}). For $k=3$ a single Worsey--Farin macro-element has no witness, and Worsey--Farin edge stars give rank $2$ in all our tests (ranks modulo a prime; Appendix~\ref{app:repro}) \status{numerical}; the trace analysis of Appendix~\ref{sh:sec:k3} has not been redone for split faces.
\end{remark}

\begin{proposition}[Sharpness witnesses on barycentric Worsey--Farin splits]\label{sh:wf}
Let $\hat{\mathcal Y}_k$ be the continuous piecewise-$P_k$ divergence-free fields on the barycentric Worsey--Farin split of $\hat T$ vanishing on $\partial\hat T$, and $\hat{\mathcal Y}^1_k$ the subspace with zero mean of $\gh$ on each of the three sub-faces of $\hat F$.
\begin{enumerate}[label=(\roman*)]
\item $k=3$: $\dim\hat{\mathcal Y}_3=18$, $\dim\hat{\mathcal Y}^1_3=12$, and both the quadratic edge-moment map $\hat E$ and the first-moment map vanish identically on $\hat{\mathcal Y}^1_3$. There is no single-macro witness of degree $3$.
\item $k=4$: $\dim\hat{\mathcal Y}^1_4=72$, $\hat E$ is onto with $\norm{\hat R_{\rm WF}}=54.001475715$ and the first-moment map is onto with right-inverse norm $9.709257103$ (both certified). Hence (M3$^3$) holds on barycentric Worsey--Farin refinements for every $k\ge4$, with $\kappa_0$ as in Proposition~\ref{ct:kappa} and $\norm{\hat R}$ replaced by $54.0015$, and Theorem~\ref{sh:lb} gives the lower bound $c\,\hG^{3/2}\norm{|\hh|\partial_Nu}_{L^2(\Gamma)}-C\hG^2$ for every solution of $\GS$ or $\CNS$ that exists.
\end{enumerate}
\status{computer-assisted, exact}
\end{proposition}
For $k\ge4$ this is a sharpness witness: together with an upper bound it would give the exact rate $\hG^{3/2}$ on barycentric Worsey--Farin splits. The upper bound needs the reformulated stability input of Remark~\ref{st:rem:status}, which we do not have.

\section{The leak limit and drag in three dimensions}\label{sec:leak}

This section identifies the $H^1$ constant of the printed method, uniformly in the penalty, and treats drag, lift and torque.

Corollary~\ref{er:Bp} says that the printed method replaces no-slip by the leak condition $(\mu/h)\,u_h\cdot n\approx p-\pbar$. It is natural to expect that, after scaling by $\mu/h$, the part of the error driven by the pressure converges to the continuous Stokes flow $U$ with Dirichlet data $-(p-\pbar)n$ on the wall, the \emph{leak flow}. We prove this for every penalty above the coercivity threshold, with no upper bound on $\gamma$ (Theorem~\ref{lk:up}), and we show that the rate $\hG^{1/2}$ is exact on Alfeld refinements (Theorem~\ref{lk:low}). The rate $\hG^{1/2}$ has a simple origin: the discrete wall data differ from $-(p-\pbar)n$ by a defect of amplitude $\hG$ that oscillates on the scale $\hG$, and such a defect has $H^{1/2}$-norm of order $\hG^{1/2}$ on a surface of any dimension. The dimension shows up in the constant: the defect contains a face-mean part, which is absent in two dimensions (Remark~\ref{lk:rem:facemean}). The same tools give the drag, lift and torque errors, by a reciprocity argument with a dual Stokes flow.

The standing assumptions are as follows. Throughout, (IS3) holds (on Alfeld refinements with $k\ge3$ it is Lemma~\ref{st:IS3}), $k\ge3$ and $\gamma\ge\gamma_0$. Put $q:=(p-\pbar)|_\Gamma$ and $\lambda:=\mu/h=\gamma/\hG$. For $F\in\FG$ let $x_c$ be its centroid, $x_F^\ast:=\pi(x_c)$, $q_F:=q(x_F^\ast)$, $P_F$ the orthogonal projection onto the plane of $F$, and $S_F:=P_FS_{x_F^\ast}P_F$, where $S$ is the shape operator ($D_Xn=SX$, $\ip{SX}Y=\hh(X,Y)$). $|\cdot|_{\rm F}$ is the Frobenius norm and
\[
G_\hh:=\bigl\|\,q\,|\hh|_{\rm F}\bigr\|_{L^2(\Gamma)}\qquad(\text{on the unit sphere }|\hh|_{\rm F}=\sqrt2,\ G_\hh=\sqrt2\,G).
\]

\begin{definition}[Leak flow]\label{lk:U}
$(U,P)\in H^1(\Omega)^3\times L^2(\Omega)$ is the Stokes flow with $-\Delta U+\nabla P=0$, $\div U=0$ in $\Omega$, $U=-qn$ on $\Gamma$ and $U=0$ on $\Sigma$.
\end{definition}
It exists because $\int_\Gamma U\cdot n=-\int_\Gamma q=0$ (the global mean; per-component fluxes need not vanish, Remark~\ref{er:rem:global}). Since $q\in C^k(\Gamma)$ with $k\ge3$, boundary Schauder theory for the Stokes system at the $C^{k+3}$ surface $\Gamma$ gives $(U,P)\in C^{2,\alpha}\times C^{1,\alpha}$ on a collar of $\Gamma$; no regularity at $\Sigma$ is used. Lemma~\ref{geom:ext}, applied on the collar, gives divergence-free extensions $\Ut\in C^2$, $\Pt\in C^1$ to $\Omega\cup\mathcal T_{r/2}$, and $F_U:=-\Delta\Ut+\nabla\Pt$ vanishes on $\Omega$, is bounded and supported in $\overline{\Oh\setminus\Omega}$, so $|(F_U,v)|\le C\hG^2\norm{\nabla v}$ for $v\in\VR$. The two-dimensional counterpart of every result below is in \cite{PadhiLong2D}.

Let $\Tr:=\{t\in C^0(\Gh)^3:\ t|_F\in P_k(F)^3\ \forall F\in\FG,\ \int_{\Gh}t\cdot n=0\}$ with $L^2(\Gh)$-orthogonal projection $\Pi_T$, let $\delta_R\in\Zh$ solve $N_h(\delta_R,v)=\Rc(v)$ for $v\in\Zh$ (as in the proof of Theorem~\ref{er:C}), and $X:=\lambda\delta_R$. Recall $\Yh=\Zh\cap\VO$ and $\Yh^1$ from Section~\ref{sec:sharp}.

\subsection{Tools}
\begin{lemma}[Lift of zero-flux traces]\label{lk:lift}
Every $t\in\Tr$ is the trace of some $L\in\Zh$ with $L|_\Sigma=0$ and $\norm{\nabla L}\le C\hG^{-1/2}\norm t_{\Gh}$, $C=C'(1+\sqrt3\beta^{-1})$. Consequently $\{v|_{\Gh}:v\in\Zh\}=\Tr$.
\end{lemma}
\begin{proof}
Let $L_0\in\Vh$ carry the nodal values of $t$ at the Lagrange nodes on $\Gh$ and vanish at all other nodes; its support, the union of the tetrahedra meeting $\Gh$, misses $\Sigma$ for $h\le h_\ast$. A node $z\in F$ has $|t(z)|\le C|F|^{-1/2}\norm t_{L^2(F)}$, and $\norm{\nabla\phi_z}^2_T\le C\diam T$; each node lies in boundedly many tetrahedra, so $\norm{\nabla L_0}^2\le C\sum_F(\diam F)^{-1}\norm t^2_F\le C(c_0\hG)^{-1}\norm t^2_{\Gh}$. Then $\div L_0\in\Pih$ and $\int_{\Oh}\div L_0=\int_{\Gh}t\cdot n=0$, and (IS3) gives $v\in\VO$ with $\div v=\div L_0$, $\norm{\nabla v}\le\sqrt3\beta^{-1}\norm{\nabla L_0}$. Put $L:=L_0-v$. Conversely, $v\in\Zh$ has a continuous piecewise-$P_k$ trace with $\int_{\Gh}v\cdot n=\int_{\Oh}\div v=0$.
\end{proof}

\begin{lemma}[Projection identity]\label{lk:proj}
Let $g:=-(\pt-\pbar)n_F$ on each $F\in\FG$ and $t^\ast:=\Pi_Tg$. Then $\Rc(v)=\ip gv=\ip{t^\ast}v$ and $N_h(X,v)=\lambda\ip{t^\ast}v$ for all $v\in\Zh$.
\end{lemma}
\begin{proof}
$\Rc(v)=-\ip{\pt-\pbar}{v\cdot n}=\ip gv$, and $v|_{\Gh}\in\Tr$ by Lemma~\ref{lk:lift}, so $\ip{g-\Pi_Tg}v=0$.
\end{proof}

\begin{lemma}[Normal data are invisible to $\dn$ on $\Zh$]\label{lk:canc}
Let $\sigma$ be Lipschitz on a neighbourhood of $\Gh$ and $g_\sigma:=\sigma n_F$ on each $F$. For $v\in\Zh$, $|\ip{g_\sigma}{\dn v}|\le C\norm\sigma_{W^{1,\infty}}\norm v_{\Gh}$; for $v\in\Yh$ it vanishes.
\end{lemma}
\begin{proof}
This is the normal part of the proof of Lemma~\ref{er:cancel}: $\dn v\cdot n_F=-\div_Fv_F$, integration by parts on $F$, and edge terms with coefficient $|t_E\times(n_F-n_{F'})|\le C\hG$. There is no tangential part, because $g_\sigma$ is exactly normal on every face, so no $\hG^{1/2}\norm{\nabla v}$ term appears. If $v\in\Yh$, $v=0$ on $\Gh$ and $\dn v\cdot n_F=\div v=0$ (Lemma~\ref{sh:normal}).
\end{proof}

\begin{lemma}[Tangential defect of the leak flow on a face]\label{lk:defect}
For $x\in F\in\FG$, with $\tau_F:=P_F\bigl(n(x^\ast_F)-n_F\bigr)$, $|\tau_F|\le C\hG$,
\[
P_F\Ut(x)=-q_F\bigl(\tau_F+S_F(x-x_c)\bigr)+O(\hG^2),\qquad\zeta:=\Ut-g=O(\hG)\ \text{on }\Gh,\qquad\norm\zeta_{\Gh}\le C\hG .
\]
\end{lemma}
\begin{proof}
$|x-x^\ast|\le C\hG^2$ (Lemma~\ref{geom:face}(a)), so $\Ut(x)=-q(x^\ast)n(x^\ast)+O(\hG^2)$. Expand $n(x^\ast)=n(x^\ast_F)+S_{x^\ast_F}(x^\ast-x^\ast_F)+O(\hG^2)$ with $x^\ast-x^\ast_F=P_F(x-x_c)+O(\hG^2)$, and use $P_Fn_F=0$ and $q(x^\ast)=q_F+O(\hG)$. For $\zeta$: $\zeta=-q_F(n(x^\ast)-n_F)+O(\hG^2)$ and $|n(x^\ast)-n_F|\le C\hG$ (Lemma~\ref{geom:face}(b)).
\end{proof}

\begin{remark}[The face-mean defect]\label{lk:rem:facemean}
The tangential defect on $F$ is affine, $-q_FS_F(x-x_c)$, \emph{plus the constant} $-q_F\tau_F=O(\hG)$. On the unit sphere $S=I$ on tangent planes and $\tau_F=-(x_c-o_F)+O(\hG^2)$, with $o_F$ the circumcentre of $F$ (Lemma~\ref{geom:sag}(b),(c)); the sign is fixed by $n$ pointing into the obstacle. So the defect is the radial field $-q(x-o_F)$ about the circumcentre, and its face mean is proportional to the centroid--circumcentre offset, nonzero on every non-equilateral face. In two dimensions the circumcentre of a chord is its midpoint, the face mean is $O(\hG^2)$, and the defect is an odd sawtooth. The constant does not change any rate below: the witnesses of Theorem~\ref{lk:low} have $\int_F\dn v=0$, and Lemma~\ref{lk:canc} does not see normal data. It does enter the constants (Corollary~\ref{lk:K}) and the size of $\Gc(w_\phi)$ in Lemma~\ref{dr:slip}. It is the same offset that makes the transposed term of Proposition~\ref{er:transposed} of order $\hG$ in the $\norm{\nabla v}$-dual norm.
\end{remark}

\subsection{The leak limit, uniformly in the penalty}
The proof compares the scaled pressure response $X=(\mu/h)\delta_R$ with a discrete Stokes--Dirichlet field $X^D$ whose wall trace is the projected leak datum. The penalty terms cancel exactly between the two, which is why no upper bound on $\gamma$ is needed; what remains is controlled by the cancellation lemma for normal data and the size of the tangential defect.

Let $\chi$ be a cutoff with $\chi=1$ on $\mathcal T_{r/4}$ and $\chi=0$ off $\mathcal T_{r/2}$, $I_h$ the $P_k$ Lagrange interpolant, $I_{SZ}$ the $P_k$ Scott--Zhang interpolant preserving the zero trace on $\Sigma$, and $I_h^\sharp\Ut:=I_h(\chi\Ut)+I_{SZ}((1-\chi)\Ut)$. For $h\le h_\ast$, $I_h^\sharp\Ut=I_h\Ut$ on the tetrahedra meeting $\Gh$ and $I_h^\sharp\Ut=0$ on $\Sigma$. Put
\[
E^0_U:=h+\norm{\nabla\bigl((1-\chi)\Ut-I_{SZ}((1-\chi)\Ut)\bigr)}_{\Oh}.
\]
$E^0_U\to0$ always (density and stability of $I_{SZ}$), and $E^0_U\le C(h+h^s\norm U_{H^{1+s}(\Omega)})$ if $U\in H^{1+s}(\Omega)$, $0<s\le1$. Only this term sees the regularity of $U$ at the edges and corners of $\Sigma$ (Remark~\ref{set:rem:corner}).

\begin{lemma}\label{lk:zU}
There is $z_U\in\Zh$ with $\norm{\nabla(\Ut-z_U)}\le CE^0_U$ and $\norm{\Ut-z_U}_{L^\infty(\Gh)}\le C\hG^2$.
\end{lemma}
\begin{proof}
$w:=I^\sharp_h\Ut$ vanishes on $\Sigma$, and $m:=\int_{\Gh}w\cdot n=\int_{\Gh}(w-\Ut)\cdot n$ satisfies $|m|\le C\hG^2$. With $b$ and $c_b$ from the proof of Lemma~\ref{er:approx}, $a:=m/(c_b|\Gh|)$ and $\norm{\nabla(ab)}\le C\hG^{3/2}$; (IS3) gives $v\in\VO$ with $\div v=\div(w-ab)$ and $\norm{\nabla v}\le\sqrt3\beta^{-1}(\norm{\nabla(w-\Ut)}+\norm{\nabla(ab)})$. Put $z_U:=w-ab-v$. On $\Gh$, $z_U=I_h\Ut-ab$.
\end{proof}

\begin{lemma}\label{lk:size}
$\norm{t^\ast-g}_{\Gh}\le C\hG$ and $\norm{t^\ast-\Ut}_{\Gh}\le C\hG$.
\end{lemma}
\begin{proof}
$z_U|_{\Gh}\in\Tr$, so $\norm{\Pi_Tg-g}\le\norm{z_U-g}\le\norm{z_U-\Ut}+\norm\zeta\le C\hG$; and $t^\ast-\Ut=(t^\ast-g)-\zeta$.
\end{proof}

Let $X^D\in\Zh$ be the discrete Stokes--Dirichlet field: $X^D|_{\Gh}=t^\ast$ and $a_h(X^D,v)=0$ for $v\in\Yh$. It exists and is unique: the affine set $\{z\in\Zh:z|_{\Gh}=t^\ast\}$ is nonempty by Lemma~\ref{lk:lift}, its direction is $\Yh$, and $a_h(y,y)=\norm{\nabla y}^2$ on $\Yh$.

\begin{lemma}\label{lk:XD}
$\norm{\nabla(X^D-\Ut)}\le C(\hG^{1/2}+E^0_U)$.
\end{lemma}
\begin{proof}
$W:=X^D-z_U$ has trace $t^\ast-z_U\in\Tr$ of norm $\le C\hG$; its lift $L$ has $\norm{\nabla L}\le C\hG^{1/2}$, and $y:=W-L\in\Yh$. Green's formula gives $a_h(\Ut,y)=(F_U,y)$, so $\norm{\nabla y}^2=a_h(\Ut-z_U,y)-(F_U,y)-a_h(L,y)$ and $\norm{\nabla y}\le2\norm{\nabla(\Ut-z_U)}+C\hG^2+2\norm{\nabla L}$.
\end{proof}

\begin{theorem}[Three-dimensional leak limit, uniform in the penalty]\label{lk:up}
For every $\gamma\ge\gamma_0$, with no upper bound on $\gamma$, and $D:=X-X^D\in\Zh$,
\[
\tnorm D\le C\Bigl[\Bigl(\frac\hG\gamma\Bigr)^{1/2}\bigl(1+\hG^{-1/2}E^0_U\bigr)+\hG^{1/2}\Bigr],\qquad\norm{\nabla(X-\Ut)}_{\Oh}\le C\bigl(\hG^{1/2}+E^0_U\bigr).
\]
In particular $\norm{\nabla\delta_R}\le C\,h/\mu$ uniformly in $\gamma\ge\gamma_0$, and $(\mu/h)\delta_R\to\Ut$ in $H^1$.
\end{theorem}
\begin{proof}
By Lemma~\ref{lk:proj} and $X^D|_{\Gh}=t^\ast$ the $\lambda$-terms cancel exactly: $N_h(D,v)=-a_h(X^D,v)+\ip{\dn X^D}v+\ip{t^\ast}{\dn v}=:\ell(v)$ on $\Zh$.
(i) Let $L_v\in\Zh$ lift $v|_{\Gh}$. Since $v-L_v\in\Yh$, $a_h(X^D,v)=a_h(\Ut,L_v)+a_h(X^D-\Ut,L_v)$. Green's formula ($\div L_v=0$, $L_v|_\Sigma=0$, $\int_{\Gh}L_v\cdot n=0$) gives $a_h(\Ut,L_v)=(F_U,L_v)+\ip{(D(\Ut)-(\Pt-c)I)n}v$, which is at most $C\norm v_{\Gh}$; and $|a_h(X^D-\Ut,L_v)|\le C(1+\hG^{-1/2}E^0_U)\norm v_{\Gh}$ by Lemmas~\ref{lk:XD} and~\ref{lk:lift}.
(ii) $\norm{\dn X^D}_{\Gh}\le\norm{\dn I^\sharp_h\Ut}_{\Gh}+C_I(c_0\hG)^{-1/2}\norm{\nabla(X^D-I^\sharp_h\Ut)}\le C(1+\hG^{-1/2}E^0_U)$ (Lemmas~\ref{st:inverse}, \ref{lk:XD}; $|\nabla I_h\Ut|\le C$ on the tetrahedra at $\Gh$).
(iii) $\ip{t^\ast}{\dn v}=\ip g{\dn v}+\ip{t^\ast-g}{\dn v}$: Lemma~\ref{lk:canc} with $\sigma=-(\pt-\pbar)$ bounds the first by $C\norm v_{\Gh}$, and Lemma~\ref{lk:size} with $\norm{\dn v}_{\Gh}\le(c_0\hG)^{-1/2}\tnorm v$ the second by $C\hG^{1/2}\tnorm v$.
With $\norm v_{\Gh}\le(\hG/\gamma)^{1/2}\tnorm v$ and coercivity (Lemma~\ref{st:coercive}; $c_1$ does not depend on $\mu\ge\mu_0$), take $v=D$. Then $\norm{\nabla(X-\Ut)}\le\tnorm D+\norm{\nabla(X^D-\Ut)}$ and $(\hG/\gamma)^{1/2}\le\gamma_0^{-1/2}\hG^{1/2}$.
\end{proof}

\subsection{A matching lower bound for every penalty}
The lower bound uses the same device as Section~\ref{sec:sharp}, with one change: the tangential defect of the leak flow on a face is affine in position, $-q_FS_F(x-x_c)$, so the relevant witnesses must see \emph{first} moments of $\dn v$ rather than the quadratic edge moments. On one Alfeld macro-element with $k\ge4$ they do, by an exact rank computation; for $k=3$ they do not.
\begin{lemma}[Identity on $\Yh$]\label{lk:id}
For every $\mu>0$ for which $\delta_R$ exists and every $v\in\Yh$, with $W:=X-\Ut$, $a_h(W,v)-\ip W{\dn v}=\ip{\Ut}{\dn v}-(F_U,v)$.
\end{lemma}
\begin{proof}
$N_h(\delta_R,v)=\Rc(v)=0$ and $N_h(\delta_R,v)=a_h(\delta_R,v)-\ip{\delta_R}{\dn v}$ for $v\in\Yh$; and $a_h(\Ut,v)=(F_U,v)$.
\end{proof}

\begin{lemma}[Single-macro leak witness]\label{lk:wit}
Let $\Th$ be an Alfeld refinement, $k\ge4$, $F\in\FG$, $T=T_F$ its macro-tetrahedron with reference map $x=A\hat x+a$ as in Section~\ref{sh:sec:transport}, $A_F:=A|_{\R^2\times\{0\}}$, $\ell:=\diam F$ and $h_T$ the height of $T$ over $F$. For every symmetric $S$ on the plane of $F$ there is $v\in\Yh^1$, supported in $T$, with $\norm{\nabla v}=1$ and
\begin{gather*}
S:\M^1_F(v)\ \ge\ \kappa_L(T)\,|S|_{\rm F}\,\ell^{3/2},\qquad \M^1_F(v):=\int_F\dn v\otimes(x-x_c)\,dA,\\
\kappa_L(T):=\frac{\sigma_{\min}(A_F)^2(\det A)^{1/2}}{h_T^2\norm A\norm{A^{-1}}\norm{\hat R_1}\,\ell^{3/2}},
\end{gather*}
where $\norm{\hat R_1}=3.645838299$ (certified; the interval has width below the last digit) is the norm of the right inverse of the first-moment map $\hat E^1:\hat v\mapsto\int_{\hat F}\gh\otimes(\hat x-\hat x_c)\in\R^{2\times2}$ on $\hat{\mathcal Y}^1_4$, from the Frobenius norm to $\norm{\hat\nabla\cdot}_{L^2(\hat T)}$. \status{computer-assisted, exact}
\end{lemma}
\begin{proof}
The exact computation (Bernstein--B\'ezier, over $\mathbb Q$) gives $\operatorname{rank}\hat E^1=4$ on $\hat{\mathcal Y}^1_4$ ($\dim24$) and brackets $\lambda_{\min}(\hat E^1\hat G^{-1}\hat E^{1\mathsf T})=7.523244811512\cdot10^{-2}$ by exact $LDL^{\mathsf T}$ tests. Let $\hat M:=A_F^{\mathsf T}SA_F/|A_F^{\mathsf T}SA_F|_{\rm F}$ and $\hat v:=\hat R_1\hat M$, so $\norm{\hat\nabla\hat v}\le\norm{\hat R_1}$. Its Piola image $v$ lies in $\Yh^1$ (proof of Theorem~\ref{sh:alfeld}). By Lemma~\ref{sh:piola}(b) with the affine weights $(x-x_c)_b=(A_F(\hat x-\hat x_c))_b$, $\int_F\partial_\nu v\otimes(x-x_c)=h_T^{-2}A_F\hat MA_F^{\mathsf T}$, and $\dn=-\partial_\nu$ (the sign is absorbed by $v\mapsto-v$). So $|S:\M^1_F(v)|=h_T^{-2}\tr(A_F^{\mathsf T}SA_F\hat M)=h_T^{-2}|A_F^{\mathsf T}SA_F|_{\rm F}\ge h_T^{-2}\sigma_{\min}(A_F)^2|S|_{\rm F}$. Normalise with Lemma~\ref{sh:piola}(c).
\end{proof}
At $k=3$ the same map has rank $1$ and its image is spanned by the skew matrix $\bigl(\begin{smallmatrix}0&1\\-1&0\end{smallmatrix}\bigr)$ (exact), so $S:\M^1_F(v)=0$ for every symmetric $S$: there is no single-macro leak witness of degree $3$, as there is none in two dimensions \cite{PadhiLong2D}.

\begin{theorem}[Lower bound with curvature weight]\label{lk:low}
Let $\Th$ be an Alfeld refinement, $k\ge4$, $\kappa_L:=\min_F\kappa_L(T_F)$ and $\Lambda_0=2+C_{PT}C_I$ (Corollary~\ref{sh:abstract}). For every $\mu>0$ for which $\delta_R$ exists (in particular every $\gamma\ge\gamma_0$),
\[
\norm{\nabla(X-\Ut)}_{\Oh}\ \ge\ c_L\,G_\hh\,\hG^{1/2}-C\hG^{3/2},\qquad c_L:=\Bigl(\frac{4}{\sqrt3}\Bigr)^{1/2}\frac{\kappa_L\,c_0^{3/2}}{\Lambda_0}.
\]
$c_L$ is explicit (Appendix~\ref{sec:constants}): $c_L=5.20\cdot10^{-3}$ when all macro-tetrahedra at $\Gh$ are regular and $c_0=1$. \status{computer-assisted, exact}
\end{theorem}
\begin{proof}
Take $v_F$ from Lemma~\ref{lk:wit} with $S=S_F$, signed so that $-S_F:\M^1_F(v_F)\ge\kappa_L|S_F|_{\rm F}\ell_F^{3/2}$, and $v:=\sum_Fq_F|S_F|_{\rm F}v_F\in\Yh^1$. The supports are distinct macro-tetrahedra (Lemma~\ref{geom:one}), so $\norm{\nabla v}^2=\Sigma:=\sum_Fq_F^2|S_F|^2_{\rm F}$, and on $F$ only $v_F$ is seen. By Lemma~\ref{lk:defect}, $\dn v\perp n_F$ (Lemma~\ref{sh:normal}) and $\int_F\dn v_F=0$, the constant $\tau_F$ drops out and
\begin{align*}
\ip{\Ut}{\dn v}&=\sum_Fq_F|S_F|_{\rm F}\bigl(-q_FS_F:\M^1_F(v_F)\bigr)+\sum_F|q_F||S_F|_{\rm F}\,O(\hG^2)\norm{\dn v_F}_{L^1(F)}\\
&\ge\kappa_Lc_0^{3/2}\hG^{3/2}\Sigma-C\hG^{3/2}\Sigma^{1/2},
\end{align*}
using $\norm{\dn v_F}_{L^1(F)}\le C\ell_F^{1/2}$ and $\#\FG\le C\hG^{-2}$. By Lemma~\ref{lk:id}, Lemma~\ref{sh:poinc} and $|a_h(W,v)|\le2\norm{\nabla W}\norm{\nabla v}$, $\Lambda_0\norm{\nabla W}\Sigma^{1/2}\ge\ip\Ut{\dn v}-C\hG^2\Sigma^{1/2}$. Finally $|F|\le\frac{\sqrt3}4\hG^2$ gives $\Sigma\ge\frac4{\sqrt3}\hG^{-2}\sum_F|F|q_F^2|S_F|^2_{\rm F}=\frac4{\sqrt3}\hG^{-2}(G_\hh^2+O(\hG))$, since the $\pi(F)$ tile $\Gamma$ with Jacobian $1+O(\hG^2)$. If $G_\hh^2\le C\hG$ the claim is trivial after enlarging $C$.
\end{proof}

\begin{corollary}[The rate of the leak limit is exactly $\frac12$]\label{lk:two}
On Alfeld refinements with $k\ge4$, if $G_\hh>0$, then for all $\gamma\in[\gamma_0,\infty)$ and $\hG\le h_1$
\[
\tfrac12c_LG_\hh\hG^{1/2}\le\Bigl\|\nabla\Bigl(\frac\mu h\delta_R-\Ut\Bigr)\Bigr\|\le C(\hG^{1/2}+E^0_U).
\]
The two sides have the same order whenever $E^0_U=O(\hG^{1/2})$, for example if $U\in H^{3/2}(\Omega)$ and $h\le C\hG$.
\end{corollary}
The weight is the second fundamental form: on flat parts of $\Gamma$ there is no $\hG^{1/2}$ term. In two dimensions the weight is the curvature, and the rate is again $\frac12$ \cite{PadhiLong2D}.

\subsection{The geometric part, uniformly in \texorpdfstring{$\gamma$}{gamma}, and the \texorpdfstring{$H^1$}{H1} constant}
Put $\varepsilon_1:=h^k+h^{k+1}\hG^{-1/2}$.
\begin{lemma}\label{lk:lin}
Let $z\in\Wh$ be the field of Lemma~\ref{er:approx}, $\delta_G\in\Zh$ with $N_h(\delta_G,v)=\Gc(v)-N_h(\ut-z,v)$ on $\Zh$, and $u^L_h:=z-\delta_G\in\Wh$. For every $\gamma\ge\gamma_0$, $\norm{\nabla(\ut-u^L_h)}\le C(\hG^{3/2}+\varepsilon_1)$.
\end{lemma}
\begin{proof}
(1) The construction of Lemma~\ref{er:approx} gives $\norm{\nabla(\ut-z)}\le C\varepsilon_1$ (the $\gamma$-dependence of $\epsd$ comes only from the penalty weight) and $\norm{\ut-z}_{\Gh}\le C(\hG^{k+1}+h^{k+1})$; also $\tnorm{\ut-u^L_h}\le C[(1+\gamma^{1/2})\hG^{3/2}+\epsd]$ (proof of Theorem~\ref{er:C}).
(2) For $v\in\Yh$, $\Gc(v)=-\ip\ut{\dn v}-(F,v)$, so $a_h(\ut-u^L_h,v)=-\ip{u^L_h}{\dn v}-(F,v)$.
(3) Let $m_\Sigma:=\int_\Sigma\gI\cdot\nu$, $a_\Sigma:=-m_\Sigma/(c_b|\Gh|)$ ($|a_\Sigma|\le Ch^{k+1}$) and $\Wh^D:=\{w\in\Wh:w|_{\Gh}=a_\Sigma b\}$. The trace $(z-a_\Sigma b)|_{\Gh}$ has zero flux and norm $\le C(\hG^2+h^{k+1})$; subtracting its lift (Lemma~\ref{lk:lift}) gives $z_D\in\Wh^D$ with $\norm{\nabla(\ut-z_D)}\le C(\varepsilon_1+\hG^{3/2})$. Let $u^D\in\Wh^D$ solve $a_h(\ut-u^D,v)=-(F,v)$ on $\Yh$; comparing with $z_D$, $\norm{\nabla(\ut-u^D)}\le C(\varepsilon_1+\hG^{3/2})$.
(4) $w:=u^L_h-u^D\in\Zh$ satisfies $a_h(w,v)=\ip{u^L_h}{\dn v}$ on $\Yh$ and has trace $u^L_h-a_\Sigma b$; with its lift $L_w$ and $y:=w-L_w\in\Yh$, $\norm{\nabla y}\le C\hG^{-1/2}\norm{u^L_h}_{\Gh}+2\norm{\nabla L_w}$, so $\norm{\nabla w}\le C\hG^{-1/2}(\norm{u^L_h}_{\Gh}+h^{k+1})$.
(5) $\norm{u^L_h}_{\Gh}\le\norm\ut_{\Gh}+(\hG/\gamma)^{1/2}\tnorm{\ut-u^L_h}\le C(\hG^2+\hG^{1/2}\gamma^{-1/2}\epsd)$, and $\gamma^{-1/2}\epsd\le C\varepsilon_1$.
\end{proof}

\begin{corollary}[$\GS$ in $H^1$, uniformly in $\gamma$]\label{lk:gen}
For every $\gamma\ge\gamma_0$,
\[
\Bigl\|\nabla\Bigl(\ut-u_h-\frac h\mu\Ut\Bigr)\Bigr\|_{\Oh}\le C\frac h\mu\bigl(\hG^{1/2}+E^0_U\bigr)+C\bigl(\hG^{3/2}+\varepsilon_1\bigr),\qquad
\norm{\nabla(\ut-u_h)}\le C\Bigl(\frac h\mu+\hG^{3/2}+\varepsilon_1\Bigr).
\]
Hence, if $G>0$, $\norm{\nabla(\ut-u_h)}\,\mu/(hG)\to K:=\norm{\nabla U}_{L^2(\Omega)}/G$ whenever $h/\mu\to0$ and $\frac\mu h(\hG^{3/2}+\varepsilon_1)\to0$; for the geometric part the condition is $\gamma\hG^{1/2}\to0$. The second bound improves Theorem~\ref{er:C}, whose right-hand side contains $(1+\gamma^{1/2})\hG^{3/2}+\epsd$.
\end{corollary}
\begin{proof}
$\ut-u_h=(\ut-u^L_h)+\delta_R$ with $\delta_R=\frac h\mu X$; Lemma~\ref{lk:lin} and Theorem~\ref{lk:up}. $\norm{\nabla\Ut}_{\Oh}\to\norm{\nabla U}_\Omega$ since $|\Oh\triangle\Omega|\le C\hG^2$.
\end{proof}
So the $H^1$ constant of the leak law of Corollary~\ref{er:Bp} is the Dirichlet energy of the Stokes flow driven by the leak, as in two dimensions \cite{PadhiLong2D}. In test P of Section~\ref{sec:numerics}, $\ut=0$ and $\Gc\equiv0$, so $u_h=-\delta_R$ and Theorem~\ref{lk:up} applies for every $\gamma$.

\begin{remark}[General data: the lower bound needs small $\gamma$]\label{lk:rem:gen}
For $\GS$ with $\ut\ne0$, the identity on $\Yh$ for $W':=\frac\mu h(\ut-u_h)-\Ut$ acquires the term $-\frac\mu h[\ip\ut{\dn v}+(F,v)]$. With the witness of Theorem~\ref{lk:low}, $\ip\ut{\dn v}=\ell_W(v)+O(\hG^{5/2})\norm{\nabla v}$ with $\ell_W$ the quadratic edge moments of Lemma~\ref{sh:lead}, of size $\hG^{3/2}\norm{\nabla v}$, so the extra term is $O(\gamma\hG^{1/2})\norm{\nabla v}$, the same order as the leak term. In two dimensions it is $O(\hG^2)$, because the sag is even and the witness odd on each edge. In three dimensions the witness cannot be chosen with vanishing quadratic moments on one Alfeld macro: the joint map (first moments, quadratic edge moments) on $\hat{\mathcal Y}^1_4$ has rank $6$, not $10$ (exact; on barycentric Worsey--Farin macros, rank $9$). So all we claim is $\norm{\nabla W'}\ge(c_LG_\hh-C\gamma\norm{|\hh|W}_{L^\infty})\hG^{1/2}-C(1+\gamma)\hG^{3/2}$, which is informative only for small $\gamma$. The limit statement of Corollary~\ref{lk:gen} is unaffected.
\end{remark}

\begin{corollary}[The constant converges at rate one, with a first-order three-dimensional term]\label{lk:K}
Let $K_h:=\norm{\nabla X}_{\Oh}/G$ and $d_h:=\norm{\nabla(X-\Ut)}/\norm{\nabla U}$. With $\sigma_U:=\dn\Ut-\Pt n$ on $\Gh$ and $\Theta_\zeta:=-\ip{\sigma_U}{\zeta}_{\Gh}$,
\begin{gather*}
\bigl|K_h^2-K^2(1+d_h^2)-2\Theta_\zeta/G^2\bigr|\le C\bigl(\hG^2+\hG\gamma^{-1/2}(1+\hG^{-1/2}E^0_U)\bigr),\\
\Theta_\zeta=\sum_F|F|\,q_F\,\sigma_U(x_F^\ast)\cdot\tau_F+O(\hG^2)=O(\hG).
\end{gather*}
Hence $|K_h-K|\le C(\hG+(E^0_U)^2)$ uniformly in $\gamma$, as in two dimensions; but the two-dimensional identity $K_h^2=K^2(1+d_h^2)+O(\hG^2)$ acquires the first-order term $2\Theta_\zeta/G^2$, which on the sphere is $O(\hG^2)$ if and only if the weighted centroid--circumcentre offsets $\sum_F|F|q_F\sigma_U\cdot(x_c-o_F)$ are.
\end{corollary}
\begin{proof}
$\norm{\nabla X}^2=\norm{\nabla\Ut}^2_{\Oh}+2(\nabla\Ut,\nabla W)+\norm{\nabla W}^2$ and $(\nabla\Ut,\nabla W)=(F_U,W)+\ip{\sigma_U}W$ (Green; $\div W=0$, $W|_\Sigma=0$, zero flux). On $\Gh$, $W=D+(\Pi_T\Ut-\Ut)-\Pi_T\zeta$. Now $\ip{\sigma_U}{\Pi_T\zeta}=\ip{\sigma_U}\zeta+\ip{\Pi_T\sigma_U-\sigma_U}\zeta$ with $\norm{\Pi_T\sigma_U-\sigma_U}\le C\hG$ (facewise smooth, jumps $O(\hG)$ across edges), and $\ip{\sigma_U}{\zeta}=-\Theta_\zeta$. By Lemma~\ref{lk:defect}, $\zeta|_F=-q_F(n(x^\ast)-n_F)+O(\hG^2)$, whose tangential face mean is $-q_F\tau_F$, while the affine part $S_F(x-x_c)$ integrates to zero against constants. Finally $|\ip{\sigma_U}{\Pi_T\Ut-\Ut}|\le C\hG^2$ and $|\ip{\sigma_U}D|\le C\norm D_{\Gh}$ (Theorem~\ref{lk:up}).
\end{proof}

\begin{remark}[What $\dn v\cdot n_F=-\div_Fv_F$ changes]\label{lk:rem:diff}
(1) The rate-$1$ mechanism of Theorem~\ref{er:C} and Lemma~\ref{lk:canc} survive: integration by parts on faces produces edge terms weighted by $|t_E\times(n_F-n_{F'})|=O(\hG)$, just as vertex terms were weighted by the tangent jump in two dimensions. (2) The upper rate $\hG^{1/2}$ is dimension-independent: a trace defect of amplitude $\hG$ oscillating on the scale $\hG$ has $H^{1/2}$ norm of order $\hG^{1/2}$ on a surface of any dimension. (3) The weight of the lower bound changes from the curvature to $|\hh|_{\rm F}$, because $\dn v$ is tangential on $F$ and pairs with the tangential defect $-qS_F(x-x_c)$ through a $2\times2$ moment matrix; only its symmetric part is seen, and the $k=3$ single-macro image is skew. (4) The face means of Remark~\ref{lk:rem:facemean} change no rate, only constants.
\end{remark}

\subsection{Drag, lift and torque}
Gjerde and Scott's motivating application is the drag on an obstacle. For a rigid motion $\psi$ (a translation for drag and lift, a rotation for torque), the force functional $J_\psi$ can be computed variationally from the discrete solution. For $\GS$ its error is dominated by the leak and equals $-(h/\mu)K_\psi$ to leading order, where the constant $K_\psi$ is obtained from the leak flow by reciprocity with a dual Stokes flow (Lemma~\ref{dr:recip}). For $\CNS$ the leak is absent and the error is of higher order.

Let $\RM:=\{\psi(x)=a+b\times x\}$; then $D(\psi)=0$, $\nabla\psi$ is skew, $\div\psi=0$ and $\int_{\Gh}\psi\cdot n=\int_\Gamma\psi\cdot n=0$. Put $J_\psi:=\int_\Gamma\sigma(u,p)n\cdot\psi$ with $\sigma(u,p)=D(u)-pI$ (drag and lift for $b=0$, torque for $a=0$). With $\Phi_\psi:=\frac12a\times x-\frac12|x|^2b$ ($\curl\Phi_\psi=\psi$) and a cutoff $\chi_0$ equal to $1$ on $\mathcal T_{r/2}$ and $0$ near $\Sigma$, let $V_\psi:=\curl(\chi_0\Phi_\psi)$, smooth, divergence-free and equal to $\psi$ near $\Gamma$. Let $\Lambda$ be the union of all tetrahedra meeting $\Gh$. For $h\le h_\ast$, $I_hV_\psi=\psi$ on $\Lambda$.

\begin{lemma}\label{dr:strip}
For $\hG\le h_\ast$, $\Oh\triangle\Omega\subset\Lambda$.
\end{lemma}
\begin{proof}
Points of $\Oh\triangle\Omega$ lie within $C\hG^2$ of $\Gh$. Near $\Gh$ the height above $\Gh$ in the frame (Gr) is affine up to $O(\hG^2)$ on sets of diameter $C\hG$; on a tetrahedron without vertices on $\Gh$ it is at least $c\hG-C\hG^2>C\hG^2$, because an interior vertex $z$ has $\dist(z,\partial\Oh)\ge ch_z$ and $h_z\ge c\hG$ near $\Gh$.
\end{proof}

\begin{definition}[Discrete drag]\label{dr:def}
For $w\in\VR$ let $\omega_h(w):=a_h(u_h,w)-(p_h,\div w)-(\ft,w)$, and $J^N_h(\psi):=\omega_h(I_hV_\psi)$. By the method equation, $J^N_h(\psi)=\int_{\Gh}t_h\cdot\psi+\ip{u_h}{\dn\psi}$ with the Nitsche flux $t_h=\dn u_h-\theta(p_h-\pbG(p_h))n-\frac\mu hu_h$ ($\theta=0$ for $\GS$, $1$ for $\CNS$), and $J^N_h(\psi)=\omega_h(w)$ for every $w\in\VR$ equal to $\psi$ on $\Lambda$. Let $v_0\in\VO$ with $\div v_0=\div I_hV_\psi$ ((IS3); $\norm{\nabla v_0}\le Ch^k$), $v_\psi:=I_hV_\psi-v_0\in\Zh$, and $\Theta_f(\psi):=(\ft,\psi)_{\Oh\setminus\Omega}-(f,\psi)_{\Omega\setminus\Oh}=O(\hG^2)$.
\end{definition}

\begin{lemma}[Error identity]\label{dr:id}
For $\GS$ and $\CNS$, with $e_u=\ut-u_h$,
\begin{gather*}
J^N_h(\psi)-J_\psi+\Theta_f(\psi)=-a_h(e_u,v_\psi)+\rho_0,\qquad\rho_0:=\ip{u_h}{\dn v_0}+(F,v_0),\\
|\rho_0|\le Ch^k\hG^{-1/2}\norm{u_h}_{\Gh}+C\hG^2h^k .
\end{gather*}
\end{lemma}
\begin{proof}
For $v_0\in\VO$ both methods give $\omega_h(v_0)=\ip{u_h}{\dn v_0}$. For $w\in\VR$, $\tilde\omega(w):=a_h(\ut,w)-(\pt,\div w)-(\ft,w)=\ip{\sigma(\ut,\pt)n}w-(F,w)$. For $w=I_hV_\psi$, which equals $\psi$ on $\Lambda\supset\Oh\triangle\Omega$ (Lemma~\ref{dr:strip}), the divergence theorem on the two signed regions between $\Gh$ and $\Gamma$, $\div(\sigma\psi)=(F-\ft)\cdot\psi$ there ($\sigma$ symmetric, $\nabla\psi$ skew, $F=0$ on $\Omega$), and $\int\psi\cdot n=0$ give $\tilde\omega(I_hV_\psi)=J_\psi-\Theta_f(\psi)$; and $\tilde\omega(v_0)=-(F,v_0)$. Since $\div v_\psi=0$, $\omega_h(v_\psi)-\tilde\omega(v_\psi)=-a_h(e_u,v_\psi)$.
\end{proof}

\begin{lemma}[Reciprocity]\label{dr:recip}
Let $(Z_\psi,R_\psi)$ be the Stokes flow in $\Omega$ with $Z_\psi|_\Gamma=\psi$ and $Z_\psi|_\Sigma=0$. Then $J_\psi(U):=\int_\Gamma\sigma(U,P)n\cdot\psi=\int_\Gamma q\,(R_\psi-\bar R_\psi)=:K_\psi$, with $\bar R_\psi$ the mean of $R_\psi$ over $\Gamma$.
\end{lemma}
\begin{proof}
Both pairs are smooth near $\Gamma$ and vanish on $\Sigma$; the weak Green formula ($\ip{\sigma(U,P)n}{W}_\Gamma=\frac12(DU,DW)-(P,\div W)$ for $W\in H^1$ with $W|_\Sigma=0$, smooth near $\Gamma$) gives $\ip{\sigma(U,P)n}{Z_\psi}_\Gamma=\frac12(DU,DZ_\psi)=\ip{\sigma(Z_\psi,R_\psi)n}U_\Gamma$, with no $H^2(\Omega)$ regularity. On $\Gamma$, $Z_\psi-\psi$ vanishes and is divergence-free, so $\dn(Z_\psi-\psi)\cdot n=0$, and $\dn\psi\cdot n=0$ by skewness; so $\sigma(Z_\psi,R_\psi)n\cdot n=-R_\psi$, and $U=-qn$ with $\int_\Gamma q=0$ gives the claim.
\end{proof}

\begin{theorem}[$\GS$: drag, lift and torque]\label{dr:gs}
For every $\gamma\ge\gamma_0$ and $\psi\in\RM$,
\[
J^N_h(\psi)=J_\psi-\frac h\mu K_\psi-\Theta_f(\psi)+\mathcal R_h,\qquad|\mathcal R_h|\le C\frac h\mu\bigl(\hG^{1/2}+E^0_U\bigr)+C\bigl(\hG^{3/2}+\varepsilon_1+h^k\hG^{-1/2}\bigr).
\]
For fixed $\gamma$, bounded $\rho$ and $E^0_U=O(h^{1/2})$: $J^N_h-J_\psi=-\frac h\mu K_\psi+O(h^{3/2})$, so the rate is exactly $1$ when $K_\psi\ne0$. The remainder is uniform in $\gamma$.
\end{theorem}
\begin{proof}
Lemma~\ref{dr:id}; $e_u=(\ut-u^L_h)+\delta_R$ (Lemma~\ref{lk:lin}), $a_h(\delta_R,v_\psi)=\frac h\mu a_h(\Ut,v_\psi)+\frac h\mu a_h(X-\Ut,v_\psi)$ (Theorem~\ref{lk:up}), and by Green's formula ($\div v_\psi=0$, $v_\psi|_{\Gh}=\psi$, $v_\psi=\psi-v_0$ on $\Lambda\supset\Oh\triangle\Omega$, $\div\sigma(\Ut,\Pt)=-F_U$), $a_h(\Ut,v_\psi)=(F_U,v_\psi)+\ip{\sigma(\Ut,\Pt)n}{\psi}_{\Gh}=J_\psi(U)-(F_U,v_0)$ with $|(F_U,v_0)|\le C\hG^2h^k$. Then Lemma~\ref{dr:recip}, and $\norm{u_h}_{\Gh}\le C$.
\end{proof}

For $\CNS$ we need the weighted normal slip. Let $w_\phi:=I_h(\chi_0(\phi\circ\pi)(n\circ\pi))$ for $\phi\in C^{k+1}(\Gamma)$ with $\int_\Gamma\phi=0$.
\begin{lemma}[Weighted normal slip of $\CNS$]\label{dr:slip}
Let $\gamma\ge\max(1,\gamma_2)$ and $Y:=(1+\gamma^{1/2})\hG^{3/2}+\epsd$. Then
\[
|\ip{\phi\circ\pi}{e_u\cdot n}_{\Gh}|\le C_\phi\Bigl[\frac h\mu\bigl(\hG^{-1/2}Y+h^k+\gamma\hG^3+\hG^2|\bar e_p|\bigr)+\hG\Bigl(\frac h\mu\Bigr)^{1/2}Y\Bigr],
\]
with $\bar e_p$ the mean of $e_p$ over $\Oh$ and $|\bar e_p|\le C$ if $\gamma\hG\le1$.
\end{lemma}
\begin{proof}
The two-dimensional proof \cite{PadhiLong2D} with three changes. (i) $|\Gc(w_\phi)|\le C_\phi(\hG+\gamma\hG^3)$ instead of $C_\phi(\hG^2+\gamma\hG^3)$: the transposed term is $\sum_F|F|\phi\,\dn u\cdot\tau_F+O(\hG^2)$ (Proposition~\ref{er:transposed}, Lemma~\ref{lk:defect}), generically $O(\hG)$; the penalty term is $O(\gamma\hG^3)$ because $\ut=d\,\partial_Nu(x^\ast)+O(d^2)$ is tangential to $\Gamma$ while $w_\phi$ is normal; and $|(F,w_\phi)|\le C\hG^2$. Since $\hG\le\hG^{-1/2}Y$ for $\gamma\ge1$, the change is absorbed. (ii) The cancellation for fields that are not divergence-free holds: $\dn\delta\cdot n_F=(\div\delta)|_{T_F}-\div_F\delta_F$, then the edge bookkeeping of Lemma~\ref{er:cancel}. (iii) $|\bar e_p|\le C$: test \eqref{er:cnserr} with $w_1:=I_h(\chi_0\,n\circ\pi)$; every term is bounded when $\gamma\hG\le1$, and $\int_{\Gh}w_1\cdot n=|\Gamma|+O(\hG^2)$.
\end{proof}

\begin{assumption}[(E$_Z$)]\label{dr:EZ}
With $\tilde Z_\psi$ the collar extension of Lemma~\ref{geom:ext} and $\tilde\zeta:=V_\psi-\tilde Z_\psi$, $E_Z:=\inf_{z\in\Zh}\tnorm{\tilde\zeta-z}\le C\hG^{1/2}$.
\end{assumption}
(E$_Z$) holds if $Z_\psi\in H^{3/2}(\Omega)$ and $h\le C\hG$: near $\Gamma$ everything is smooth and $\tilde\zeta=O(\hG^2)$ on $\Gh$, so only the $H^1$ approximation away from $\Gamma$ matters. It is therefore a regularity statement for the dual Stokes flow at $\Sigma$. For Lipschitz $\Sigma$ it is expected from $H^{3/2}$ regularity of the Stokes system on Lipschitz domains, and for the box from Dauge's $H^2$ theory on convex polyhedra \cite{Dauge1989}; we have checked neither (Remark~\ref{set:rem:corner}). Without (E$_Z$), $E_Z\to0$ still, and Theorem~\ref{dr:cns}(b) holds with $\gamma^{1/2}\hG^{3/2}E_Z=o(\hG^{3/2})$.

\begin{theorem}[$\CNS$: drag, lift and torque]\label{dr:cns}
Let $\gamma\ge\max(1,\gamma_2)$, $\gamma\hG\le1$, $\psi\in\RM$, and
\[
\Lambda_h(\psi):=\ip\ut{\dn(\tilde Z_\psi-\psi)}_{\Gh}=\sum_F\int_Fd\,\bigl(\partial_Nu\cdot\dn(Z_\psi-\psi)\bigr)(x^\ast)\,dA+O(\hG^3).
\]
Then (a) $|J^N_h-J_\psi+\Theta_f|\le C(Y+h^k\hG^{-1/2})$, the rate $\frac32$; and (b)
\[
|J^N_h(\psi)-J_\psi+\Theta_f(\psi)+\Lambda_h(\psi)|\le C\bigl[\gamma^{-1/2}\hG^2+\gamma\hG^3+\gamma^{1/2}\hG^{3/2}E_Z+\epsd+h^k\hG^{-1/2}\bigr].
\]
Modulo (E$_Z$), $|J^N_h-J_\psi|\le C(\hG^2+\epsd+h^k\hG^{-1/2})$: the rate is $2$.
\end{theorem}
\begin{proof}
The two-dimensional proof \cite{PadhiLong2D}, with $v_\psi$ of Definition~\ref{dr:def} and the extra $\rho_0$ of Lemma~\ref{dr:id}. The auxiliary problem $N_h(\zeta_h,w)+B^\ast(\pi_h,w)=a_h(v_\psi,w)$ is the $\CNS$ system for $(\tilde\zeta,-\tilde R_\psi)$, since $a_h(v_\psi,w)=(-\Delta V_\psi,w)+a_h(v_\psi-V_\psi,w)$ ($D(V_\psi)=0$ near $\Gh$, $\norm{\nabla(v_\psi-V_\psi)}\le Ch^k$); Theorem~\ref{er:D} gives existence and $\tnorm{\tilde\zeta-\zeta_h}\le C[(1+\gamma^{1/2})\hG^{3/2}+E_Z+h^k]$. The remaining dimension-dependent inputs are: $|\tilde\zeta|\le C\hG^2$ and $|(\nabla\ut)^{\mathsf T}n_F|\le C\hG$ on $\Gh$; $\tilde\zeta\cdot n_F=d\,\dn\zeta(x^\ast)\cdot n_F+O(\hG^4)=O(\hG^3)$ by Lemma~\ref{geom:wall} applied to $\zeta$; $\norm{e_p-\pbG(e_p)}_{\Gh}\le C\hG^{-1/2}Y$ (proof of Theorem~\ref{er:D} and Lemma~\ref{st:inverse}); and Lemma~\ref{dr:slip} with $\phi=R_\psi-\bar R_\psi$.
\end{proof}
With $d=Q_F+O(\ell^3)$ (Lemma~\ref{geom:face}(a)), $\int_FQ_F=\frac{|F|}{24}\sum_{i<j}w^F_{ij}$. On the unit sphere, $\Theta_f+\Lambda_h\approx-\sum_F\frac{|F|}{24}\bigl(\sum_{i<j}\ell_{ij}^2\bigr)\bigl(f\cdot\psi+\partial_Nu\cdot\dn(Z_\psi-\psi)\bigr)(x^\ast_F)$: the leading drag error of $\CNS$ is the shape derivative of $J_\psi$ under the chordal sag. Whether rate $2$ is sharp is not proved.

\section{The penalty threshold scales with \texorpdfstring{$\rho$}{rho}}\label{sec:penalty}

Sufficiency is Lemma~\ref{st:coercive}: $\mu\ge\mu_0=4\kappa C_I^2\rho$ implies coercivity of $N_h$ on $\VR$. This section shows that a threshold proportional to $\rho$ is also necessary, on the divergence-free space $\Zh$ on which the methods actually need coercivity, under conditions on the cells at a smallest boundary face.

\paragraph{Scaling.} For $v\in\Zh$ put $A:=a_h(v,v)$, $B:=\ip{\dn v}v$ and $C:=\norm v^2_{\Gh}$, so that $N_h(v,v)=A-2B+(\mu/h)C$. Under dilation by $\ell$, $A$ and $B$ scale like $\ell$ and $C$ like $\ell^2$. For a field of length scale $\ell$,
\begin{equation}\label{pen:scale}
N_h(v,v)=\frac C\ell\Bigl[\frac{\mu\ell}h-\frac{\ell(2B-A)}C\Bigr],
\end{equation}
and the quotient $\ell(2B-A)/C$ is dilation invariant. So $N_h$ is indefinite on $\Zh$ as soon as some $v\in\Zh$ has $\ell(2B-A)/C>\mu\ell/h$.

\subsection{A witness in one Alfeld macro-element}
Let $\Th$ be an Alfeld refinement, $T=\conv(P_0,\dots,P_3)$ a macro-tetrahedron with $F=\conv(P_1,P_2,P_3)\in\FG$ and $T\cap\Sigma=\emptyset$, and
\[
W_k(T):=\bigl\{v\in C^0(T)^3:\ v|_S\in P_k^3\ \forall S\in T^A,\ \div v=0,\ v=0\text{ on }\partial T\setminus F\bigr\},
\]
with $A_T,B_T,C_T$ the forms above restricted to $T$ and $F$, and
\[
\lambda_T:=\diam F\cdot\sup_{v\in W_k(T),\,C_T(v)>0}\frac{2B_T-A_T}{C_T}.
\]

\begin{lemma}[Piola invariance]\label{pen:piola}
Let $\Phi(\hat x)=J\hat x+b$ map $\hat T$ onto $T$ with $\hat F\mapsto F$. Then $\hat v\mapsto(\det J)^{-1}J\,\hat v\circ\Phi^{-1}$ is a bijection $W_k(\hat T)\to W_k(T)$. Hence $\dim W_k(T)$ does not depend on the shape. Pulled back to $W_k(\hat T)$, $A_T$ and $B_T$ are rational functions of $J$, and $C_T$ is a rational function of $J$ times $|J^{-\mathsf T}\hat\nu|$, the square root of a rational function; all three are continuous where $\det J\ne0$. So a strict inequality $(2B_T-A_T-cC_T)(v)>0$ at a fixed pulled-back witness persists on an open set of shapes.
\end{lemma}

\begin{proof}
$\Phi$ maps barycentre to barycentre, hence $\hat T^A$ onto $T^A$; continuity, degree and vanishing on faces are preserved, and $\div v=(\det J)^{-1}(\widehat\div\,\hat v)\circ\Phi^{-1}$. In $B_T$ the normalisation of $n$ cancels against Nanson's factor; in $C_T$ it does not.
\end{proof}

\begin{theorem}[Single-cell necessity]\label{pen:single}
If $\mu\diam F/h<\lambda_T$, then $N_h$ is indefinite on $\Zh$, for every $k\ge3$. In particular, if the macro-tetrahedron on a smallest face $F_\ast$ of $\Gh$ has $\lambda_T\ge\lambda_0>0$, coercivity on $\Zh$ requires $\mu\ge\lambda_0\rho$. Exact rational certificates give, for $k=3$:
\begin{center}\footnotesize
\begin{tabular}{lcccc}
\toprule
shape (apex) & $\dim W_3$ & $\sup\frac{2B-A}C$ & $\lambda_T$ & exact certificate\\
\midrule
near-regular, $(\frac12,\frac7{24},\frac{13}{16})$ & 8 & 32.2481 & 32.50 & $2B-A-\frac{28060}{871}C>0$\\
right corner$^\dagger$, $(0,0,1)$ & 8 & 22.9834 & 32.50 & $2B-A-\frac{9850}{429}C>0$\\
flat, height $\frac12$ & 8 & 61.4449 & 61.92 & $2B-A-\frac{51071}{832}C>0$\\
flat, height $\frac14$ & 8 & 134.356 & 135.40 & $2B-A-\frac{133014}{991}C>0$\\
skew, $(\frac32,\frac13,\frac34)$ & 8 & 32.3856 & 32.64 & $2B-A-\frac{14656}{453}C>0$\\
tall, height $2$ & 8 & $-10.381$ & $-10.46$ & $A-2B-10C\succ0$\\
tall, height $3$ & 8 & $-57.454$ & $-57.90$ & $A-2B-57C\succ0$\\
\bottomrule
\end{tabular}
\end{center}
The base is $(0,0,0),(1,0,0),(\frac12,\frac78,0)$ except for $^\dagger$, where it is $(0,0,0),(1,0,0),(0,1,0)$; flat and tall cells have their apex above $(\frac12,\frac7{24},0)$. By Lemma~\ref{pen:piola}, each positive certificate also holds on an open set of shapes around the tabulated one. \status{computer-assisted, exact}
\end{theorem}

\begin{proof}
A field $v\in W_k(T)$ extended by zero is continuous (it vanishes on the faces shared with other cells), piecewise $P_k$, zero on $\Sigma$ and divergence free, so $v\in\Zh$, and its forms are $A_T,B_T,C_T$. Apply \eqref{pen:scale} with $\ell=\diam F$ to the maximiser, and $h/\diam F_\ast=\rho$.
\end{proof}
\cert{notes/v6/td\_pen/alfeld\_tet.py, verify\_alfeld.py, neg\_cert.py}{alfeld\_k3.log, verify\_*.log, neg\_cert.log}
The null space is computed exactly over $\mathbb Q$; the suprema are found by floating-point bisection on the $8$-dimensional pencil, and the \emph{certificates} are rational witnesses checked in exact arithmetic. An independent quadrature check reproduces the near-regular quotient $32.24809277$ against $32.24809278$ from the exact pipeline. For tall cells, the exact $LDL^{\mathsf T}$ factorisation shows that $A-2B-10C$ (height $2$) is positive definite on $W_3(T)$, so no single-cell witness of degree $3$ exists (higher degrees were not tested); a height scan puts the change of sign at height $\approx1.69$ times the base size (the regular tetrahedron has height $0.816$). This is the three-dimensional analogue of the tall-cell limitation in two dimensions \cite{PadhiLong2D}.

\subsection{A continuum witness}
\paragraph{Half-space constant.} Take the fluid in $\{z>0\}$, $n=(0,0,-1)$ and $\dn=-\partial_z$. Let $g(t)=t(1-t)^6$ on $[0,1]$ (zero for $t\ge1$), $\phi(s)=(1-s^2)^6$ on $[-1,1]$, and
\[
\mathcal F:=Y\phi(x/X)\phi(y/X)g(z/Y),\qquad w:=\curl(\mathcal Fe_y)=(-\partial_z\mathcal F,0,\partial_x\mathcal F).
\]
Then $\div w=0$ and $w\cdot n=0$ on $z=0$; at $z=0$, $w_1=-\phi\phi$ and $\dn w\cdot w=12\phi^2\phi^2/Y>0$. Exact integration gives, for $X=2Y$,
\[
\hat\lambda:=\frac{Y(2B-A)}C=\frac{2327591}{151008}=15.41369\ldots,
\]
with $15.53940$, $15.58222$, $15.62292$ for $X/Y=3,4,8$, and $\hat\lambda\to-2g'(0)g''(0)-\int_0^1|g''|^2=\frac{172}{11}=15.636\ldots$ as $X/Y\to\infty$ (the one-dimensional limit). \status{exact}
\cert{notes/v6/td\_pen/continuum3d.py}{continuum3d.log}

\begin{enumerate}[label=(LQ)]
\item Let $F_\ast\in\FG$ have minimal diameter, $x_0\in F_\ast$, and $Y:=C_q\diam F_\ast/\kappa$. Every tetrahedron of $\Th$ meeting $B(x_0,4Y)$ has diameter at most $C_q\diam F_\ast=\kappa Y$, and $B(x_0,4Y)\cap\Sigma=\emptyset$.
\end{enumerate}

\begin{theorem}[Continuum-witness necessity]\label{pen:cont}
Assume (IS3) with constant $\beta$ (on Alfeld refinements, $k\ge3$, Lemma~\ref{st:IS3}). There are $\kappa_0,Y_0>0$, depending on $C_q$, $\beta$, $k$, $\Gamma$ and shape regularity, such that under (LQ) with $\kappa\le\kappa_0$ and $Y\le Y_0$, $N_h$ is indefinite on $\Zh$ whenever $\mu<15h/Y$. Hence the coercivity threshold on $\Zh$ satisfies $\mu^\ast\ge(15\kappa_0/C_q)\rho$. No condition on the shapes of the boundary cells beyond shape regularity is needed.
\end{theorem}

\begin{proof}
Constants $C$ depend on the profile, $k$, $\Gamma$ and shape regularity.
\begin{enumerate}
\item \emph{Field.} Use the frame (Gr) at $x_0^\ast:=\pi(x_0)$, with $z$ pointing into $\Omega$, $X=2Y$, and the field $w$ above; for $z<0$ extend $g$ polynomially and multiply by a cutoff equal to $1$ on $z\ge-Y/2$. Then $w\in W^{4,\infty}$ with $|D^jw|\le C_jY^{-j}$, $\div w=0$ and $\operatorname{supp}w\subset B(x_0,4Y)$. Near $x_0$, $\Gh$ is a graph $z=\gamma_h(x,y)$ with $|\gamma_h|\le CY^2$ and $|\nabla\gamma_h|\le CY$ (Corollary~\ref{geom:graph} and (Gr)).
\item \emph{Comparison with the half-space.} On $\operatorname{supp}w$, $\Oh$ and $\{z>0\}$ differ in a set of volume $\le CY^4$ where $|\nabla w|^2\le CY^{-2}$, and moving the integrands of $B$ and $C$ from $z=0$ to $z=\gamma_h$ (with $n_F$ in place of $(0,0,-1)$) changes them by $O(1)$ and $O(Y)$ pointwise. So $|A_{\Oh}(w)-Y\hat A|+|B_{\Gh}(w)-Y\hat B|\le CY^2$ and $|C_{\Gh}(w)-Y^2\hat C|\le CY^3$, with $\hat A,\hat B,\hat C$ the half-space values at $Y=1$.
\item \emph{Discrete field.} Let $I$ be degree-$3$ Lagrange interpolation, so $Iw\in\Vh$ and $\norm{w-Iw}_{L^\infty}\le C\kappa^4$, $\norm{\nabla(w-Iw)}_{L^\infty}\le C\kappa^3/Y$. Let $b:=\sum_{F\subset\Gh\cap B(x_0,Y)}b_Fn_F$; then $\int_{\Gh}b\cdot n\ge cY^2$, $|\nabla b|\le CC_q/(\kappa Y)$ and $\norm{\nabla b}^2\le CC_q^2Y/\kappa$. Since $\int_{\Gh}w\cdot n=\int_{\Oh}\div w=0$, $m:=\int_{\Gh}Iw\cdot n$ satisfies $|m|\le CY^2\kappa^4$; with $a:=m/\int_{\Gh}b\cdot n$, $|a|\le C\kappa^4$. (IS3) gives $v_0\in\VO$ with $\div v_0=\div(Iw-ab)$ and $\norm{\nabla v_0}\le\beta^{-1}\sqrt3\bigl(\norm{\nabla(Iw-w)}+|a|\norm{\nabla b}\bigr)\le C\beta^{-1}(\kappa^3+C_q\kappa^{7/2})Y^{1/2}$. Put $v:=Iw-ab-v_0\in\Zh$.
\item \emph{Perturbation.} On $\Gh$, $v=Iw-ab$. So $|C(v)^{1/2}-C_{\Gh}(w)^{1/2}|\le C\kappa^4Y$ and $|A(v)^{1/2}-A_{\Oh}(w)^{1/2}|\le C(1+\beta^{-1})C_q\kappa^3Y^{1/2}$. In $B$, the terms $\dn(Iw-w)$ and $a\dn b$ change $B$ by at most $C(1+C_q)\kappa^3Y$, and by the inverse trace inequality and $\sum_F|F|/\diam F\le CC_qY/\kappa$, $|\ip{\dn v_0}{Iw-ab}|\le C(C_qY/\kappa)^{1/2}\norm{\nabla v_0}\le C\beta^{-1}C_q^{3/2}\kappa^{5/2}Y$.
\item \emph{Conclusion.} $Y(2B(v)-A(v))/C(v)\ge\hat\lambda-C[(1+\beta^{-1})C_q^{3/2}\kappa^{5/2}+Y]$. Choose $\kappa_0,Y_0$ so that the bracket is at most $0.4$; then $\mu<15h/Y$ gives $N_h(v,v)<0$ by \eqref{pen:scale}. Finally $Y=C_q\diam F_\ast/\kappa$ and $h/\diam F_\ast=\rho$.\qedhere
\end{enumerate}
\end{proof}

\begin{corollary}[The threshold scales with $\rho$]\label{pen:rho}
On Alfeld refinements with $k\ge3$, the coercivity threshold $\mu^\ast$ of $N_h$ on $\Zh$ satisfies $c\,\rho\le\mu^\ast\le4\kappa C_I^2\rho$ whenever the mesh is locally quasi-uniform over $O(1/\kappa_0)$ layers at a smallest boundary face in the sense of (LQ), or the macro-cell on a smallest boundary face has a certified $\lambda_T>0$ (Theorem~\ref{pen:single}).
\end{corollary}

The constant $\kappa_0$ in Theorem~\ref{pen:cont} is not explicit, and the statement without (LQ) or a shape condition is open, as in two dimensions. If commuting projections onto the Alfeld velocity space with $W^{1,\infty}$-local bounds were used in step~3 in place of the interpolate-and-correct construction \cite{FuGuzmanNeilan2020}, (IS3) would not be needed; we have not checked those bounds.

In our computations (Section~\ref{sec:numerics}) the threshold on $\Zh$ grows with $h/\hG$: $\mu^\ast=76.3$ and $148.0$ at the two sphere levels for $k=3$, $193.6$ at the first level for $k=4$.

\section{Steady Navier--Stokes flow}\label{sec:ns}

Gjerde and Scott's application is cylinder drag, a Navier--Stokes computation \cite{GjerdeScott2024}. Consider
\[
-\nu\Delta u+(u\cdot\nabla)u+\nabla p=f,\qquad\div u=0\ \text{in }\Omega,\qquad u|_\Gamma=0,\quad u|_\Sigma=g,
\]
and, with $c_1(w;u,v):=((w\cdot\nabla)u,v)$ or $c_s(w;u,v):=\frac12[((w\cdot\nabla)u,v)-((w\cdot\nabla)v,u)]$ (denoted $c$), the methods: find $u_h\in\Vh$ with $u_h|_\Sigma=\gI$ and $p_h\in\Pih$ such that for all $v\in\VR$, $q\in\Pih$,
\[
\nu N_h(u_h,v)+c(u_h;u_h,v)-(p_h,\div v)\;\bigl[+\ip{p_h-\pbG(p_h)}{v\cdot n}\bigr]=(\ft,v),\qquad(q,\div u_h)=0 .
\]
Without the bracket this is $\GS$, with it $\CNS$. The wall penalty is then $\nu\mu/h$ (see \cite{PadhiLong2D} for the conversion to an unscaled penalty).

\paragraph{Assumptions.} (NS1) $(u,p)$ is a solution with $u\in C^{k+1}(\overline\Omega)$, $p\in C^k(\overline\Omega)$, $\ft$ is a smooth extension of $f$, and the assumptions of Section~\ref{sec:setting} hold. (NS2) $\ut,\pt$ are the extensions of Lemma~\ref{geom:ext}; the defect $F:=\ft+\nu\Delta\ut-(\ut\cdot\nabla)\ut-\nabla\pt$ vanishes on $\Omega$, is supported in $\overline{\Oh\setminus\Omega}$ and is bounded. Let $\Aop(w,v):=\nu N_h(w,v)+c(w;\ut,v)+c(\ut;w,v)$ be the linearisation at $\ut$, $\norm\ell_{\ast,X}:=\sup_{0\ne v\in X}\ell(v)/\tnorm v$, and $B^\ast$ as in the proof of Theorem~\ref{er:D}. The discrete stability conditions are:
\begin{enumerate}
\item[(L)] there is $\beta_L>0$, independent of $h$, with $\beta_L\tnorm\delta\le\norm{\Aop(\delta,\cdot)}_{\ast,\Zh}$ for all $\delta\in\Zh$;
\item[(L$^\ast$)] there is $\beta_\ast>0$, independent of $h$, such that for every linear functional $\ell$ on $\VR$ the problem $(\delta,\xi)\in\Zh\times\Pih$, $\Aop(\delta,v)+B^\ast(\xi,v)=\ell(v)$ for all $v\in\VR$, has a unique solution with $\beta_\ast\tnorm\delta\le\norm\ell_{\ast,\VR}$.
\end{enumerate}
They are the discrete forms of ``$u$ lies on a nonsingular branch'' for $\GS$ and $\CNS$. The strategy is the classical one for nonsingular branches \cite{BRR1980}: a fixed-point argument around the linearisation at $\ut$, in which the Stokes estimates of Section~\ref{sec:error} control the consistency error and (L) or (L$^\ast$) replaces coercivity. The work lies in checking that every dimension-dependent ingredient of the two-dimensional proof has a three-dimensional replacement; the proof below lists them. Let $X_3:=h/\mu+(1+\gamma^{1/2})\hG^{3/2}+\epsd$.

\begin{corollary}[Navier--Stokes in three dimensions]\label{ns:cor}
Assume (NS1)--(NS2) and (IS3); on Alfeld refinements with $k\ge3$ (IS3) holds by Lemma~\ref{st:IS3}.
\begin{enumerate}[label=(\roman*)]
\item \emph{Small data.} Let $N$ be the constant in $|c(w;u,v)|\le N\norm w_{H^1}\norm u_{H^1}\norm v_{H^1}$ and $C_1$ that in $\norm v_{H^1(\Oh)}\le C_1\tnorm v$ for $v\in\VR$; both are uniform in $h$. If $\gamma\ge\gamma_0$ and $\lambda:=NC_1^2\norm\ut_{H^1(\Oh)}/(\nu c_1)\le\frac14$, with $c_1$ the coercivity constant of Lemma~\ref{st:coercive}, then (L) holds with $\beta_L=\nu c_1/2$, and (L$^\ast$) holds for $\gamma\ge\gamma_2'$ with $\beta_\ast=\nu/C_\ast$, $\gamma_2'$ and $C_\ast$ independent of $\nu$ and $h$.
\item \emph{Nonsingular solutions.} If $u$ is a nonsingular solution (the linearisation at $u$ is an isomorphism of the divergence-free $H^1_0$ fields onto their dual), then there is $h_0$, independent of $\mu,\gamma,\rho$, such that (L), and (L$^\ast$) for $\gamma\ge\gamma_2'$, hold for $h\le h_0$.
\item \emph{$\GS$.} Under (L), for either convective form and $X_3\le c_\ast$: $\GS$ has a locally unique solution, $\norm{\nabla(\ut-u_h)}\le C[h/\mu+(1+\gamma^{1/2})\hG^{3/2}+\epsd]$, $\tnorm{\ut-u_h}\le\nu^{-1}(h/\mu)^{1/2}G+CX_3$, and under the hypotheses of Corollary~\ref{er:Bp}
\[
\ut-u_h=\frac h{\nu\mu}v_\varphi+\vartheta_h,\qquad\tnorm{\vartheta_h}=o\bigl((h/\mu)^{1/2}\bigr),
\]
so that $u_h\cdot n\approx(h/\nu\mu)(p-\pbar)$ on $\Gh$ with the Navier--Stokes pressure.
\item \emph{$\CNS$.} Under (L$^\ast$) and $(1+\gamma^{1/2})\hG^{3/2}+\epsd\le c_\ast$, $\CNS$ has a locally unique solution with $\tnorm{\ut-u_h}\le C[(1+\gamma^{1/2})\hG^{3/2}+\epsd]$, and $\norm{\ut-u_h}_{H^1(\Oh)}\le C(\hG^{3/2}+h^k)$ for bounded $\gamma$ and $\hG\ge h^2$.
\end{enumerate}
\end{corollary}

\begin{proof}
The two-dimensional proofs \cite{PadhiLong2D} consist of a fixed-point argument around the linearisation at $\ut$, the Stokes estimates of Section~\ref{sec:error}, and, for (ii), a compactness argument. We list every dimension-dependent input and its three-dimensional replacement.
\begin{enumerate}
\item \emph{Sobolev constants.} $H^1\subset L^4$ in three dimensions; transferred to $\Oh$ by the bi-Lipschitz map $\Phi_h$ of Lemma~\ref{geom:korn}, the constant $N$ in $|c(w;u,v)|\le N\norm w_{H^1}\norm u_{H^1}\norm v_{H^1}$ is uniform in $h$. The convective boundary identities and the wall P\'eclet estimate are dimension-free (divergence theorem; $\norm v^2_{\Gh}\le(h/\mu)\tnorm v^2$).
\item \emph{Error equation.} The identity is dimension-free. The geometric term $\nu\Gc$ is bounded by Lemma~\ref{er:G} (the Navier--Stokes defect enters only through $\norm F_\infty$ and its support), and the convective consistency term by $C\hG^4\norm{\nabla v}$, since $\norm\ut_{L^\infty(\Gh)}\le C\hG^2$.
\item \emph{(i).} Coercivity, continuity and the inverse traces are Lemmas~\ref{st:inverse}, \ref{st:coercive}; the inf-sup condition is (IS3); the square count and the uniqueness of the constant use the face bubble (Lemma~\ref{st:divrange}). For (L), $\nu N_h(\delta,\delta)\ge\nu c_1\tnorm\delta^2$ for $\gamma\ge\gamma_0$, and each of the two convective terms of $\Aop(\delta,\delta)$ is at most $\lambda\nu c_1\tnorm\delta^2\le\frac14\nu c_1\tnorm\delta^2$. For (L$^\ast$), $\Aop$ is bounded by $\nu(M+c_1/2)$ in the energy norm, so the absorption has the same form as in two dimensions with every term carrying the factor $\nu$, and $\gamma_2'$ depends only on $C_I,c_0,\beta,c_1,M$.
\item \emph{(ii).} The compactness argument uses: a strip of volume $|S_h|\le C\hG^2$ (Corollary~\ref{geom:graph}); trace transfer by $\Phi_h$; discrete test fields $I_h\varphi-w_h$ with $w_h\in\VO$ from (IS3) and $|\ip{\delta_j}{\dn w_h}|\le(\rho/\mu)^{1/2}C_I\norm{\nabla w_h}$; Rellich's theorem and compact traces on the fixed Lipschitz domain $\Omega$; density of smooth solenoidal fields in the divergence-free $H^1_0$ fields on Lipschitz domains in $\R^3$ \cite{Temam}; and compactness of $L^2\hookrightarrow V'$. All hold in three dimensions.
\item \emph{(iii).} Steps of the $\GS$ proof: the approximation step uses Lemmas~\ref{er:G} and~\ref{er:approx}; the leak step uses $m(v)=\Rc(v)-\ip{v_\varphi}v$, bounded only in the energy dual, $|m(v)|\le C_p\hG(h/\mu)^{1/2}\tnorm v$ (proof of Theorem~\ref{er:C}), which suffices because (L) is stated in the energy dual; the load $\ell$ is bounded by Lemma~\ref{er:cancel} with $S=w$ and Lemma~\ref{er:vphi}, and the convective term by $C(h/\mu)\norm{\nabla v}$ since $\norm{v_\varphi}_{H^1}\le C$. This gives $\tnorm r\le C\beta_L^{-1}(h/\mu+\hG^2+\epsd)\le CX_3$. The leading-term statement follows as Corollary~\ref{er:Bp}, with \eqref{er:riemann}.
\item \emph{(iv).} $|B^\ast(\eta,v)|\le2\norm\eta_{\Gh}(h/\mu)^{1/2}\tnorm v$ with $\norm\eta_{L^\infty(\Gh)}\le C\hG^k$; the fixed point is as for $\GS$ with $\beta_\ast$; the $H^1$ statement uses the flux proviso of Theorem~\ref{er:D}.\qedhere
\end{enumerate}
\end{proof}

\begin{remark}\label{ns:rem}
An explicit threshold $h_0$ in (ii) by a Schatz-type duality argument \cite{Schatz1974} is possible in two dimensions under Stokes $H^2$-regularity \cite{PadhiLong2D}. In three dimensions the same argument needs the corresponding regularity on $\Omega$, which for a box minus an obstacle involves the edges and corners of the box \cite{Dauge1989}; it also needs a divergence-free extension of the $H^2$ dual solution across $\Gamma$, for which the Piola--Hestenes construction of Lemma~\ref{geom:ext} applies with Sobolev coefficients. We have not written this out, and the edges and corners of the box limit the regularity (Remark~\ref{set:rem:corner}). %% VERIFY: Dauge1989 H^2 x H^1 on convex polyhedra; statement not checked.
The Oseen leak limit and the drag statements are in Section~\ref{ns:sec:oseen}; none of them uses this global regularity.
\end{remark}

\subsection{The Oseen leak limit and drag}\label{ns:sec:oseen}
For Navier--Stokes flow the leak of the printed method is driven by the same wall pressure mismatch, but it is transported by the linearised flow. The limit of the scaled error is therefore the \emph{Oseen} leak flow, the solution of the linearisation at $u$ with Dirichlet data $-(p-\pbar)n$ on the wall, and the Stokes arguments of Section~\ref{sec:leak} carry over once the discrete Oseen operators are uniformly stable.
Assume (NS1)--(NS2), (IS3), and that $u$ is a nonsingular solution (Corollary~\ref{ns:cor}(ii)); write (H$_{\rm ns}$) for the last assumption. Let $K(w,v):=c(w;\ut,v)+c(\ut;w,v)$, so $\Aop=\nu N_h+K$. Besides (L) and (L$^\ast$) we use two further discrete stability conditions:
\begin{enumerate}
\item[(L$_D$)] there is $\beta_D>0$, independent of $h$, with $\beta_D\norm{\nabla y}\le\sup_{0\ne v\in\Yh}(\nu a_h(y,v)+K(y,v))/\norm{\nabla v}$ for all $y\in\Yh$ (the discrete Oseen--Dirichlet problem);
\item[(L$^\dagger$)] the condition (L$^\ast$) for the adjoint linearisation $(\delta,v)\mapsto\Aop(v,\delta)$.
\end{enumerate}
The compactness argument of Corollary~\ref{ns:cor}(ii) gives (L), (L$^\ast$) and (L$^\dagger$) for $\gamma\ge\gamma_2'$, and (L$_D$), for $h\le h_0$ with $h_0$ independent of $\mu,\gamma,\rho$; the two-dimensional proofs are in \cite{PadhiLong2D}, and the dimension-dependent inputs are those listed in the proof of Corollary~\ref{ns:cor}.
Let $(U_O,P_O)$ be the Oseen leak flow: $-\nu\Delta U_O+(U_O\cdot\nabla)u+(u\cdot\nabla)U_O+\nabla P_O=0$, $\div U_O=0$ in $\Omega$, $U_O|_\Gamma=-qn$ with $q=(p-\pbar)|_\Gamma$ the Navier--Stokes wall pressure, and $U_O|_\Sigma=0$. It exists by (H$_{\rm ns}$) after lifting the data with Lemma~\ref{er:w}, and it is $C^{2,\alpha}$ on a collar of $\Gamma$ (local bootstrap with the Stokes Schauder estimate; no global regularity is used). $\Ut_O$, $E^0_O$ and $F_O$ are defined as $\Ut$, $E^0_U$ and $F_U$ in Section~\ref{sec:leak}. Let $U_\infty:=\norm\ut_{W^{1,\infty}(\Oh)}$ and $C_F$ the constant in $\norm w_{L^2(\Oh)}\le C_F\norm{\nabla w}$ for $w$ vanishing on $\Sigma$ (uniform in $h$ by Lemma~\ref{geom:korn}).

\begin{theorem}[Oseen and Navier--Stokes versions]\label{ns:oseen}
For $h\le h_0$ and either convective form:
\begin{enumerate}[label=(\roman*)]
\item \emph{Oseen leak limit, all $\gamma\ge\gamma_0$.} With $\Aop(\delta_R,v)=\Rc(v)$ on $\Zh$ and $X:=\frac{\nu\mu}h\delta_R$: $\norm{\nabla(X-\Ut_O)}\le C(\hG^{1/2}+E^0_O)$.
\item \emph{Lower bound} (Alfeld, $k\ge4$, every $\mu$ for which $\delta_R$ exists). $\norm{\nabla(X-\Ut_O)}\ge c_O\,G_\hh\hG^{1/2}-C\hG^{3/2}$ with $c_O:=\nu\Lambda_0c_L/(\nu\Lambda_0+C_F(1+C_F)U_\infty)$, $G_\hh$ and $c_L$ as in Theorem~\ref{lk:low}.
\item \emph{$\GS$, uniformly in $\gamma$.} If $X_u:=h/\mu+\hG^{3/2}+\varepsilon_1\le c_\ast$, then $\GS$ has a locally unique solution with $\norm{\nabla(\ut-u_h)}\le CX_u$ and
\[
\Bigl\|\nabla\Bigl(\frac{\nu\mu}h(\ut-u_h)-\Ut_O\Bigr)\Bigr\|\le C(\hG^{1/2}+E^0_O)+C\frac\mu h\bigl(\hG^{3/2}+\varepsilon_1+X_u^2\bigr).
\]
So the Navier--Stokes $H^1$ constant $K_{NS}=\norm{\nabla U_O}/G$ is attained whenever $\gamma\hG^{1/2}\to0$ and $\frac\mu h\varepsilon_1\to0$.
\item \emph{$\GS$ drag.} $J^N_h(\psi)=J_\psi-\frac h{\nu\mu}K^\dagger_\psi-\Theta_f(\psi)+O\bigl(\frac h\mu(\hG^{1/2}+E^0_O)+\hG^{3/2}+\varepsilon_1+X_u^2+h^k\hG^{-1/2}\bigr)$, where $K^\dagger_\psi=\int_\Gamma q(R^\dagger_\psi-\bar R^\dagger_\psi)$ and $(Z^\dagger_\psi,R^\dagger_\psi)$ is the adjoint Oseen flow with data $\psi$ on $\Gamma$ and $0$ on $\Sigma$.
\item \emph{$\CNS$ drag.} Theorem~\ref{dr:cns} holds with $\nu$ and $Z^\dagger_\psi$ in place of $Z_\psi$ (the term $\Lambda^\dagger_h$), under $\gamma\ge\max(1,\gamma_2')$, $\gamma\hG\le1$ and $Y\le c_\ast$; the rate $2$ holds modulo the analogue (E$_Z^\dagger$) of (E$_Z$).
\end{enumerate}
\end{theorem}
\begin{proof}
(i) is Theorem~\ref{lk:up} with $a_h$ replaced by $\nu a_h+K$: the Oseen--Dirichlet field $X^D$ uses (L$_D$), the energy step uses (L), and $|K(w,v)|\le C_K\norm w\norm{\nabla v}+\frac14U_\infty\hG^2\norm w_{\Gh}\norm v_{\Gh}$ (the convective boundary identity; dimension-free apart from $C_F$, via $H^1\subset L^4$ and $\Phi_h$). The lift, projection and cancellation lemmas are Lemmas~\ref{lk:lift}--\ref{lk:canc}. (ii) Lemma~\ref{lk:id} becomes $\nu a_h(W,v)+K(W,v)-\nu\ip W{\dn v}=\nu\ip{\Ut_O}{\dn v}-(F_O,v)$ on $\Yh$, with $|K(W,v)|\le C_F(1+C_F)U_\infty\norm{\nabla W}\norm{\nabla v}$ ($W$ divergence-free and zero on $\Sigma$); then the proof of Theorem~\ref{lk:low}, since Lemma~\ref{lk:defect} holds for $U_O$. (iii) Lemma~\ref{lk:lin} with (L$_D$) in place of Korn's inequality, plus the contraction of Corollary~\ref{ns:cor}(iii). (iv) Lemma~\ref{dr:id} with $\omega_h$ including $c(u_h;u_h,w)$ (dimension-free; $|\ut|\le C\hG^2$ on $\Oh\triangle\Omega$), and the reciprocity of Lemma~\ref{dr:recip} with $((u\cdot\nabla)U,Z)=-((u\cdot\nabla)Z,U)$, again in weak form. (v) The proof of Theorem~\ref{dr:cns} with (L$^\dagger$) for the auxiliary problem; the dual pair is smooth near $\Gamma$ by local regularity, and Lemma~\ref{dr:slip} acquires the terms $K(e_u,w_\phi)$ and $c(e_u;e_u,w_\phi)$, bounded as in two dimensions \cite{PadhiLong2D}.
\end{proof}
The quantitative threshold $h_0$ still needs Stokes $H^2$-regularity on $\Omega$ (Remark~\ref{ns:rem}), which is not used anywhere in Theorem~\ref{ns:oseen}.

\section{Computations}\label{sec:numerics}

All results of this section are numerical. They are a first three-dimensional implementation, run on a shared machine with 7\,GB of memory and two cores; the finer levels needed for asymptotic rates are listed in Appendix~\ref{app:repro} with the exact commands.

\subsection{Solver and meshes}
The solver (\texttt{code3d/svn3d.py}) uses continuous $P_k$ velocities, $k=3,4$, and discontinuous full-$P_{k-1}$ pressures on every sub-tetrahedron of an Alfeld refinement. The forms are those of Section~\ref{sec:setting}, with $h$ the largest tetrahedron diameter and $n$ pointing into the obstacle on $\Gh$; $\GS$ omits and $\CNS$ includes the mean-free traction. The pressure is eliminated by the iterated-penalty method (penalty $10^4$; the velocity matrix is factorised once with MKL PARDISO, and the $\CNS$ traction enters only the right-hand side), stopped when $\norm{\div u_h}<10^{-13}$ and the momentum residual is below $10^{-11}$, or at round-off stagnation. A nonsymmetric LU variant gives the same $u_h$ to $10^{-12}$.

\emph{Sphere meshes.} An equiangular $N\times N$ grid on each face of the box $R=(-2,2)^3$ is projected radially onto the unit sphere; the projected points are the vertices of a convex inscribed polyhedron $P_h$ with $12N^2$ faces. Radial layers at $r_j=|b|^{j/m}$, $m=N/2$, between $\partial P_h$ and $\partial R$ give prisms, which are split into three tetrahedra by a minimum-index rule and then Alfeld-refined. Every boundary face is a whole face of $P_h$ or lies on $\partial R$. Shape regularity is uniform ($\min\text{vol}/\diam^3=0.0074$--$0.0079$ for $N=2,\dots,24$), but $h/\hG=3.26,4.14,4.79,5.24$ for $N=2,4,6,8$: the family is not self-similar at small $N$, which contaminates rates measured in $\hG$.

\emph{Test problems.} (box) the cube $(-1,1)^3$ with a smooth manufactured solution and fully consistent Nitsche conditions on all faces; (P) $u=0$, $f=\nabla p$ with the wall pressure of the next test, so that $\ut=0$, $\Gc\equiv0$ and the $\GS$ output is exactly the traction response; (sph) Stokes flow past the fixed unit sphere with uniform stream $U=(1,0.4,-0.3)$, with nonconstant wall pressure; (rot) $u=(1-r^{-3})\,\omega\times x$, $p=0$, so that $u=0$ and $p$ is constant on $\Gamma$ ($G=0$).

\subsection{Verification}
\emph{Patch test.} A $P_k$ solution is reproduced to $10^{-9}$--$10^{-11}$ in $H^1$ for $k=3,4$, with Nitsche and with strong conditions. \emph{Orientation.} On the $N=2$ sphere mesh the computed normal points into $P_h$ on $\Gh$ and outward on the box (checked independently).

\begin{center}\small
\begin{tabular}{rrrcccccc}
\toprule
$k$ & $N$ & vel.\ dofs & $h$ & $H^1$ error & rate & $L^2$ error & rate & $\norm{\div u_h}$\\
\midrule
3&2&3\,189&1.732&$3.015\cdot10^{-1}$&--&$2.439\cdot10^{-2}$&--&$1.3\cdot10^{-13}$\\
3&3&10\,290&1.155&$9.780\cdot10^{-2}$&2.78&$5.265\cdot10^{-3}$&3.78&$1.7\cdot10^{-13}$\\
3&4&23\,871&0.866&$4.298\cdot10^{-2}$&2.86&$1.677\cdot10^{-3}$&3.98&$2.3\cdot10^{-13}$\\
3&5&46\,038&0.693&$2.252\cdot10^{-2}$&2.90&$6.812\cdot10^{-4}$&4.04&$3.0\cdot10^{-13}$\\
3&6&78\,897&0.577&$1.322\cdot10^{-2}$&2.92&$3.250\cdot10^{-4}$&4.06&$3.9\cdot10^{-13}$\\
4&2&7\,083&1.732&$3.820\cdot10^{-2}$&--&$3.266\cdot10^{-3}$&--&$6.7\cdot10^{-13}$\\
4&3&23\,115&1.155&$8.207\cdot10^{-3}$&3.79&$4.653\cdot10^{-4}$&4.81&$1.0\cdot10^{-12}$\\
4&4&53\,907&0.866&$2.649\cdot10^{-3}$&3.93&$1.138\cdot10^{-4}$&4.90&$1.5\cdot10^{-12}$\\
\bottomrule
\end{tabular}
\end{center}
\noindent\textbf{Table 1.} Box test, $\mu=400$. The rates approach $k$ in $H^1$ and $k+1$ in $L^2$, and $u_h$ is divergence-free to round-off.

\emph{Stability.} For the pair $(\VO,\Pih\cap L^2_0)$ with the vector Laplacian, a dense eigenvalue computation finds exactly one zero eigenvalue (the constants) and no spurious modes in every case; the discrete inf-sup constant is $0.142$, $0.149$, $0.150$ on the box for $N=1,2,3$, and $0.069$ on the sphere mesh $N=2$. Uniformity in $h$ is therefore seen only on the box.

\emph{Coercivity.} The threshold $\mu_Z^\ast$ of $N_h$ on $\Zh$ (from the inertia of the penalised matrix, independent of the penalty $10^5$--$10^7$, validated against a dense eigenvalue computation) is $76.3$ ($N=2$) and $148.0$ ($N=4$) for $k=3$, and $193.6$ ($N=2$) for $k=4$. It grows with $h/\hG$, as $\mu_0\propto\rho$ predicts (Section~\ref{sec:penalty}). A local sufficient threshold on $\Vh$, computed cell by cell, is $216$, $341$, $421$, $476$ for $N=2,4,6,8$ ($k=3$). All sphere runs below use $\mu\ge200>\mu_Z^\ast$.

\subsection{The leak law}
\begin{center}\small
\begin{tabular}{rrcccccc}
\toprule
$\mu$ & $N$ & $\norm{\nabla(\ut-u_h)}$ & $\dfrac{\norm{\nabla(\ut-u_h)}}{(h/\mu)G}$ & slip & leak & energy & $\norm{\div u_h}$\\
\midrule
200&2&$1.34\cdot10^{-1}$&2.34&0.913&0.891&1.001&$1.0\cdot10^{-15}$\\
200&4&$9.21\cdot10^{-2}$&2.41&0.949&0.943&1.012&$1.4\cdot10^{-15}$\\
400&2&$7.26\cdot10^{-2}$&2.54&0.941&0.923&0.992&$8.7\cdot10^{-16}$\\
400&4&$4.78\cdot10^{-2}$&2.50&0.970&0.965&1.002&$7.5\cdot10^{-16}$\\
1000&2&$3.30\cdot10^{-2}$&2.90&0.963&0.949&0.991&$2.3\cdot10^{-16}$\\
1000&4&$2.03\cdot10^{-2}$&2.66&0.984&0.981&1.001&$3.0\cdot10^{-16}$\\
\bottomrule
\end{tabular}
\end{center}
\noindent\textbf{Table 2.} Test P, $\GS$, $k=3$. slip $:=\norm{\ut-u_h}_{\Gh}/((h/\mu)G)$, leak $:=\ip{u_h\cdot n}{p-\pbar}_{\Gh}/((h/\mu)G^2)$, energy $:=\tnorm{\ut-u_h}(\mu/h)^{1/2}/G$, with $G$ computed on $\Gh$. $N=2$ and $N=4$ have $9\,096$ and $70\,026$ velocity unknowns.

The $\GS$ error is of order $h/\mu$ (the $H^1$ ratio is roughly constant, near $2.5$), and the slip, leak and energy constants approach $1$, with $1-\text{slip}$ decreasing with $h/\mu$: this is Corollary~\ref{er:Bp} in three dimensions. The $H^1$ ratio is not universal (in two dimensions it converges to a domain-dependent constant \cite{PadhiLong2D}). $\CNS$ has no leak: at $\mu=400$ its error on test P is $2.7\cdot10^{-3}$ ($N=2$) and $2.6\cdot10^{-4}$ ($N=4$), at the level of the pressure projection. At $\mu=200$ the $N=4$ $\CNS$ solve returns an error of order $1$: $\mu=200$ is below the well-posedness threshold $\gamma_2$ of Theorem~\ref{er:D} on that mesh, although it exceeds $\mu_Z^\ast$. We report $\CNS$ only for $\mu\ge400$.

\subsection{Sphere tests}
\begin{center}\small
\begin{tabular}{llrcccccc}
\toprule
test & $N$ & $\hG$ & $\GS$ $H^1$ & $\CNS$ $H^1$ & interp.\ $H^1$ & $\GS$ $\norm{e}_{\Gh}$ & $D$: $H^1$ & $D$: slip / leak\\
\midrule
sph&2&0.919&1.269&1.265&1.633&0.445&0.138&1.254 / 0.978\\
sph&4&0.524&0.338&0.325&0.192&0.110&0.064&0.978 / 0.925\\
rot&2&0.919&2.209&2.225&2.390&0.867&--&--\\
rot&4&0.524&0.541&0.549&0.251&0.223&--&--\\
\bottomrule
\end{tabular}
\end{center}
\noindent\textbf{Table 3.} Tests sph and rot, $k=3$, $\mu=400$. ``interp.'' is the $H^1$ error of the nodal interpolant of $\ut$; $D:=u_h^{\GS}-u_h^{\CNS}$, with slip and leak ratios defined as in Table~2. $\norm{\div u_h}$ is $10^{-13}$--$10^{-10}$ in all runs.

\subsection{Interpretation}
Confirmed at these levels: pointwise incompressibility; the optimal rates $h^k$ and $h^{k+1}$ on the box for $k=3,4$; a stable pair without spurious pressure modes; the growth of the coercivity threshold with $h/\hG$; the leak law of $\GS$ with constant tending to $1$ in slip and energy (test P at three penalties, and $D$ in test sph, with slip $0.98$ and leak $0.93$ at $N=4$); and the removal of the leak by $\CNS$ (the leak correlation of $\CNS$ in test sph is $-0.38$ at $N=2$ and $-0.06$ at $N=4$, against $0.60$ and $0.87$ for $\GS$).

\emph{Not confirmed: the rates.} Only two sphere levels fit in memory, and they are pre-asymptotic. The observed $\hG$-rates of $2.4$--$2.5$ for both methods in tests sph and rot are dominated by the $P_3$ approximation error (the interpolation error drops from $1.63$ to $0.19$, and at $N=2$ it is larger than the finite element error) and by the $O(\hG^2)$ trace mismatch of $\ut$ on $\Gh$. At $N=4$ the finite element error is already $1.7$--$2.2$ times the interpolation error, so the geometric $\hG^{3/2}$ part is taking over; measuring its rate needs $N\ge8$. In test sph the $\GS$ leak ($D$, $0.064$ in $H^1$ at $N=4$) is still small compared with the geometric error ($0.33$), so the rate $1$ of $\GS$ and the rate $3/2$ of $\CNS$ in total $H^1$ error will separate only at much finer levels. The drift of $h/\hG$ also affects rates in $\hG$. We therefore make no claim about asymptotic rates on the sphere from these computations; the commands for $N=6,\dots,12$ (up to $1.85$ million unknowns) are in Appendix~\ref{app:repro}.

\section{What is not proved}\label{sec:open}

\begin{enumerate}
\item \emph{Unsplit meshes.} (IS3) on general unsplit tetrahedral meshes is, to our knowledge, open for the Scott--Vogelius pair in three dimensions (Remark~\ref{st:rem:status}). All results that need existence of $u_h$ are therefore unconditional only on Alfeld refinements; on unsplit meshes they hold under (IS3). The lower bound of Section~\ref{sec:sharp} for $4\le k\le8$ on unsplit meshes is not proved (Proposition~\ref{sh:bubble} needs $k\ge9$).
\item \emph{Worsey--Farin splits.} (IS3) as formulated, with the full space of discontinuous $P_{k-1}$ pressures, is false when the wall faces are subdivided (Proposition~\ref{st:wf}); this is known \cite{GuzmanLischkeNeilan2022,FabienGuzmanNeilanZytoon2022}. So none of our upper bounds, and not the leak analysis of Section~\ref{sec:leak}, is proved on Worsey--Farin meshes. That would need: a pressure space with the singular-edge conditions of \cite{FabienGuzmanNeilanZytoon2022}; a uniform-in-$h$ inf-sup constant for that pair on the perturbed domains $\Oh$, as Lemma~\ref{st:IS3} provides for Alfeld; and a lift lemma onto the constrained trace space. The lower bound holds for every solution that exists, for $k\ge4$ on barycentric Worsey--Farin splits (Proposition~\ref{sh:wf}) and for $k\ge9$ in general. For $k=3$ there is no single-macro witness, and the trace analysis of Section~\ref{sh:sec:k3} has not been redone for split faces; the $k=3$ edge-star ranks are only numerical (Remark~\ref{sh:rem:wf}).
\item \emph{Explicit constants.} On Alfeld refinements with $k\ge4$ the witness constants are explicit (Section~\ref{sec:constants}). Not explicit: the constant $\kappa_0$ of Theorem~\ref{k3:witness} (degree $3$, general surfaces), which comes from compactness (the sphere criterion of Theorem~\ref{k3:sphere} is explicit up to $C_E(\theta)$); the constants $\kappa_0,Y_0$ of Theorem~\ref{pen:cont}; and the (IS3) constant, because $C_\Pi$ and $C_B(\Omega)$ are not explicit and $\hat C_{\rm loc}$ is only computed in floating point. The single-cell penalty bound $\lambda_T\ge20$ is certified on one explicit box of shapes (with the floating-point caveat of Section~\ref{sec:constants}), not over larger classes.
\item \emph{Penalty necessity in general.} Necessity of $\mu\gtrsim\rho$ without local quasi-uniformity (LQ) or a certified cell at a smallest face is open, as in two dimensions.
\item \emph{Leak and drag.} The leak lower bound of Theorem~\ref{lk:low} is proved only on Alfeld refinements with $k\ge4$; for $k=3$ there is no single-macro leak witness, and multi-macro witnesses would be needed. For general data at a fixed penalty the rate-$\frac12$ lower bound is open (Remark~\ref{lk:rem:gen}). The rate $2$ of the $\CNS$ drag (Theorem~\ref{dr:cns}) is proved modulo (E$_Z$), a regularity statement for the dual Stokes flow at $\Sigma$, and in the Oseen case modulo (E$_Z^\dagger$); its sharpness is not proved. The Gjerde--Scott volume functional and pointwise tractions are not treated.
\item \emph{Regularity at the outer boundary.} The smoothness of $(u,p)$ on $\overline\Omega$ is a hypothesis on the solution: the edges and corners of the box limit generic regularity (Remark~\ref{set:rem:corner}). Stokes $H^2$-regularity on $\Omega$, needed for an explicit threshold $h_0$ in Corollary~\ref{ns:cor}(ii) (Remark~\ref{ns:rem}), is reduced to a convex-polyhedron theorem attributed to \cite{Dauge1989} whose statement we have not checked; (E$_Z$) rests on the same kind of input.
\item \emph{Citations checked indirectly.} The wording of \cite[Theorem~3.1]{GuzmanNeilan2018} was read through a text-extraction tool; \cite[Proposition~6.5]{Neilan2015} only through \cite{FarrellMitchellScott2024}; the pressure spaces of \cite{GuzmanLischkeNeilan2022} only through its abstract and \cite{FabienGuzmanNeilanZytoon2022}; the Bogovski\u{\i} regularity used in Lemma~\ref{er:w} and the boundary Schauder theory for the leak flow in Section~\ref{sec:leak} are cited, not re-proved.
\item \emph{Numerics.} The sphere computations are pre-asymptotic; the $\hG^{3/2}$ rates of Theorems~\ref{er:D} and~\ref{sh:k3main}, the rate $\hG^{1/2}$ of the leak limit, and the separation of the $\GS$ and $\CNS$ rates have not been observed in three dimensions (Section~\ref{sec:numerics}).
\end{enumerate}


\appendix
\section{Degree-three witnesses on edge stars}\label{app:k3}\label{sh:sec:k3}

This appendix proves Theorem~\ref{k3:witness}, the local condition (M3$^3$) for $k=3$ on Alfeld refinements, which gives the lower bound of Theorem~\ref{sh:k3main}. By Proposition~\ref{sh:k3single} a single macro-element is not enough, so the witnesses live on the star of a macro edge $[z,q]$ joining a wall vertex $z$ to an interior vertex $q$. The argument has three steps. (1) We describe exactly which traces on the wall faces of the star are attained by discrete divergence-free fields (the Ring Trace Theorem~\ref{k3:trace}); the description is in terms of a few vectors attached to the interior faces of the star. (2) For a ``vertex move'', in which these vectors are all multiples of one vector $\delta$, the edge moments have a closed form (Theorem~\ref{k3:vertex}); on the sphere this already gives an explicit witness (Theorem~\ref{k3:sphere}). (3) On a general surface the weights of the vertex move can cancel at saddle points, and two further families of moves are needed; a nondegeneracy theorem for flat stars (Theorem~\ref{k3:nondeg}) and a compactness argument give Theorem~\ref{k3:witness}.

\subsection{Edge stars and their traces}
Throughout this appendix $\Th$ is an Alfeld refinement with minimal angle $\theta$ of the macro mesh, and $k=3$ (larger $k$ follow, since the witnesses are also piecewise $P_k$ and the moments do not depend on $k$).

\paragraph{Edge stars.} Let $z$ be a vertex of $\Gh$ and $q$ a macro vertex with $[z,q]$ a macro edge and $q\notin\Gh$. For the apex $q_F$ of $T_F$ this holds for each vertex $z$ of $F$: $q_F$ has distance at least $c\,\diam F$ from the plane of $F$, on the side of $\Oh$, while $\Gh$ lies within $C(\diam F)^2$ of that plane near $F$ (as in the proof of Theorem~\ref{geom:mesh}). The relative interior of $[z,q]$ lies in $\Oh$, so the \emph{edge star} $R=R(z,q)$, the union of the macro-tetrahedra containing $[z,q]$, has a link that is a cycle $x_0,\dots,x_{m-1}$ with $3\le m\le m_{\max}(\theta)$:
\[
T_j=\conv(q,z,x_j,x_{j+1}),\qquad f_j=\conv(q,z,x_j)\ \text{(interior faces)},\qquad F_j:=\conv(z,x_j,x_{j+1}).
\]
The faces of the $T_j$ on $\Gh$ are among the $F_j$; let $I:=\{j:F_j\in\FG\}$. Let $Y(R)$ be the continuous piecewise-$P_3$ divergence-free fields on the Alfeld refinement of $R$ vanishing on $\partial R$, and $Y^1(R):=\{v\in Y(R):\int_{F_i}\partial_\nu v=0,\ i\in I\}$. Extension by zero maps $Y^1(R)$ into $Y_h^1$.

\paragraph{Traces.} Fix $i\in I$, write $F=F_i=\conv(z,a,b)$, $T=T_i$ with barycentre $B$, and $H:=\dist(q,\text{plane of }F)$. The Alfeld sub-tetrahedra of $T$ at $z$ are $s_1=\conv(B,z,a,b)\supset F$, $s_2=\conv(B,q,z,b)$ and $s_3=\conv(B,q,z,a)$. Since $v|_{s_1}$ vanishes on $F$, $\nabla v|_{s_1}=\gh\otimes\nabla\lambda_B^{s_1}$ on $F$ with $\gh(x):=\partial_{B-x}v(x)\in P_2(F;\R^3)$, and $\partial_\nu v=\frac4H\gh$. In this appendix we work with
\[
\hat\M_F(v):=\frac4H\int_FQ_F\,\gh\,dA=-\M_F(v),
\]
which is harmless because only the sign of the witnesses changes. With $m_{xy}$ the midpoint of $[x,y]$ and $n_S$ a normal of the plane of $S$, define
\begin{align*}
U_F&:=\{\gh\in P_2(F;\R^3):\gh\cdot n_F\equiv0,\ \gh(a)=\gh(b)=0,\ \gh(m_{ab})\cdot n_{qab}=0\}\qquad(\dim7),\\
S_F&:=\{\gh\in U_F:\gh(z)=0,\ \gh(m_{za})\cdot n_{f_a}=0,\ \gh(m_{zb})\cdot n_{f_b}=0\}\qquad(\dim3),
\end{align*}
with $f_a=\conv(q,z,a)$, $f_b=\conv(q,z,b)$. The map $\pi_F:U_F\to T_F\times\R\times\R$, $\gh\mapsto(\gh(z),\gh(m_{za})\cdot n_{f_a},\gh(m_{zb})\cdot n_{f_b})$, is onto with kernel $S_F$ (here $T_F$ also denotes the tangent plane of $F$).

\begin{lemma}[Necessary relations]\label{k3:nec}
Let $v\in Y(R)$ and $\delta:=\partial_{q-z}v(z)$, which is common to all $T_j$. Then:
\begin{enumerate}[label=(\alph*)]
\item $\gh_i\in U_{F_i}$ for every $i\in I$.
\item On $f_j$, in its barycentrics, $v|_{f_j}=\lambda_q\lambda_z(p_q\lambda_q+\delta\lambda_z+p_j\lambda_x)$ with $p_q\in\R^3$ common to all $j$. The fluxes $A_j\cdot(\frac{\delta+p_q}{30}+\frac{p_j}{60})$, $A_j:=\frac12(q-z)\times(x_j-z)$, are equal for all $j$. Along $[z,x_j]$, $h_j:=\partial_{q-x}(v|_{f_j})$ satisfies $h_j(m_{zx_j})=\frac14(\delta+p_j)$.
\item $\delta_i:=\partial_{B-z}v|_{T_i}(z)$ satisfies $\delta_i\cdot\nabla\lambda_B^{s_1}=0$ and $\delta\cdot\nabla\lambda_q^s+\delta_i\cdot\nabla\lambda_B^s=0$ for $s=s_2,s_3$, and $\gh_i(z)=\delta_i$.
\item $\gh_i\cdot\nabla\lambda_B^{s_3}+h_i\cdot\nabla\lambda_q^{s_3}=0$ on $[z,a]$, and likewise on $[z,b]$ with $s_2$ and $h_{i+1}$.
\end{enumerate}
\end{lemma}

\begin{proof}
(a) At $a$, the three edges of $T$ at $a$ lie in $\partial R$, so in each sub-tetrahedron at $a$, $\nabla v(a)=\delta'\otimes\nabla\lambda_B$ with a common $\delta'$; $\div v=0$ makes $\delta'$ orthogonal to three independent normals, so $\gh(a)=0$, and likewise $\gh(b)=0$. On $[a,b]$, in $\conv(B,q,a,b)$, $\nabla v=\gh\otimes\nabla\lambda_B$ and $\div v=0$ gives $\gh\cdot n_{qab}=0$; tangency is $\div v|_{s_1}=\gh\cdot\nabla\lambda_B^{s_1}=0$. (b) $v|_{f_j}\in P_3$ vanishes on $[z,x_j]$ and $[q,x_j]$, and its restriction to $[z,q]$ is shared. With $\int_T\lambda_q^2\lambda_z=|T|/30$ and $\int_T\lambda_q\lambda_z\lambda_x=|T|/60$, and $\int_{\partial T_j}v\cdot n=0$ with $\partial T_j\setminus\partial R=f_j\cup f_{j+1}$, the fluxes agree. (c), (d): at $z$ in $s_2$, $v$ vanishes on $[z,b]$, so $\nabla v(z)=\delta\otimes\nabla\lambda_q+\delta_i\otimes\nabla\lambda_B$ with zero trace; similarly in $s_3$, $s_1$ and along $[z,a]$.
\end{proof}

\begin{lemma}[Single macro-element]\label{k3:single}
For one Alfeld macro-tetrahedron $T$ with $v=0$ on $\partial T$, the trace space $\{\gh\}$ is exactly $S_F$, and $\gh\mapsto\int_F\gh$ has rank $2$ on $S_F$. \status{computer-assisted}
\end{lemma}

\begin{proof}
$\subseteq$: Lemma~\ref{k3:nec} with $\delta=0$ and $h\equiv0$; then $\delta_i=0$. $\supseteq$: the trace rank on $Y(T)$ is at least $3$ (a rank computation modulo a prime, which is a lower bound for the rank over $\mathbb Q$; the rank $39$ of the divergence is its a priori maximum). The Piola map multiplies $\gh$ by a constant matrix and preserves $S_F$, so the statement is shape independent. The mean statement is exact on the reference triangle.
\end{proof}
\cert{notes/v6/td\_k3/trace1.py, ref\_K.py}{trace1.log, ref\_K.log}

\begin{lemma}[Realisability]\label{k3:suff}
Given $\delta,p_q,p_0,\dots,p_{m-1}\in\R^3$ with equal fluxes as in Lemma~\ref{k3:nec}(b), and $s_i\in S_{F_i}$ ($i\in I$), there is $v\in Y(R)$ with these interior-face data such that, for each $i\in I$, $\gh_i$ differs from a fixed element of $U_{F_i}$ with $\pi_{F_i}$-value determined by (c), (d) by exactly $s_i$.
\end{lemma}

\begin{proof}
Let $v_0$ be the continuous piecewise-$P_3$ field whose Bernstein coefficients on the domain points of the faces $f_j$ are those of the prescribed $v|_{f_j}$ ($\delta/3$ at $\frac{q+2z}3$, $p_q/3$ at $\frac{2q+z}3$, $p_j/6$ at $\frac{q+z+x_j}3$), and zero elsewhere. Then $v_0=0$ on $\partial R$ and $\int_{T_j}\div v_0=0$ by the equal fluxes. On each Alfeld macro-element, $\div$ maps the piecewise-$P_3$ fields vanishing on $\partial T_j$ onto the mean-free piecewise $P_2$ (Fact~\ref{st:gn}; exact rank $39$ on the reference element), which gives $w_j$ with $\div w_j=-\div v_0$ on $T_j$ and $w_j=0$ on $\partial T_j$. $v_1:=v_0+\sum_jw_j\in Y(R)$ has the prescribed data, and adding single-macro fields (Lemma~\ref{k3:single}) adjusts the $S_{F_i}$-components independently.
\end{proof}

\begin{theorem}[Ring Trace Theorem]\label{k3:trace}
The joint trace space $\{(\gh_i)_{i\in I}:v\in Y(R)\}$ consists exactly of the tuples with $\gh_i\in U_{F_i}$ and $\pi_{F_i}(\gh_i)$ given by Lemma~\ref{k3:nec}(c), (d) for some data $(\delta,p_q,(p_j))$ with all fluxes equal.
\end{theorem}

\begin{proof}
Lemma~\ref{k3:nec} and Lemma~\ref{k3:suff}.
\end{proof}
An exact implementation of this reduced model agrees with the full finite element computation in dimension and moment rank on eight configurations, including the rank-$0$ octahedral fan below (\texttt{notes/v6/td\_k3/validate\_reduced.log}).

For $i\in I$ let $n_i$ be a unit normal of $F_i$, $P_i:=I-n_i\otimes n_i$, $H_i:=|(q-z)\cdot n_i|$, and
\[
r_i:=(a_i-z)+(b_i-z)+P_i(q-z),\qquad G_i\delta:=P_i\delta-\frac{n_i\cdot\delta}{(q-z)\cdot n_i}\,r_i,\qquad c_i:=\frac{|F_i|\,w^{F_i}_{ab}}{H_i},
\]
where $w^{F_i}_{ab}$ is the weight of the edge of $F_i$ opposite $z$.

\begin{lemma}[Jet map]\label{k3:jet}
In Lemma~\ref{k3:nec}(c), $\gh_i(z)=\delta_i=\frac14G_i\delta$.
\end{lemma}

\begin{proof}
The solution is unique, so it suffices to check $\delta=q-z$ and $\delta\in T_{F_i}$. For $\delta=q-z$: $\nabla\lambda_q^s\cdot(q-z)=1$, and $y=-(B-z)+t(q-z)$ satisfies $\nabla\lambda_B^s\cdot y=-1$; tangency forces $t=\frac14$ (heights $H/4$ and $H$), so $\delta_i=-\frac14((a-z)+(b-z))=\frac14G_i(q-z)$. For $\delta\in T_{F_i}$: $\varphi:=\frac14\lambda_B^s+\lambda_q^s$ is affine, vanishes at $z$ and at $b$ (resp.\ $a$), and equals $\frac14$ at $B$ and $1$ at $q$, so $\varphi=\dist_\pm(\cdot,F)/H$ and $\nabla\varphi\perp T_F$; hence $\delta_i=\frac14\delta$.
\end{proof}

\begin{lemma}[Face moment of the vertex representative]\label{k3:K}
Let $e\in T_F$. The set $\{\gh\in U_F:\pi_F(\gh)=(e,0,0),\ \int_F\gh=0\}$ is a line parallel to $S_F\cap\{\int\gh=0\}$, and it contains an element $r_F(e)$, linear in $e$, with $\int_FQ_F\,r_F(e)\,dA=-\frac{|F|\,w^F_{ab}}{360}\,e$. \status{computer-assisted, exact}
\end{lemma}

\begin{proof}
On the reference triangle (exact): a representative has moment $Ke$ with $K=\frac1{720}\begin{psmallmatrix}-w_{ab}&2(w_{ab}-w_{za})\\0&-3w_{ab}+2w_{zb}\end{psmallmatrix}$ in the basis $(a-z,b-z)$, and the line direction has moment $D=\frac1{720}(w_{za}-w_{ab},\,w_{ab}-w_{zb})$. Since $K\binom01+2D=-\frac{w_{ab}}{720}\binom01$, $K\equiv-\frac{w_{ab}}{720}I$ modulo $D$; fix $r_F$ accordingly. Transport: $\int_FQ\gh\,dA=2|F|A\int_{\rm ref}Q\hat e$ for $\gh=A\hat e$.
\end{proof}
\cert{notes/v6/td\_k3/ref\_K.py; independent re-derivation notes/v6/referee8/ref8\_lemmaK.py}{ref\_K.log, ref8\_lemmaK.log}

\begin{theorem}[Vertex witness]\label{k3:vertex}
For every edge star $R(z,q)$ with $I\ne\emptyset$, any geometry and any real weights, and every $\delta\in\R^3$, there is $v_\delta\in Y^1(R)$ with
\[
\sum_{i\in I}\hat\M_{F_i}(v_\delta)=-\frac1{360}\sum_{i\in I}c_i\,G_i\,\delta,\qquad\norm{\nabla v_\delta}_R\le C_E(\theta)\,|\delta|\,(\diam R)^{1/2}.
\]
\end{theorem}

\begin{proof}
Take $p_q=-\delta/2$ and $p_j=-\delta$ for all $j$. All fluxes are $A_j\cdot(\frac\delta{60}-\frac\delta{60})=0$ and $h_j(m_{zx_j})=0$, so both $\nu$-components of $\pi_{F_i}$ vanish (Lemma~\ref{k3:nec}(d)). By Lemmas~\ref{k3:suff} and~\ref{k3:jet} there is $v\in Y(R)$ with $\pi_{F_i}(\gh_i)=(\frac14G_i\delta,0,0)$ for all $i\in I$. Subtracting single-macro fields (Lemma~\ref{k3:single}) makes $\gh_i=r_{F_i}(\frac14G_i\delta)$, which has mean zero, and Lemma~\ref{k3:K} gives the moment. The energy bound: the Bernstein data are $O(|\delta|)$, and the Alfeld right inverse and the single-macro right inverses are fixed reference maps transported affinely, so their norms depend only on $\theta$ ($m\le m_{\max}(\theta)$); scaling gives $(\diam R)^{1/2}$.
\end{proof}
\cert{notes/v6/td\_k3/vertex\_formula\_general.py, fe\_vertex\_check.py; notes/v6/referee8/ref8\_vertex\_fe.py}{vertex\_formula\_general.log, fe\_vertex\_check.log, ref8\_vertex\_fe\_eqw.log}
In the exact reduced model the formula holds in all $30$ test cases. Where the weights are generic, the test ``modulo the span of the single-macro directions'' can be vacuous; the non-vacuous checks are the equal-weight cases and an independent full finite element computation on three non-flat stars inscribed in the sphere with distinct per-face equal weights, where the finite element moment agrees with the closed form to $10^{-15}$.

\paragraph{Flat stars and the octahedral degeneracy.} If all $F_i$ lie in one plane $T$, then $G_i|_T=I$ and the tangential moment of $v_\delta$ is $-\frac1{360}(\sum_ic_i)\delta$, a multiple of the identity. For the octahedral fan ($z=e_3$, $q=2e_3$, link $\pm e_1,\pm e_2$, all four faces in $\Gh$, equal weights), $\sum_ic_iP_TG_i|_T=0$ exactly and the full finite element rank is $0$; its faces are tilted by $\arccos(1/\sqrt3)$ against $T$. On shape-regular inscribed meshes with small $\hG$ such tilts do not occur.

\subsection{The sphere: an explicit criterion}
For a vertex $z$ of $\Gh$ with unit normal $n_z$ of $\Gamma$ and an edge star $R=R(z,q)$ put
\[
\tau_z:=\max_{i\in I}\angle(n_{F_i},n_z),\qquad\rho_z:=\max_{i\in I}|r_i|/H_i .
\]
On the unit sphere $\angle(n_F,z)=\arcsin r_F$ for every vertex $z$ of $F$, so $\tau_z\le C(\theta)\hG$, while $\rho_z\le C(\theta)$.

\begin{theorem}[(M3$^3$) for $k=3$ on the sphere, explicit]\label{k3:sphere}
In the model case, let $\Th$ be an Alfeld refinement and $k\ge3$. Suppose every $F\in\FG$ has a vertex $z$ with $\tau_z(\tau_z+\rho_z)\le\frac12$ for $R(z,q_F)$; this holds for $\hG\le h_\ast(\theta)$. Then (M3$^3$) holds for $\hG\le h_\ast(\theta)$ with $\omega_F=R(z,q_F)$, $K_{\rm ov},C_s\le C(\theta)$ and $\kappa_0=c(\theta)/C_E(\theta)$.
\end{theorem}

\begin{proof}
For unit $t'\in n_z^\perp$: $|n_i\cdot t'|\le\sin\tau_z$ and $|P_it'|^2\ge1-\sin^2\tau_z$, so
\[
t'\cdot G_it'=|P_it'|^2-\frac{(n_i\cdot t')(r_i\cdot t')}{(q-z)\cdot n_i}\ge1-\tau_z(\tau_z+\rho_z)\ge\tfrac12 .
\]
Given a unit $t\in T_{y_F}\Gamma$, let $t'$ be its normalised projection onto $n_z^\perp$; $|t-t'|\le C\hG$, so $t\cdot G_it'\ge\frac12-C\hG(1+\rho_z)\ge\frac14$ for $\hG\le h_\ast(\theta)$. On the sphere all $c_i$ have the same sign ($w_{ab}=-\ell_{ab}^2$); take $\delta:=\pm t'$ in Theorem~\ref{k3:vertex} with the sign that makes $t\cdot\sum_i\hat\M_{F_i}(v_\delta)=\frac1{360}\sum_i|c_i|\,t\cdot G_it'\ge\frac1{1440}|c_F|\ge c(\theta)(\diam F)^3$, using $\ell_{ab}\ge c\diam F$, $|F|\ge c(\diam F)^2$ and $H\le C\diam F$. Divide by $\norm{\nabla v_\delta}\le C_E(\diam F)^{1/2}$. A macro-tetrahedron lies in $R(z,q_F)$ only if $[z,q_F]$ is one of its six edges, so $K_{\rm ov}\le C(\theta)$.
\end{proof}
The only non-explicit constant is $C_E(\theta)$, a finite combination of right-inverse norms of fixed reference maps. On icosphere meshes with a radial prism layer (Alfeld-refined), the criterion holds from the second or third level on, and the matrix $\sum_ic_iP_TG_i|_T$ is nonsingular at every level (\texttt{notes/v6/td\_k3/icosphere\_check.log}). \status{numerical}

\subsection{General surfaces}
On a general surface the weights $w_{ab}$ of the vertex move can cancel, $\sum_ic_i=0$, at saddle points. Two further families of moves in the Ring Trace Theorem repair this. For a face $F=F_i=\conv(z,a,b)$ write $e_a=a-z$, $e_b=b-z$, $P_a=\delta+p_a$, $P_b=\delta+p_b$ (the spoke data), $u_a=\nabla\lambda_q^{s_3}$, $u_b=\nabla\lambda_q^{s_2}$, $w=(w_{za},w_{zb},w_{ab})$ its weights, and
\[
D_F:=(w_{za}-w_{ab})\,e_a+(w_{ab}-w_{zb})\,e_b,
\]
the single-macro direction of Proposition~\ref{sh:k3single}.

\begin{lemma}[Per-face moment]\label{k3:face}
For every $v\in Y^1(R)$ with interior-face data $(\delta,p_q,(p_j))$ and every $i\in I$,
\[
\hat\M_{F_i}(v)\equiv\frac{8|F_i|}{H_i}\Bigl[-\frac{w_{ab}}{2880}G_i\delta-\frac{w_{za}-w_{zb}}{2880}\Bigl((P_a\cdot u_a)e_b-(P_b\cdot u_b)e_a\Bigr)\Bigr]\pmod{\R D_{F_i}},
\]
and every value of the right-hand side, plus any multiple of $D_{F_i}$ chosen independently for each $i$, is attained. \status{computer-assisted, exact}
\end{lemma}

\begin{proof}
On the reference triangle $z=0$, $a=e_1$, $b=e_2$, for $\gh\in U_F\cap\{\int\gh=0\}$ with $\gh(z)=e$, $\tilde\alpha:=\gh(m_{za})\cdot\nabla\mu_b$ and $\tilde\beta:=\gh(m_{zb})\cdot\nabla\mu_a$, an exact computation gives $\int_{\rm ref}Q\gh\equiv-\frac{w_{ab}}{720}e+\frac{w_{za}-w_{zb}}{180}(\tilde\alpha\hat e_2-\tilde\beta\hat e_1)$ modulo $D$. Under $\gh=A\gh_{\rm ref}$ the quantities $\tilde\alpha,\tilde\beta$ are affine invariants, so $\int_FQ\gh\equiv2|F|[-\frac{w_{ab}}{720}\gh(z)+\frac{w_{za}-w_{zb}}{180}(\tilde\alpha e_b-\tilde\beta e_a)]$ modulo $\R D_F$. By Lemma~\ref{k3:jet}, $\gh(z)=\frac14G_i\delta$; by Lemma~\ref{k3:nec}(d) with $h_i(m_{za})=\frac14P_a$, and since the tangential part of $\nabla\lambda_B^{s_3}$ equals $4\nabla\mu_b$, $\tilde\alpha=-\frac1{16}P_a\cdot u_a$, and likewise $\tilde\beta=-\frac1{16}P_b\cdot u_b$. Multiply by $4/H_i$. Realisability is Theorem~\ref{k3:trace}, and Lemma~\ref{k3:single} leaves exactly the $D_{F_i}$-freedom on each face.
\end{proof}
\cert{notes/v6/td\_k3b/ref\_nu.py, face\_formula.py}{ref\_nu.log, face\_formula.log (29/29 faces; rerun with a fresh seed in notes/v6/referee9)}
The check compares face by face modulo the single line $\R D_{F_i}$, so it is never vacuous.

A star is \emph{$\Gamma$-flat} if its $\Gamma$-faces lie in one plane $T$. Put $z=0$, $T=\{x_3=0\}$, $q=(q',H)$ with $H>0$; let $J$ be the rotation by $\frac\pi2$ in $T$, and orient so that $s_i\wedge s_{i+1}=2|F_i|>0$ for $i\in I$, where $s_j:=x_j$ are the spokes. Let $\hh$ be a symmetric form on $T$ with weights $w^F_{ij}=\hh(e_{ij},e_{ij})$; put $\omega_j:=\hh(s_j,s_j)$, $w_i:=w^{F_i}_{ab}$, $\hat\tau_i:=\omega_i-\omega_{i+1}$, write $P_j=(\pi_j,\nu_j)\in T\times\R$, and $\sigma_j:=Js_j\cdot(\pi_j-\nu_jq'/H)$.

\begin{lemma}[Flat reduction]\label{k3:flat}
On a $\Gamma$-flat star:
\begin{enumerate}[label=(\alph*)]
\item $P_a\cdot u_a=-\sigma_a/(2|F|)+\nu_a/H$ and $P_b\cdot u_b=\sigma_b/(2|F|)+\nu_b/H$, hence, modulo $\operatorname{span}\{D_{F_i}\}$,
\[
\sum_{i\in I}\hat\M_{F_i}\equiv\frac1{360H}\sum_{i\in I}\Bigl[-|F_i|w_iG_i\delta+\hat\tau_i\Bigl(\tfrac12(\sigma_is_{i+1}+\sigma_{i+1}s_i)-\tfrac{|F_i|}H(\nu_is_{i+1}-\nu_{i+1}s_i)\Bigr)\Bigr].
\]
\item The flux condition for spoke $j$ reads $A_j\cdot P_j=\frac H2\sigma_j$. The admissible data are exactly: $\nu_j\in\R$ free for each $\Gamma$-spoke, and $\sigma_j=\gamma-s_j\wedge\eta$ with $(\gamma,\eta)\in\R\times\R^2$ free; the data $\delta$, $(\nu_j)$, $(\gamma,\eta)$ and the $D$-coefficients are mutually independent.
\end{enumerate}
\end{lemma}

\begin{proof}
(a) In the flat case $u_a=(-\nabla\mu_b,(1+\nabla\mu_b\cdot q')/H)$, $u_b=(-\nabla\mu_a,(1+\nabla\mu_a\cdot q')/H)$, $\nabla\mu_b=Js_a/(2|F|)$ and $\nabla\mu_a=-Js_b/(2|F|)$; substitute into Lemma~\ref{k3:face}. (b) $A_j=\frac12q\times x_j=\frac12(HJs_j,q'\wedge s_j)$ gives $A_j\cdot P_j=\frac H2\sigma_j$. Lemma~\ref{k3:nec}(b) with $p_j=P_j-\delta$ gives $A_j\cdot P_j=60\Phi-A_j\cdot\psi$ with $\psi:=\delta+2p_q$ and $\Phi$ the common flux, both free; so $\sigma_j=\frac{120\Phi}H-s_j\wedge(\psi'-\psi_nq'/H)$. Given $\nu_j$ and $\sigma_j$, the component $Js_j\cdot\pi_j$ is fixed; the component along $s_j$ is invisible.
\end{proof}
\cert{notes/v6/td\_k3b/flat\_algebra.py}{flat\_algebra.log}

\begin{definition}
A $\Gamma$-face $F=\conv(z,a,b)$ is \emph{degenerate at $z$} if $w_{ab}=0$ and $w_{za}=w_{zb}$.
\end{definition}

\begin{lemma}\label{k3:degface}
$F$ is degenerate at $z$ iff $\hh=0$ or $\hh=\pm\nu\otimes\nu$ with $\nu\perp b-a$. In particular no face is degenerate at any vertex if $\det\hh\ne0$, and if $\hh\ne0$, $F$ is nondegenerate at the vertex opposite any edge $e$ with $\hh(e,e)\ne0$; such an edge exists.
\end{lemma}

\begin{proof}
With $e=b-a$, $\hh(e_b,e_b)=\hh(e_a,e_a)+2\hh(e_a,e)+\hh(e,e)$, so degeneracy means $\hh(e,e)=\hh(e_a,e)=0$. If $\det\hh\ne0$, the $\hh$-orthogonal complement of a null vector $e$ is $\R e$, forcing $e_a\parallel e$. If $\hh=\pm\nu\otimes\nu$ both conditions reduce to $\nu\cdot e=0$. A quadratic form vanishing on three pairwise non-parallel vectors is zero.
\end{proof}

\begin{theorem}[Star Nondegeneracy]\label{k3:nondeg}
Let $R$ be a $\Gamma$-flat edge star, $\hh\ne0$ a symmetric form on $T$ with weights $\hh(e,e)$, and assume the $\Gamma$-spokes are pairwise distinct as directions. If some $\Gamma$-face of $R$ is nondegenerate at $z$, then $v\mapsto P_T\sum_{i\in I}\hat\M_{F_i}(v)$ maps $Y^1(R)$ onto $T$.
\end{theorem}

\begin{proof}
Suppose $0\ne\lambda\in T$ annihilates the image; put $\ell_j:=\lambda\cdot s_j$, $\mu:=J\lambda$ and $\tau_i:=|F_i|\hat\tau_i$, and write $\mathbf 1_i=[i\in I]$. By Lemma~\ref{k3:flat} the families are independent, so $\lambda$ annihilates each:
\begin{align*}
&(\delta_T)\ \textstyle\sum_{i\in I}|F_i|w_i=0\quad(\delta\in T,\ G_i|_T=I);\qquad(D)\ \lambda\cdot D_{F_i}=0\ \ \forall i\in I;\\
&(\nu_j)\ \tau_{j-1}\ell_{j-1}\mathbf 1_{j-1}=\tau_j\ell_{j+1}\mathbf 1_j\ \text{ for every $\Gamma$-spoke $j$};\\
&(\sigma)\ \textstyle\sum_{i\in I}\hat\tau_i\bigl[\gamma(\ell_i+\ell_{i+1})+L(s_i)\ell_{i+1}+L(s_{i+1})\ell_i\bigr]=0\quad\forall\gamma\in\R,\ L\in T^\ast .
\end{align*}
\emph{Step 1: $\tau\equiv0$.} Let $\varpi_i:=\tau_i\ell_i\ell_{i+1}$. Multiplying $(\nu_j)$ by $\ell_j$ shows that $\varpi$ is constant on each run of consecutive $\Gamma$-faces and zero at the ends of a non-cyclic run. If $I$ is the whole cycle, $(\sigma)$ with $\gamma=0$, $L=\lambda\cdot$ gives $2\sum_i\varpi_i/|F_i|=0$. So $\varpi\equiv0$, and $\tau_i=0$ whenever $\ell_i\ell_{i+1}\ne0$. Let $Z$ be the set of $\Gamma$-spokes with $\ell_p=0$, i.e.\ $s_p\parallel\mu$; $|Z|\le2$, two such spokes point in opposite directions, and no face has both spokes in $Z$. For $p\in Z$, $(\nu_p)$ defines $X_p:=\tau_{p-1}\ell_{p-1}\mathbf 1_{p-1}=\tau_p\ell_{p+1}\mathbf 1_p$, which vanishes unless both adjacent faces are $\Gamma$-faces, and the only possibly nonzero $\tau$'s are $\tau_{p-1}=X_p/\ell_{p-1}$ and $\tau_p=X_p/\ell_{p+1}$. With $k_p:=|F_{p-1}|^{-1}+|F_p|^{-1}$, $(\sigma)$ with $L=0$ and with $L=\mu\cdot$ gives $\sum_{p\in Z}k_pX_p=0$ and $\sum_{p\in Z}(\mu\cdot s_p)k_pX_p=0$. The numbers $\mu\cdot s_p$ are nonzero and of opposite signs if there are two, so $X_p=0$.
\emph{Step 2: $w_i=0$ for all $i\in I$.} Now $\omega_i=\omega_{i+1}=:\Omega_i$, and $(D)$ reads $(\Omega_i-w_i)(\ell_i-\ell_{i+1})=0$. Suppose $w_i\ne0$. If $w_i=\Omega_i$, then $\hh$ takes the same nonzero value $c$ on $s_i$, $s_{i+1}$ and $s_{i+1}-s_i$, so $\hh=c\begin{psmallmatrix}1&1/2\\1/2&1\end{psmallmatrix}$ in that basis is definite; then all $w_{i'}$ have the sign of $c$, contradicting $(\delta_T)$. Otherwise $\ell_i=\ell_{i+1}$, i.e.\ $s_{i+1}-s_i\parallel\mu$, and $w_i=|s_{i+1}-s_i|^2\hh(\hat\mu,\hat\mu)$. So all nonzero $w_i$ have the sign of $\hh(\hat\mu,\hat\mu)$, and $(\delta_T)$ forces them to vanish. Hence every $\Gamma$-face is degenerate at $z$, a contradiction.
\end{proof}
An exact search over all real $\hh\ne0$ (eliminating $\hh$ through $3\times3$ minors) on $7$ structured and $3034$ random $\Gamma$-flat stars found degeneracy only when every $\Gamma$-face is degenerate, and full finite element ranks confirm the predictions on seven stars, including a saddle with $\sum_iw_{ab}=0$ (\texttt{notes/v6/td\_k3b/degen*.log, fe\_cross.log}). \status{numerical, exact arithmetic}

For $F\in\FG$ let $z_F$ be the vertex of $F$ opposite an edge $e_F$ that maximises $|\hh_{y_F}(P_{y_F}e,P_{y_F}e)|$ over the edges of $F$, and $\omega_F:=R(z_F,q_F)$.

\begin{theorem}[Uniform witness, $k=3$]\label{k3:witness}
On Alfeld refinements there are $\kappa_0=\kappa_0(\theta,\theta_0)>0$, $C_1$ and $h_\ast$ such that for $\hG\le h_\ast$, (M3$^3$) holds with $\omega_F=R(z_F,q_F)$ and $K_{\rm ov},C_s\le C(\theta)$.
\end{theorem}

\begin{proof}
Let $\ell:=\diam F$ and $\hh_0:=\hh_{y_F}\circ P_{y_F}$.
\emph{(1) Weights.} Every $\Gamma$-face $F'$ of $\omega_F$ lies within $C\ell$ of $y_F$, so $w^{F'}_{ij}=\hh_0(e_{ij},e_{ij})+O(\ell^3)$ (Lemma~\ref{geom:face}(a), $\Gamma\in C^3$). Moments are linear in the weights with $|\hat\M_{F'}(v)|\le C|w|\ell^{1/2}\norm{\nabla v}$, so replacing the weights costs $C\ell^{7/2}\norm{\nabla v}$, the $C_1$ term. If $\hh_0=0$ there is nothing to prove; otherwise write $\hh_0=|\hh_0|\hat\hh$ with $|\hat\hh|=1$.
\emph{(2) The parameter map.} For admissible data $x=(\delta,p_q,(p_j),(s_i))$, Lemmas~\ref{k3:suff} and~\ref{k3:single} give $v(x)\in Y^1(R)$ with $\norm{\nabla v(x)}\le C_E(\theta)|x|\ell^{1/2}$, where $C_E$ is uniform for configurations in a compact set of normalised shapes. Let $\Pi(\mathcal C;\hat\hh)$ be the linear map from admissible data to $\ell^{-3}P_t\sum\hat\M_{F'}(v(x))\in\R^2$, for the normalised configuration $\mathcal C$ and weights $\hat\hh(e,e)$. By Lemma~\ref{k3:face}, which holds for every geometry, $\Pi$ is an explicit rational function of $\mathcal C$, linear in $\hat\hh$, on an admissible subspace of constant dimension (each flux equation contains its own $p_j$ with $A_j\ne0$). So $s(\mathcal C,\hat\hh):=\sigma_2(\Pi)$, the second singular value of a map into $\R^2$, is continuous.
\emph{(3) Compactness.} Let $\mathcal K$ be the set of normalised $\Gamma$-flat configurations with $m\le m_{\max}(\theta)$, containing the face $F_{i_0}$ of diameter $1$, macro-tetrahedra of minimal angle at least $\theta/2$, $\Gamma$-spokes separated by angles at least $\theta_0/2$ and consistently oriented, $|\hat\hh|=1$ and $|\hat\hh(\hat e_F,\hat e_F)|\ge c_1(\theta_0)$. The last condition can be imposed because $\hh\mapsto(\hh(e,e))_e$ is invertible with inverse bounded by $C(\theta_0)\ell^{-2}$ (Lemma~\ref{geom:face}(d)), so the maximising edge has $|\hat\hh(\hat e_F,\hat e_F)|\ge c_1$. The spoke separation holds because the $\Gamma$-spokes at $z$ are distinct edges of the embedded surface $\Gh$ around $z$ (Theorem~\ref{geom:orient}) and (G1) bounds their angles below. $\mathcal K$ is compact. On $\mathcal K$, $F_{i_0}$ is nondegenerate at $z$ (Lemma~\ref{k3:degface}), so $s>0$ by Theorem~\ref{k3:nondeg}, and $s\ge s_0(\theta,\theta_0)>0$.
\emph{(4) Actual meshes.} Projecting the $\Gamma$-vertices of $\omega_F$ onto $T_{y_F}\Gamma$ moves them by $O(\ell^2)$, i.e.\ $O(\ell)$ after normalisation, and $\hh_0(Pe,Pe)=\hh_0(e,e)$; the distinction between $t\in T_{y_F}\Gamma$ and the plane of the projected faces is also $O(\ell)$. By uniform continuity of $\Pi$ on a compact neighbourhood of $\mathcal K$, $\sigma_2(\Pi_{\rm actual})\ge s_0/2$ for $\ell\le h_\ast$. For a unit $t$ choose $x\propto\Pi^{\mathsf T}t$ normalised; then $t\cdot\sum\hat\M\ge\frac{s_0}{2C_E}|\hh_{y_F}|\ell^{5/2}-C\ell^{7/2}$.
\end{proof}
The constant $\kappa_0$ comes from compactness and is not explicit; the explicit criterion of Theorem~\ref{k3:sphere} is available on the sphere.

\section{Explicit constants}\label{sec:constants}

The constants in Theorems~\ref{sh:alfeld}, \ref{sh:lb} and~\ref{lk:low} and the single-cell penalty constant can be made explicit on Alfeld refinements. The (IS3) constant cannot, at present; we say why. The values below are certified lower bounds for the stated shapes, not estimates of the true constants: the pre-asymptotic sphere data of Section~\ref{sec:numerics} are orders of magnitude larger.

\subsection{The witness constants}
\begin{proposition}[Explicit $\kappa_0$ and $\kappa_L$]\label{ct:kappa}
Let $k\ge4$, $T=T_F$ a macro-tetrahedron of an Alfeld refinement, $A$ its reference map (any labelling with $\det A>0$), $\ell=\diam F$, $h_T$ the height of $T$ over $F$, $D=\diam T$, $r$ the inradius, $c_s:=r/D$, and $\theta$ the smaller of the minimal angles of $F$ and of its projection onto $T_{y_F}\Gamma$ ($\theta\ge\theta_0/2$ for $\hG\le h_\ast$). Then (M3$^3$)(iii) holds with
\[
\kappa_0(T)=\frac{c_w(F)\,(\det A)^{1/2}}{2h_T^2\norm A\norm{A^{-1}}^2\norm{\hat R}\,\ell^{1/2}},\qquad c_w(F):=\frac{\sigma_{\min}\bigl(H\mapsto(H(u_i,u_i))_{i=1}^3\bigr)}{\ell^2}\ge\frac{\sin^4\theta}{\sqrt{10}},
\]
where $\sigma_{\min}$ is taken with the Frobenius norm on symmetric $H$ and $u_1,u_2,u_3$ are the projected edge vectors, and
\[
\norm{\hat R}=91.678155469\quad\text{(certified; the interval has width below the last digit)}.
\]
In closed form,
\[
\kappa_0(T)\ \ge\ \frac{\sin^5\theta\;c_s^{11/4}\,(16\pi/\sqrt3)^{1/4}}{\sqrt{30}\,\norm{\hat R}},\qquad
\kappa_L(T)\ \ge\ \frac{\sin^5\theta\;c_s^{13/4}\,(16\pi/\sqrt3)^{3/4}}{\sqrt6\,\norm{\hat R_1}} .
\]
Since $\hat{\mathcal Y}_4\subset\hat{\mathcal Y}_k$, both hold for every $k\ge4$. \status{computer-assisted, exact} for $\norm{\hat R}$ and $\norm{\hat R_1}$; the rest is proved.
\end{proposition}
\begin{proof}
The first formula is Proposition~\ref{sh:joint} with $|P_Ft|=1-O(\hG^2)$ (the defect goes into $C_1$) and, by the definition of $c_w$, $|w^F|\ge c_w\ell^2|\hh_{y_F}|_{\rm F}\ge c_w\ell^2|\hh_{y_F}|$, up to the $O(\ell^3)$ of Lemma~\ref{geom:face}(d). For the bound on $c_w$ let $E=[e_1,e_2]$ be two edges and $w_3=(e_1+e_2)^{\mathsf T}H(e_1+e_2)$. Then $E^{\mathsf T}HE=\bigl(\begin{smallmatrix}w_1&(w_3-w_1-w_2)/2\\ (w_3-w_1-w_2)/2&w_2\end{smallmatrix}\bigr)$ has Frobenius norm at most $\sqrt{5/2}\,|w|$, and $|E^{\mathsf T}HE|_{\rm F}\ge\sigma_{\min}(E)^2|H|_{\rm F}$, so $|H|_{\rm F}\le|E^{\mathsf T}HE|_{\rm F}/\sigma_{\min}(E)^2$; the bound is in the Frobenius norm, as the definition of $c_w$ requires. Finally $\sigma_{\min}(E)^2\ge(2|F|)^2/\tr(E^{\mathsf T}E)\ge\frac12\ell^2\sin^4\theta$, using $|F|\ge\frac12\ell^2\sin^2\theta$ (sine rule). For the closed forms: $\det A=2|F|h_T$, $(2|F|)^{1/2}\ge\ell\sin\theta$, $\norm A\le\sqrt3D$, $\norm{A^{-1}}\le\sqrt2/(2r)$ \cite[Theorem~3.1.3]{Ciarlet1978book} (the reference diameter is $\sqrt2$), $h_T\le D$, and $\ell/D\ge(16\pi/\sqrt3)^{1/2}c_s^{3/2}$ from $\frac43\pi r^3\le|T|=\frac13|F|h_T\le\frac13\frac{\sqrt3}4\ell^2D$. For $\kappa_L$, $\sigma_{\min}(A_F)^2\ge\frac12\ell^2\sin^4\theta$ in the same way.
\end{proof}
\cert{notes/v7/td/kappa\_cert.py 4; kappa\_explicit.py}{kappa\_k4.log; kappa\_explicit.log}
The value $\norm{\hat R}=91.678155469$ is the certified form of the floating-point value $91.7$ quoted after Theorem~\ref{sh:alfeld}; it is obtained, like $\norm{\hat R_1}$, by bracketing the smallest eigenvalue of $\hat E\hat G^{-1}\hat E^{\mathsf T}$ with exact $LDL^{\mathsf T}$ tests.

\begin{center}\small
\textbf{Table C1.} Explicit constants (floating-point evaluation of the proved formulas, best labelling).\\[2pt]
\begin{tabular}{lcccccc}
\toprule
macro shape (apex over base) & $c_w(F)$ & $\kappa_0(T)$ & $\kappa_L(T)$ & $c_s$ & $\theta$ & closed-form $\kappa_0$\\
\midrule
regular & 0.866 & $2.11\cdot10^{-3}$ & $8.65\cdot10^{-2}$ & 0.204 & $60^\circ$ & $2.85\cdot10^{-5}$\\
near-regular $(\frac12,\frac7{24},\frac{13}{16})$ & 0.861 & $2.09\cdot10^{-3}$ & $8.74\cdot10^{-2}$ & 0.203 & $59.5^\circ$ & $2.73\cdot10^{-5}$\\
right corner $(0,0,1)$ & 0.331 & $1.52\cdot10^{-3}$ & $1.63\cdot10^{-1}$ & 0.149 & $45^\circ$ & $4.39\cdot10^{-6}$\\
flat, height $\frac12$ & 0.861 & $1.79\cdot10^{-3}$ & $1.18\cdot10^{-1}$ & 0.166 & $59.5^\circ$ & $1.57\cdot10^{-5}$\\
skew $(\frac32,\frac13,\frac34)$ & 0.861 & $1.33\cdot10^{-3}$ & $7.73\cdot10^{-2}$ & 0.095 & $59.5^\circ$ & $3.40\cdot10^{-6}$\\
\bottomrule
\end{tabular}
\end{center}
The bases are those of Theorem~\ref{pen:single}.

\emph{The remaining constants of the lower bounds.} In Corollary~\ref{sh:abstract}, $\Lambda_0=2+C_{PT}C_I$, where on the Alfeld sub-tetrahedron $T_F=\conv(B,F)$, of height $h_{\rm sub}=h_T/4$ over $F$ and diameter $D_{\rm sub}$,
\[
C_I^2\le k(k+2)\,\ell/h_{\rm sub},\qquad C_{PT}^2\le\frac{D_{\rm sub}}{h_{\rm sub}}\Bigl(\frac3{\pi^2}+\frac2\pi\Bigr).
\]
The first is the trace-inverse inequality of Warburton and Hesthaven \cite{WarburtonHesthaven2003} for $\dn v\in P_{k-1}$ in three dimensions, restricted to one face. The second follows by integrating $\div((x-B)|w-\bar w_{T_F}|^2)$ over $T_F$ and applying the Payne--Weinberger inequality \cite{PayneWeinberger1960} on the convex set $T_F$. For regular macros and $k=4$: $C_I\le10.84$, $C_{PT}\le2.15$, $\Lambda_0\le25.28$. Hence
\[
c_L=5.20\cdot10^{-3}\ \text{in Theorem~\ref{lk:low}},\qquad c=\frac{\kappa_0c_0^{5/2}c_A^{1/2}}{\Lambda_0}\cdot\frac12\Bigl(\frac34\Bigr)^{1/2}=2.37\cdot10^{-5}\ \text{in Theorem~\ref{sh:lb}},
\]
with $c_0=1$ and $c_A=\sqrt3/4$. For $k=3$ on general surfaces (Theorem~\ref{k3:witness}) $\kappa_0$ comes from compactness and is not explicit.

\subsection{The single-cell penalty constant over a shape class}
\begin{proposition}[$\lambda_T\ge20$ on an explicit box of shapes]\label{ct:pen}
Normalise $F=\conv(0,e_1,(a,b,0))$ and the apex $(x,y,z)$; every tetrahedron with a marked face is similar to one of these. For every $(a,b,x,y,z)$ in the box
\[
\bigl|(a,b,x,y,z)-\bigl(\tfrac12,\tfrac{\sqrt3}2,\tfrac12,\tfrac{\sqrt3}6,\sqrt{2/3}\bigr)\bigr|_\infty\le0.1
\]
around the regular tetrahedron, $\lambda_T\ge20$ for every $k\ge3$. Hence on Alfeld refinements whose macro-tetrahedron at a smallest wall face lies in this class, coercivity of $N_h$ on $\Zh$ requires $\mu\ge20\rho$ (Theorem~\ref{pen:single}). \status{computer-assisted, interval arithmetic}
\end{proposition}
\begin{proof}
The box is bisected into a tree of $879$ boxes with $440$ leaves. For each box a rational witness $\hat v\in W_3(\hat T)$ is fixed (the floating-point maximiser at the box centre, rounded to rationals); the pulled-back forms of Lemma~\ref{pen:piola} are evaluated in ball arithmetic (Arb \cite{Johansson2017}, $80$ bits) with a mean-value enclosure (forward-mode derivatives in ball arithmetic over the whole box). Every leaf gives $\diam F\cdot(2B_T-A_T)/C_T\ge20.0004$. The enclosure was cross-checked against the floating-point pipeline at a random shape ($40.5657$ both ways) and against the near-regular value $32.50$ of Theorem~\ref{pen:single}, which lies in the box. $W_3\subset W_k$ gives $k\ge3$.
\end{proof}
\cert{notes/v7/td/pen\_cert.py 20 0.1 (and 20 0.05)}{pen\_cert\_d010.log, pen\_cert\_d005.log}
\emph{Scope of the certificate.} The centres and half-widths of the boxes, including the children $m\pm r/2$ of each bisection, were computed in binary floating point. Strictly, the certified region is therefore the union of the $440$ leaf boxes \emph{as stored in floating point}, which can differ from the exact box above by gaps and overlaps of one unit in the last place between siblings. Given the margin $20.0004-20$ this is immaterial for the conclusion, but it is not covered by the certificate; a rerun with dyadic rational centres would remove the caveat. A box of half-width $0.15$ was not certified within our time budget: the depth-first mean-value scheme stalls, which says nothing about $\lambda_T$ there. The tall-cell obstruction of Theorem~\ref{pen:single} (negative $\lambda_T$ above height about $1.69$) shows that no shape-uniform positive bound exists.

\subsection{The (IS3) constant}
Lemma~\ref{st:IS3} gives $\beta^{-1}\le C_{\rm loc}+C_\Pi C_B(\Oh)(1+\sqrt3C_{\rm loc})$. It is not explicit, for three reasons.
\begin{enumerate}
\item $C_{\rm loc}$: \cite[Theorem~3.1]{GuzmanNeilan2018} states only that it depends on $k$ and shape regularity. It can be made explicit by Piola transport: for $p\in\mathring P_{k-1}(K^r)$ put $\hat p:=(\det A)\,p\circ\Phi$; the reference right inverse gives $\hat v$ with $\norm{\hat\nabla\hat v}\le\hat C_{\rm loc}(\det A)^{1/2}\norm p_K$, and the Piola image has $\div v=p$ and $\norm{\nabla v}_K\le\norm A\norm{A^{-1}}\hat C_{\rm loc}\norm p_K\le\sqrt{3/2}\,c_s^{-1}\hat C_{\rm loc}\norm p_K$. \status{numerical} $\hat C_{\rm loc}=6.18$ ($k=3$) and $6.41$ ($k=4$), from generalised eigenvalues in floating point; these could be certified by the same $LDL^{\mathsf T}$ bracketing, which we have not done.
\item $C_\Pi$ (Scott--Zhang plus face-bubble fluxes) is local and explicit in principle, but depends on the vertex patches; we have not derived a value.
\item $C_B(\Oh)\le2C_B(\Omega)$ holds only for $\hG\le h_\ast$ with a non-explicit $h_\ast$ (Lemma~\ref{st:bog}), and $C_B(\Omega)$ for a box minus a ball is not known explicitly. Explicit Bogovski\u{\i} bounds exist for domains star-shaped with respect to a ball, and $R\setminus\overline D$ is not one; decomposition bounds are explicit but very crude.
\end{enumerate}
\begin{center}\small
\textbf{Table C2.} \status{numerical} Discrete inf-sup constant of $(\VO,\Pih\cap L^2_0)$ on icosphere--Alfeld shells: $\Oh$ lies between the level-$l$ icosphere polyhedron and its dilate by $2$, with $m$ prism layers (radii $2^{j/m}$), $3$ tetrahedra per prism, Alfeld-refined. $\beta_h^{-2}$ is the largest eigenvalue of $M^{1/2}(BA^{-1}B^{\mathsf T})^+M^{1/2}$ on $L^2_0$ (ARPACK, sparse LU of the saddle system).\\[2pt]
\begin{tabular}{cccccccc}
\toprule
$k$ & $l$ & $m$ & macro tets & $\hG$ & $\dim\VO$ & $\beta_h$ (3 smallest) & $1/\beta_h$\\
\midrule
3 & 0 & 1 & 60 & 1.052 & 3252 & .1162, .1162, .1213 & 8.61\\
3 & 1 & 1 & 240 & 0.618 & 12972 & .1161, .1205, .1206 & 8.61\\
3 & 0 & 2 & 120 & 1.052 & 6780 & .1058, .1062, .1064 & 9.45\\
3 & 0 & 3 & 180 & 1.052 & 10308 & .0833, .0834, .0836 & 12.01\\
4 & 0 & 1 & 60 & 1.052 & 7578 & .1230, .1231, .1277 & 8.13\\
4 & 0 & 2 & 120 & 1.052 & 15642 & .1152, .1155, .1156 & 8.68\\
\bottomrule
\end{tabular}
\end{center}
\cert{notes/v7/td/is3\_ico.py k l m}{is3\_k3\_l0.log, is3\_more.log}
At fixed layer structure the constant does not change from $l=0$ to $l=1$ ($8.608$ to $8.612$). More layers at fixed $l$ make the prisms flatter (smaller $c_s$) and $\beta_h$ smaller, consistent with the $c_s^{-1}$ dependence of $C_{\rm loc}$. No spurious pressure modes were found. These values, $1/\beta_h\approx8.6$ ($k=3$) and $8.1$ ($k=4$) on single-layer shells, are evidence for, not a proof of, a moderate constant.

\section{Reproducibility}\label{app:repro}

All code, logs and raw records are in the repository \url{https://github.com/jaideepsaipadhi/zero-gradient-prize}. Paths are relative to its root. %% TBD: the notes/v6/* scripts may be moved to a stable location (e.g. certificates3d/) before the Zenodo release; update the paths below accordingly.

\subsection*{Exact and computer-assisted certificates}
{\small
\begin{longtable}{>{\raggedright\arraybackslash}p{0.30\textwidth}>{\raggedright\arraybackslash}p{0.38\textwidth}>{\raggedright\arraybackslash}p{0.24\textwidth}}
\toprule
Result & Command & Output\\
\midrule
Theorem~\ref{sh:alfeld} (rank $8$, $k=4$) & \path{python3 notes/v6/td_sharp/alfeld_witness.py 4} & \texttt{alfeld\_k4.log}\\
$\norm{\hat R}=91.7$ & \path{python3 notes/v6/td_sharp/k4_constant.py} & \texttt{k4\_constant.log}\\
Proposition~\ref{sh:bubble} & \path{python3 notes/v6/td_sharp/bubble_variance.py} & \texttt{bubble\_variance.log}\\
Proposition~\ref{sh:k3single} & \path{python3 notes/v6/td_sharp/alfeld_witness.py 3} & \texttt{alfeld\_k3.log}\\
Lemmas~\ref{k3:single}, \ref{k3:K} & \path{notes/v6/td_k3/trace1.py}, \texttt{ref\_K.py}; \path{notes/v6/referee8/ref8_lemmaK.py} & logs alongside\\
Theorem~\ref{k3:trace} (reduced model) & \path{notes/v6/td_k3/validate_reduced.py} & \texttt{validate\_reduced.log}\\
Theorem~\ref{k3:vertex} & \path{notes/v6/td_k3/vertex_formula_general.py}, \texttt{fe\_vertex\_check.py}; \path{notes/v6/referee8/ref8_vertex_fe.py} & logs alongside\\
Icosphere criterion & \path{notes/v6/td_k3/icosphere_check.py} & \texttt{icosphere\_check.log}\\
Lemmas~\ref{k3:face}, \ref{k3:flat} & \path{notes/v6/td_k3b/ref_nu.py}, \texttt{face\_formula.py}, \texttt{flat\_algebra.py} & logs alongside\\
Theorem~\ref{k3:nondeg} checks & \path{notes/v6/td_k3b/degen.py}, \texttt{degen\_random.py}, \texttt{degen\_gflat.py}, \texttt{fe\_cross.py} & logs alongside\\
Theorem~\ref{pen:single} & \path{python3 notes/v6/td_pen/alfeld_tet.py 3 [shape]}, \texttt{verify\_alfeld.py}, \texttt{neg\_cert.py} & \texttt{alfeld\_k3.log}, \texttt{neg\_cert.log}\\
Half-space constant & \path{python3 notes/v6/td_pen/continuum3d.py} & \texttt{continuum3d.log}\\
\bottomrule
\end{longtable}}
\noindent Each script runs in under five minutes on two cores, except the height scan and the $k=4$ single-cell pencils (\texttt{alfeld\_tet.py 4}), estimated at 10--30 minutes per shape and not run.

\subsection*{Three-dimensional solver (\texttt{code3d/})}
\begin{itemize}
\item \texttt{svn3d.py}: the solver (Section~\ref{sec:numerics}); \texttt{exact3d.py}: the exact solutions; \texttt{run3d.py}: drivers (box, sphere, interpolation); \texttt{table3d.py}: tables; \texttt{penalty3d.py}: the threshold $\mu_Z^\ast$ (PARDISO inertia); \texttt{infsup3d.py}: the discrete inf-sup constant (dense, coarse meshes only); \texttt{patch\_test.py}.
\item \texttt{local\_*.sh}: exactly the runs of Section~\ref{sec:numerics}. Records: \path{notes/v6/num3d/results.jsonl}; tables: \path{notes/v6/num3d/tables_local.txt}.
\item Requirements: Python~3 with numpy, scipy, sympy and pypardiso (MKL; \texttt{PYPARDISO\_MKL\_RT} points to \texttt{libmkl\_rt.so.3}).
\end{itemize}

\subsection*{Finer levels}
The levels that did not fit our machine are scripted in \path{code3d/run_pod3d.sh}. From the repository root:
\begin{verbatim}
sh code3d/run_pod3d.sh            # all stages, sequentially
sh code3d/run_pod3d.sh STAGE      # one of: box4 pen k3a k3b k3c pen4 k4
python3 code3d/table3d.py         # tables from results.jsonl
\end{verbatim}
Memory was measured locally for $k=3$ (PARDISO $LDL^{\mathsf T}$ of the velocity matrix, peak resident size): sphere $N=4$ ($70$k velocity unknowns) $1.0$\,GB; box $N=6$ ($79$k) $1.15$\,GB. Extrapolated with fill $\propto n^{4/3}$: sphere $N=6$ ($233$k) $5$--$6$\,GB, $N=8$ ($550$k) $\approx17$\,GB, $N=10$ ($1.07$M) $\approx40$\,GB, $N=12$ ($1.85$M) $\approx80$\,GB; $k=4$ multiplies the unknowns by about $2.4$.
\begin{center}\small
\begin{tabular}{ll}
\toprule
stage & content\\
\midrule
\texttt{box4} & $k=4$ box verification, $N=2,\dots,5$\\
\texttt{pen} & $\mu_Z^\ast$ on the sphere meshes $N=6,8$ (run first: confirms $\mu=400>\mu_Z^\ast$)\\
\texttt{k3a} & tests P, sph, rot; $k=3$, $\mu=400$, $N=6,8$\\
\texttt{k3b} & the same at $N=10,12$\\
\texttt{k3c} & test P at $\mu=1000$, $N=2,\dots,8$ (penalty independence of the leak constant)\\
\texttt{pen4}, \texttt{k4} & $\mu_Z^\ast$ for $k=4$, then tests P, sph, rot with $k=4$, $\mu=1000$, $N=2,4,6$\\
\bottomrule
\end{tabular}
\end{center}


\section*{Use of generative-AI tools declaration}
%% TBD: if authorship becomes joint, rephrase in the plural and add a CRediT statement.
The author used Anthropic's Claude Opus~5.5 as a research-support tool. It assisted with the following:
\begin{itemize}
\item developing and checking the proofs;
\item implementing, running and validating the three-dimensional finite element code and the exact-arithmetic certificates;
\item independently reviewing the arguments as a referee;
\item preparing the manuscript and the accompanying repository.
\end{itemize}
The author directed the project and retains authorship of the work. Responsibility for all mathematical statements, computations, references and conclusions rests with the author.

\section*{Data and code availability}
All code, exact-arithmetic certificates, logs and raw records are available at \url{https://github.com/jaideepsaipadhi/zero-gradient-prize} (directories \texttt{code3d/} and \texttt{notes/v6/}); Appendix~\ref{app:repro} lists the command behind each number. The two-dimensional companion paper \cite{PadhiLong2D} and its code are in the same repository.


\bibliographystyle{amsplain}
\bibliography{refs}

\end{document}
